%~Mouliné par MaN_auto v.0.27.3 2024-02-05 10:33:00
\documentclass[AHL,Unicode,longabstracts,published]{cedram}




\usepackage{amscd}
\usepackage{dsfont}
%\DeclareMathOperator{\Im}{Im}
\DeclareMathOperator{\lcm}{lcm}
\DeclareMathOperator{\Arg}{Arg}
\DeclareMathOperator{\Hom}{Hom}
\DeclareMathOperator{\SM}{SM}
\DeclareMathOperator{\Herm}{Herm}
\DeclareMathOperator{\Tr}{Tr}
\DeclareMathOperator{\Pic}{Pic}
\DeclareMathOperator{\Proj}{Proj}
\DeclareMathOperator{\Vol}{Vol}
\DeclareMathOperator{\Span}{Span}
\DeclareMathOperator{\rank}{rank}
\DeclareMathOperator{\diff}{d}
\DeclareMathOperator{\dbar}{\bar\partial}
\DeclareMathOperator{\ddbar}{\partial\bar\partial}
\DeclarePairedDelimiter{\abs}{\lvert}{\rvert}
\DeclarePairedDelimiter{\Abs}{\lVert}{\rVert}





%\AtBeginDocument{ }



\newenvironment{equation'}{
\begin{equation}
\tag{\theequation\ensuremath{'}}}{
\end{equation}
}
\newenvironment{equation_}[1]{
\begin{equation}
\tag{\ensuremath{\theequation_{#1}}}}{
\end{equation}
}
\newenvironment{eqlist}{
\begingroup \leftskip\parindent \rightskip\parindent \noindent
\begin{minipage}{\linewidth-2\parindent}} {
\end{minipage}
\endgroup}


%\let\Im\relax







\newcommand*\rS{\mathrm{S}}
\newcommand*\Id{\mathrm{Id}}
\newcommand*\Sing{\mathrm{Sing}}
\newcommand*\orb{\mathrm{orb}}
\newcommand*\FS{\mathrm{FS}}
\newcommand*\loc{\mathrm{loc}}
\newcommand{\rma}{\mathrm{a}}
\newcommand{\rmb}{\mathrm{b}}


\newcommand*\cI{\mathcal{I}}
\newcommand*\cJ{\mathcal{J}}
\newcommand*\cO{\mathcal{O}}

\newcommand*\bC{\mathbb{C}}
\newcommand*\bD{\mathbb{D}}
\newcommand*\bE{\mathbb{E}}
\newcommand*\bG{\mathbb{G}}
\newcommand*\bN{\mathbb{N}}
\newcommand*\bP{\mathbb{P}}
\newcommand*\bQ{\mathbb{Q}}
\newcommand*\bR{\mathbb{R}}
\newcommand*\bS{\mathbb{S}}
\newcommand*\bZ{\mathbb{Z}}
\newcommand*\bB{\mathbb{B}}


\newcommand*\bOne{{\mathchoice {\rm 1\mskip-4mu l} {\rm 1\mskip-4mu l} {\rm 1\mskip-4.5mu l}{\rm 1\mskip-5mu l}}}

\newcommand\bDel{\mbox{$\Delta$\newline \scalebox{.6}{\raisebox{3pt}{$\mkern-22mu\not{}\mkern12mu$}}}}
%\newcommand\bDel{\ooalign{$\Delta$\\ \scalebox{.6}{\raisebox{3pt}{\,$\not{}$}}}}
\newcommand\bDelta{\mathchoice{\bDel}{\bDel}{\scalebox{.7}{\bDel}}{\scalebox{.5}{\bDel}}}
%\newcommand\bDelta{\mbox{$\Delta\mkern-2.5mu\not{}$}}


\let\leq\leqslant
\let\geq\geqslant
\newcommand*\bydef{\coloneqq}



%%%%%%%%%%%%%%%%%%%%%%%%%%%%%%%%%%%%%%%%%%%%%%%%%%%%%%%%%%%%%%%%%%%%%%%%%%%%%%%%%%%%%%%%%%%%%%%%%%%%%%%%%%%%%%%%%%%%%%%%%%%%%%%%%%%%%%%%%%%%%%
\graphicspath{{./figures/}}

\newcommand*{\mk}{\mkern -1mu}
\newcommand*{\Mk}{\mkern -2mu}
\newcommand*{\mK}{\mkern 1mu}
\newcommand*{\MK}{\mkern 2mu}

\hypersetup{urlcolor=purple, linkcolor=blue, citecolor=red}

\newcommand*{\romanenumi}{\renewcommand*{\theenumi}{\roman{enumi}}}
\newcommand*{\Romanenumi}{\renewcommand*{\theenumi}{\Roman{enumi}}}
\newcommand*{\alphenumi}{\renewcommand*{\theenumi}{\alph{enumi}}}
\newcommand*{\Alphenumi}{\renewcommand*{\theenumi}{\Alph{enumi}}}
\let\oldtilde\tilde
\renewcommand*{\tilde}[1]{\mathchoice{\widetilde{#1}}{\widetilde{#1}}{\oldtilde{#1}}{\oldtilde{#1}}}
%\let\oldhat\hat
%\renewcommand*{\hat}[1]{\mathchoice{\widehat{#1}}{\widehat{#1}}{\oldhat{#1}}{\oldhat{#1}}}
\let\oldexists\exists
\renewcommand*{\exists}{\mathrel{\oldexists}}
\let\oldforall\forall
\renewcommand*{\forall}{\mathrel{\oldforall}}
%%%%%%%%%%%%%%%%%%%%%%%%%%%%%%%%%%%%%%%%%%%%%%%%%%%%%%%%%%%%%%%%%%%%%%%%%%%%%%%%%%%%%%%%%%%%%%%%%

\title[On the existence of logarithmic and orbifold jet differentials] {On the existence of logarithmic and orbifold jet differentials}
\alttitle{Sur l'existence des différentielles de jets logarithmiques et orbifoldes}


\subjclass{32Q45, 32H30, 14F06}
\keywords{Projective variety, directed variety, orbifold, ramification divisor, entire curve, jet differential, Green--Griffiths conjecture, algebraic differential operator, holomorphic Morse inequalities, Chern curvature, Chern form}


\author[\initial{F.} \lastname{Campana}]{\firstname{Frédéric} \lastname{Campana}}
\address{Institut de Mathématiques \'Elie Cartan,\\
Université de Lorraine,\\
B.P.\ 70239,\\
54506 Vand{\oe}uvre-lès-Nancy (France)}
\email{frederic.campana@univ-lorraine.fr}

\thanks{The third author is supported by the Advanced ERC grant ALKAGE, no 670846 from September 2015, attributed by the European Research Council.}


\author[\initial{L.} \lastname{Darondeau}]{\firstname{Lionel} \lastname{Darondeau}}
\address{Sorbonne Université \\
and Université Paris Cité,\\
CNRS, IMJ-PRG,\\
F-75005 Paris (France)}
\email{darondeau@imj-prg.fr}

\author[\initial{J.-P.} \lastname{Demailly}]{\firstname{Jean-Pierre} \lastname{Demailly}}
\address{Université Grenoble Alpes,\\
Institut Fourier,\\
100 rue des Maths,\\
38610 Gières (France)}
\email{jean-pierre.demailly@univ-grenoble-alpes.fr}

\author[\initial{E.} \lastname{Rousseau}]{\firstname{Erwan} \lastname{Rousseau}}
\address{Univ Brest,\\ 
CNRS, UMR 6205,\\
Laboratoire de Mathématiques\\
de Bretagne Atlantique,\\
F-29200 Brest (France)}
\email{erwan.rousseau@univ-brest.fr}



\begin{abstract}
We introduce the concept of directed orbifold, namely triples 
$(X,V,D)$ formed by a directed algebraic or analytic variety $(X,V)$, and a ramification divisor $D$, where $V$ is a coherent subsheaf of the tangent bundle $T_X$. In this context, we introduce an algebra of orbifold jet differentials and their sections. These jet sections can be seen as algebraic differential operators acting on germs of curves, with meromorphic coefficients, whose poles are supported by $D$ and multiplicities are bounded by the ramification indices of the components of $D$. We estimate precisely the curvature tensor of the corresponding directed structure $V\langle D\rangle$ in the general orbifold case -- with a special attention to the compact case $D=0$ and to the logarithmic situation where the ramification indices are infinite. Using holomorphic Morse inequalities on the tautological line bundle of the projectivized orbifold Green--Griffiths bundle, we finally obtain effective sufficient conditions for the existence of global orbifold jet differentials.
\end{abstract}
\begin{altabstract}
Nous introduisons le concept d'orbifoldes dirigées, à savoir les triplets $(X,V,D)$ formés par une variété dirigée algébrique ou analytique $(X,V)$, et un diviseur de ramification $D$, où $V$ est un sous-faisceau cohérent du fibré tangent $T_X$. Dans ce contexte, nous introduisons une algèbre de différentielles de jets orbifoldes et leurs sections. Ces sections peuvent être vues comme des opérateurs différentiels algébriques agissant sur les germes de courbes, à coefficients méromorphes, dont les pôles sont supportés par $D$ et les multiplicités sont bornées par les indices de ramification des composantes de $D$. Nous estimons avec précision le tenseur de courbure de la structure dirigée correspondante $V\langle D\rangle$ dans le cas orbifolde général -- avec une attention particulière pour le cas compact $D=0$ et le cas logarithmique où les indices de ramifications sont infinis. En utilisant les inégalités de Morse holomorphes sur le fibré en droites tautologique du fibré projectivisé orbifolde de Green--Griffiths, nous obtenons finalement des conditions suffisantes pour l'existence de différentielles de jets orbifoldes globales.
\end{altabstract}

\datereceived{2021-11-19}
\daterevised{2023-09-22}
\dateaccepted{2023-10-23}
\editors{X. Caruso and C. Xu}
%%%%%%%%%%%%%%%%%%%%%%%%%%%%%%%%%%%%%%%%%%%%%%%%%%%%%%%%%%%%%%%%%%%%%%%%%%%%%%%%%%%%%%%%%%%%%%%%%%%%%%%%%%%%%%%%%%%%%%%%%%%%%%%%%%%%%%%%%%%%%%%
\dateposted{2024-06-27}
\begin{document}
\maketitle

\section{Introduction and main definitions}
The present work is concerned primarily with the existence of logarithmic and orbifold jet differentials on projective varieties. For the sake of generality, and in view of potential applications to the case of foliations, we work throughout this paper in the category of directed varieties, and generalize them by introducing the concept of directed orbifold.

\begin{defi}
Let $X$ be a complex manifold or variety. A \emph{directed structure} $(X,V)$ on $X$ is defined to be a subsheaf $V\subset\cO(T_X)$ such that $\cO(T_X)/V$ is torsion free. A morphism of directed varieties $\Psi\colon(X,V)\to (Y,W)$ is a holomorphic map $\Psi\colon X\to Y$ such that $\diff\Psi(V)\subseteq\Psi^*W$. We say that $(X,V)$ is non-singular if $X$ is non-singular and $V$ is locally free, i.e., is a holomorphic subbundle of $T_X$.
\end{defi}

We refer to the \emph{absolute case} as being the situation when $V=T_X$, the \emph{relative case} when $V=T_{X/S}$ for some fibration $X\to S$, and the \emph{foliated case} when $V$ is integrable, i.e.\ $[V,V]\subset V$, that is, $V$ is the tangent sheaf to a holomorphic foliation.

We now combine these concepts with orbifold structures in the sense of Campana~\cite{Cam04}.

\begin{defi}
A \emph{directed orbifold} is a triple $(X,V,D)$ where $(X,V)$ is a directed variety and where $D=\sum(1-\tfrac{1}{\rho_j})\Delta_j$ is an effective real divisor, for some irreducible hypersurfaces $\Delta_j$ with associated ``ramification numbers'' $\rho_j\in{}]1,\infty]$. We denote by $\lceil D\rceil=\sum\Delta_j$ the corresponding reduced divisor, and by $\abs{D}=\bigcup \Delta_j$ its support.
\begin{enumerate}
\item We will say that $(X,V,D)$ is non-singular if $(X,V)$ is non-singular and $D$ is a simple normal crossing divisor such that $D$ is transverse to $V$. If $r\bydef\rank(V)$, we mean by this that there are at most $r$ components $\Delta_j$ meeting at any point $x\in X$, and that for any $p$-tuple $(j_1,\,\dotsc,\,j_p)$ of indices, $1\leq p\leq r$, we have $\dim V_x\cap \bigcap_{j=1}^pT_{\Delta_{j_\ell},x}=r-p$ at any point $x\in\bigcap_{j=1}^p\Delta_{j_\ell}$.
\item If $(X,V,D)$ is non-singular, the \emph{canonical divisor} of $(X,V,D)$ is defined to be
\[
K_{V,D}=K_V+ D
\]
(in additive notation), where $K_V=\det V^*$.
\item The so-called \emph{logarithmic case} corresponds to all multiplicities $\rho_j=\infty$ being taken infinite, so that $D=\sum\Delta_j=\lceil D\rceil$.
\end{enumerate}
\end{defi}

In case $V=T_X$, we recover the concept of orbifold introduced in~\cite{Cam04}, except possibly for the fact that we allow here $\rho_j>1$ to be real or $\infty$, (even though the case where $\rho_j$ is in $\bN\cup\{\infty\}$ is of greater interest). In the sequel, we will often denote the pair $(X,D)$ by $X\langle D\rangle$ and the logarithmic cotangent sheaves by $V^*\langle D\rangle$. It would certainly be interesting to investigate the case when $(X,V,D)$ is singular, by allowing singularities in $V$ and tangencies between $V$ and $D$, and to study whether the results discussed in this paper can be extended in some way, e.g.\ by introducing suitable multiplier ideal sheaves taking care of singularities, as was done in~\cite{Dem15} for the study of directed varieties $(X,V)$. For the sake of technical simplicity, we will refrain to do so here, and will therefore leave for future work the study of singular directed orbifolds.

\begin{defi}
Let $(X,V,D)$ be a non-singular directed orbifold. We say that $f\colon\bC\to X$ is an \emph{orbifold entire curve} if $f$ is a non-constant holomorphic map such that:
\begin{enumerate}
\item
$f$ is tangent to $V$ (i.e.\ $f'(t)\in V_{f(t)}$ at every point, or equivalently $f\colon(\bC,T_\bC)\to (X,V)$ is a morphism of directed varieties;
\item
$f(\bC)$ is not identically contained in $\abs{D}$;
\item\label{cond}
at every point $t_0\in\bC$ such that $f(t_0)\in\Delta_j$, $f$ meets $\Delta_j$ with ramification number${}\geq\rho_j$, i.e., if $\Delta_j=\{z_j=0\}$ near $f(t_0)$, then $z_j\circ f(t)$ vanishes with multiplicity${}\geq\rho_j$ at $t_0$.
\item[(\arabic{enumi}\ensuremath{'})] In the case of a logarithmic component $\Delta_j$ $(\rho_j=\infty)$, condition~\eqref{cond} is to be replaced by the assumption: $f(\bC)$ does not meet $\Delta_j$.
\end{enumerate}
\end{defi}

One can now consider a category of directed orbifolds as follows.
\begin{defi}
Consider directed non-singular orbifolds $(X,V,D)$, $(Y,W,D')$ with
\[
D=\sum\Big(1-\tfrac{1}{\rho_i}\Big)\Delta_i,\qquad D'=\sum\Big(1-\tfrac{1}{\rho'_j}\Big)\Delta'_j.
\]
A morphism $\Psi\colon(X,V,D)\to(Y,W,D')$ is a morphism $\Psi\colon(X,V)\to(Y,W)$ of directed varieties satisfying the additional following properties {\rm(a,b,c)}.
\begin{enumerate}
\item for every component $\Delta'_j$, $\Psi^{-1}(\Delta'_j)$ consists of a union of components $\Delta_i$, $i\in I(j)$, eventually after adding a number of extra components $\Delta_i$ with $\rho_i=1$;
\item in case $\rho'_j<\infty$, for every $i\in I(j)$ and $z\in\Delta_i$, if $\Delta'_j=\{y_j=0\}$ near $\Psi(z)$ and $\Delta_i=\{z_i=0\}$ near $z$, then the function $z_i \to \Psi_j(z)$ vanishes with multiplicity ${}\geq \rho'_j/\rho_i $ at $0$ where $\Psi_j:=y_j \circ \Psi $;
\item if $\Delta'_j$ is a logarithmic component ($\rho'_j=\infty$), then $\Phi^{-1}(\Delta'_j)= \bigcup_{i\,\in\,I(j)}\Delta_i$ where the $(\Delta_i)_{i\,\in\,I(j)}$ consist of logarithmic components $(\rho_i=\infty)$.
\end{enumerate}
\end{defi}

It is easy to check that, if the image of the composed morphism is not contained in the support of the divisor on the target space, the composite of directed orbifold morphisms is actually a directed orbifold morphism, and that the composition of an orbifold entire curve $f\colon\bC\to(X,V,D)$ with a directed orbifold morphism $\Psi\colon(X,V,D)\to(Y,W,D')$ produces an orbifold entire curve $\Psi\circ f\colon\bC\to(Y,W,D')$ (provided that $\Psi\circ f(\bC) \not \subset \abs{D'}$). One of our main goals is to investigate the following orbifold generalization of the Green--Griffiths conjecture.

\begin{conj}
Let $(X,V,D)$ be a non-singular directed orbifold of general type, in the sense that the canonical divisor $K_V+D$ is big. Then then exists an algebraic subvariety $Y\subsetneq X$ containing all orbifold entire curves
$f\colon\bC\to(X,V,D)$.
\end{conj}

As in the absolute case ($V=T_X$, $D=0$), the idea is to show, at least as a first step towards the conjecture, that orbifold entire curves must satisfy suitable algebraic differential equations. In Section~\ref{se:jet_differentials}, we introduce graded algebras
\[
\bigoplus_{m\,\in\,\bN}E_{k,m}V^*\langle D\rangle
\]
of sheaves of ``orbifold jet differentials''. These sheaves correspond to algebraic differential operators $P(f;f',f'',\dotsc,f^{(k)})$ acting on germs of $k$-jets of curves that are tangent to $V$ and satisfy the ramification conditions prescribed by $D$. The strategy relies on the following orbifold version of the vanishing theorem, whose proof is sketched in the appendix.

\begin{prop}[Orbifold vanishing theorem]\label{prop:vanishing}
Let $(X,V,D)$ be a projective non-singular directed orbifold, and let $A$ be an ample divisor on $X$. Then, for every orbifold entire curve $f\colon\bC\to(X,V,D)$ and every global jet differential operator $P\in H^0(X,E_{k,m}V^*\langle D\rangle \otimes\cO_X(-A))$, we have $P(f;f',f'',\dotsc,f^{(k)})=0$.
\end{prop}

The next step consists precisely of finding sufficient conditions that ensure the existence of global sections
$P\in H^0(X,E_{k,m}V^*\langle D\rangle\otimes\cO_X(-A))$. Recall that it has been shown in~\cite[Prop.~5.1]{CDR20} that the general type assumption is not a sufficient condition for the existence of global jet differentials. This contrasts with the reduced case, in which we obtain (cf.~\cite{Dem11} for the compact case; the logarithmic case is proven mutatis mutandis):
\begin{theo}[Reduced case]
When the boundary divisor $D$ is reduced, the (non-singular) directed orbifold $(X,V,D)$ admits non-zero global jet differentials vanishing on an ample divisor if and only if it is of general type.
\end{theo}

Towards a condition for the existence of global jet differentials on orbifolds, ``higher order'' orbifold structures have been introduced in~\cite{CDR20}:
\[
D^{(s)}=\sum_j\bigg(1-\tfrac{s}{\rho_j}\bigg)_{\kern-3pt+}\Delta_j.
\]
where $x_+:=\max \{x, 0\}.$

The following conjecture is proposed~\cite{CDR20}.
\begin{conj}\label{conj:orbiDem}
A smooth orbifold $(X, D)$ of dimension $n\geq2$ with smooth boundary divisor admits nonzero global jet differentials vanishing on an ample divisor if and only if
$(X, D^{(n)})$ is of general type.
\end{conj}

The results in this paper can be seen as a first step in the direction of Conjecture~\ref{conj:orbiDem}. Among more general results, we obtain
\begin{theo}\label{theo:existence}
Let $D=\sum_{j=1}^N(1-\tfrac{1}{\rho_j})\Delta_j$ a simple normal crossing orbifold divisor on $\bP^n$ with
$\deg\Delta_j=d_j$. Then there exist non-zero jet differentials of order $k$ and large degree $m$ on $\bP^n\langle D\rangle$, with a small negative twist $\cO_{\bP^n}(-m\tau)$, $\tau>0$, under any of the following two sufficient conditions:
\begin{enumerate}
\item
$\,k\geq n,\quad N\geq 1,\quad \rho_j\geq \rho>n$ and\vskip-8pt
\[
\sum_j d_j\cdot
\min\left(\!\min_j\left(\tfrac{\rho_j}{ d_j}\right),\tfrac{1}{ 2}\right)
\,\prod_{s=1}^n\left(1-\tfrac{s}{\rho}\right)>c_n
\]
where
\[
c_n
\bydef n(n^2+n-1)\,n!\,\left(\sum_{s=1}^n\tfrac{1}{ s}+\tfrac{1}{ n^3}\right)^{n-1}
\sim(2\pi)^{1/2}\, n^{n+7/2}\,e^{-n}(\gamma+\log n)^{n-1}.
\]
\item
$\,k\geq 1,\quad N\geq n,\quad \rho_j\geq \rho>1$ and for $t=\max(\max(d_j/\rho_j),2),$\vskip-10pt
\[
\sum_{J\subset\{1,\,\dotsc,\,N\},\,\abs{J}=n}
\prod_{j\,\in\,J}d_j\left(1-\tfrac{1}{ \rho_j}\right)> (2n-1)\,t\,
\left(n\,t-n-1+\sum_jd_j(1-1/\rho_j)\right)^{n-1}.
\]
\end{enumerate}

When all components $(\Delta_j)_{1\,\leq\,j\,\leq\,N}$ possess the same degrees $d_j=d\geq 1$ and ramification numbers $\rho_j\geq\rho$, we get the following simpler sufficient conditions:
\begin{enumerate}\alphenumi
%%Créer une liste de type a', b'
\item
$ k\geq n, N\geq 1, \rho>n,\quad N\min(\rho,d)\,\prod_{s=1}^n(1-\tfrac{s}{\rho})>2c_n,$
\item
$k\geq 1, N\geq n, \rho>1,\quad N\min(\rho,d)\,(1-\tfrac{1}{\rho})^n>2^n\,(2n-1)\,n^n.$
\end{enumerate}
\end{theo}

Let us recall some related results previously obtained in this orbifold setting. In the case of orbifold surfaces $(\bP^2,(1-\tfrac{1}{\rho})C)$ where $C$ is a smooth curve of degree $d$, such existence results have been obtained in~\cite{CDR20} for $k=2$, $d\geq 12$ and $\rho \geq 5$ depending on $d$. In~\cite{DR20}, the existence of jet differentials is obtained for orbifolds $(\bP^n,\sum_{i=1}^d(1-\tfrac{1}{\rho})H_i)$ in any dimension for $k=1$, $\rho \geq 3$ along an arrangement of hyperplanes of degree $d\geq 2n(\tfrac{2n }{ \rho-2} + 1)$. In~\cite{BD18}, it is established that the orbifold $(\bP^n,(1-\tfrac{1}{ d})D)$, where $D$ is a general smooth hypersurface of degree $d$, is hyperbolic i.e.\ there is no non-constant orbifold entire curve $f\colon\bC\to(\bP^n, (1-\tfrac{1}{d})D)$, if $d\geq (n+2)^{n+3}(n+1)^{n+3}$ and $\rho\geq d$.

The proof of Theorem~\ref{theo:existence} depends on a number of ingredients and on rather extensive curvature calculations. The first point is that the curvature tensor of the orbifold directed structure
$V\langle D\rangle$ can be controlled in a precise manner. This is detailed in \S~\ref{sse:orbi_metric}.

\begin{theo}\label{theo:orbi_curv_estimate}
Assume that $X$ is projective and $(X,V,D)$ is non-singular. Given an ample line bundle $A$ on $X$, let $\gamma_V$ be the infimum of real numbers $\gamma\geq 0$ such that $\gamma\,\Theta_A\otimes\Id_V -\Theta_V$ is positive in the sense of Griffiths, for suitable $C^\infty$ smooth hermitian metrics on $V$. Let $D=\sum_j(1-1/\rho_j)\Delta_j$ and select $d_j\geq 0$ such that $d_jA-\Delta_j$ is nef. Then for $\gamma>\gamma_{V,D}\bydef\max(\max(d_j/\rho_j),\gamma_V)\geq 0$ and for suitable hermitian metrics on $A$, $V$, $\cO_X(\Delta_j)$, the ``orbifold metric''

\begin{multline}\label{theo1.10.a}
\tag{a}
\abs{u}_{h_{V\langle D\rangle,\varepsilon}}^2\\
\bydef
\abs{u}^2_{h_V}+\sum_{1\,\leq\,j\,\leq\,N}\varepsilon_j\,\abs{\sigma_j}^{-2+2/\rho_j}\,
\abs{\nabla_j\sigma_j(u)}^2_{h_j},~~u\in V,~~\sigma_j\in H^0\left(X,\cO_X(\Delta_j)\right)
\end{multline}
yields a curvature tensor $\gamma\,\Theta_A\otimes\Id-\Theta_{V\langle D\rangle}$ such that the associated quadratic form $Q_{V\langle D\rangle,\gamma,\varepsilon}$ on
$T_X\otimes V$ satisfies for
$\varepsilon_N\ll\varepsilon_{N-1}\ll\cdots\ll\varepsilon_1\ll 1$ the curvature estimate
\begin{multline}\label{theo1.10.b}
\tag{b} Q_{V\langle D\rangle,\gamma,\varepsilon} (z)(\xi\otimes u)\simeq\gamma\,\Theta_A(\xi,\xi)\,\abs{u}^2 -\langle\Theta_V(\xi,\xi)\cdot u,u\rangle
\\
\begin{aligned}
&+\sum_j\varepsilon_j\,\abs{\sigma_j}^{-2+2/\rho_j}\,
\big(\gamma\,\Theta_A(\xi,\xi)-\rho_j^{-1}\,\Theta_{\Delta_j}(\xi,\xi)\big)
\,\left|\nabla_j\sigma_j(u)\right|^2
\\
&+\sum_j\tfrac{\varepsilon_j\,\abs{\sigma_j}^{-2+2/\rho_j}}{ 1+\varepsilon_j\,\abs{\sigma_j}^{-2+2/\rho_j}\,
\left|\nabla_j\sigma_j\right|^2}\,\left|\,\nabla_j^2\sigma_j(\xi,u)-(1-1/\rho_j)\,
\sigma_j^{-1}\,\nabla_j\sigma_j(\xi)\nabla_j\sigma_j(u)\right|^2.
\end{aligned}
\end{multline}
Here, the symbol $\simeq$ means that the ratio of the left and right hand sides can be chosen in $[1-\alpha,1+\alpha]$ for any
$\alpha>0$ prescribed in advance.
\end{theo}

The next argument is the observation that the sheaf
$\cO_X(E_{k,m}V^*\langle D\rangle)$ is the direct image of a certain tautological rank $1$ sheaf $\cO_{X_k(V\langle D\rangle)}(m)$ on the ``orbifold $k$-jet bundle'' \text{$X_k(V\langle D\rangle)\to X$}. Choosing hermitian metrics according to Theorem~\ref{theo:orbi_curv_estimate}, one then gets a hermitian metric on $\cO_{X_k(V\langle D\rangle)}(1)$ associated with an ``orbifold Finsler metric'' on the bundle $J_kV$ of $k$-jets of holomorphic curves $f\colon(\bC,0)\to(X,V)$. In normalized coordinates
$(z_1,\,\dotsc,\,z_n)$ on $X$ and on $V$, the latter can be expressed as
\[
\left(
\sum_{s=1}^k\varepsilon_s^{2b}\left(\sum_{j=1}^p
\left|f_j\right|^{-2(1-s/\rho_j)_+}\left|f_j^{(s)}\right|^2+\sum_{j=p+1}^r\left|f_j^{(s)}\right|^2\right)^{2b/s}
\right)^{1/b},
\quad f\in J^kV,\; f(0)=x,
\]
at any point $x\in X$ where $\Delta_j=\{z_j=0\}$, $1\leq j\leq p$, $r=\rank V$. An application of holomorphic Morse inequalities (\cite{Dem85}, see also \S~\ref{se:Morse}, \ref{se:jet_metrics}, \ref{se:Morse_integrals}) then provides asymptotic estimates of the dimensions of the cohomology groups
\[
H^q\big(X,E_{k,m}V^*\langle D\rangle\otimes\cO_X(-A)\big)\simeq H^q(X_k(V\langle D\rangle),
\cO_{X_k(V\langle D\rangle)}(m)\otimes\pi_k^*\cO_X(-A)).
\]
This is done in several steps. Section~\ref{se:Morse_integrals} expresses the Morse integrals that need to be computed. Section~\ref{se:Chern_inequalities} establishes some general estimates of Chern forms related to the curvature tensor $\Theta_{E,h}$ of a given hermitian vector bundle $(E,h)$, under suitable positivity assumptions. More precisely, Proposition~\ref{prop:Chern_curv_inequalities} gives upper and lower bounds of integrals of the form
\[
\int_{u\,\in\,S(E)}\abs{\ell_1(u)}^2\ldots\abs{\ell_k(u)}^2\,
\langle\Theta_{E,h}(u),u\rangle_h^{p-k}\,d\mu(u)
\]
in terms of $\Tr_E\Theta_{E,h}=\Theta_{\det E,\det h}$, where $\mu$ is the unitary invariant probability measure on the unit sphere bundle $S(E)$, and the $\ell_j$ are linear forms. As far as we know, these estimates seem to be new. Sections~\ref{sse:Chern_integrals} and~\ref{se:non_proba_estimates} then proceed with the detailed calculations of the orbifold and logarithmic Morse integrals involved in the problem. It is remarkable that a large part of the calculations use Chern forms and are non cohomological, although the final bounds are purely cohomological. At~this point, we do not have a complete explanation of this ``transcendental'' phenomenon.

\section{Logarithmic and orbifold jet differentials}\label{se:jet_differentials}
\subsection{Directed varieties and associated jet differentials}
Let $(X,V)$ be a non-singular directed variety. We set $n\bydef\dim_\bC X$,
$r\bydef\rank_\bC V$, and following the exposition of~\cite{Dem97}, we denote by $\pi_k\colon J^kV\to X$ the bundle of $k$-jets of holomorphic curves tangent to
$V$ at each point. The canonical bundle of $V$ is defined to be
\[
K_V
\bydef
\det(V^*) =
\Lambda^rV^*.
\]
If $f\colon(\bC,0)\to X$, $t\mapsto f(t)$ is a germ of holomorphic curve tangent to $V$, we denote by $f_{[k]}(0)$ its $k$-jet at $t=0$. For $x_0\in X$ given, we take a coordinate system $(z_1,\dotsc,z_n)$ centered at $x_0$ such that
$V_{x_0}=\Span(\tfrac{\partial}{\partial z_\mu})_{1\,\leq \,\mu\,\leq\,r}$. Then there exists a neighborhood $U$ of $x_0$ such that
$V_{|U}$ admits a holomorphic frame $(e_\mu)_{1\,\leq\,\mu\,\leq\,r}$ of the form
\[
e_\mu(z)=\tfrac{\partial}{\partial z_\mu}+\sum_{r+1\,\leq \,\lambda\,\leq\,n} a_{\lambda\mu}(z)\tfrac{\partial}{\partial z_\lambda},\quad 1\leq\mu\leq r,
\]
with $a_{\lambda\mu}(0)=0$. Germs of curves $f\colon(\bC,0)\to X$ tangent to $V_{|U}$ are obtained by integrating the system of ordinary differential equations
\[
f'_\lambda(t)=\sum_{1\,\leq\,\mu\,\leq\,r}a_{\lambda\mu}(f(t))\,f'_\mu(t),\quad r+1\leq \lambda\leq n,
\]
when we write $f=(f_1,\,\dotsc,\,f_n)$ in coordinates. Therefore any such germ of curve $f$ is uniquely determined by its initial point $z=f(0)$ and its projection $\tilde f=(f_1,\,\dotsc,\,f_r)$ on the first $r$ coordinates. By definition, every $k$-jet $f_{[k]}\in J^kV_z=\pi_k^{-1}(z)$ is uniquely determined by its initial point $f(0)=z\simeq(z_1,\dotsc,z_n)$ and the Taylor expansion of order $k$
\begin{multline}\label{eq:Taylor}
\tilde f(t)-\tilde f(0)=t\xi_1+\tfrac{1}{ 2!}t^2\xi_2+\cdots+\tfrac{1}{ k!} t^k\xi_k+O(t^{k+1}),\\
 t\in\bD(0,\varepsilon),~\xi_s\in\bC^r,~1\leq s\leq k.
\end{multline}
Alternatively, we can pick an arbitrary local holomorphic connection $\nabla$ on $V_{|U}$ and represent the $k$-jet $f_{[k]}(0)$ by
$(\xi_1,\,\dotsc,\,\xi_k)$, where
$\xi_s=\nabla^sf(0)\in V_z$ is defined inductively by $\nabla^1 f=f'$ and $\nabla^sf=\nabla_{f'}(\nabla^{s-1}f)$. This gives a local biholomorphic trivialization of $J^kV_{|U}$ of the form
\begin{equation}\label{eq:triv}
J_kV_{|U}\to V_{|U}^{\oplus k},\qquad f_{[k]}(0)\mapsto(\xi_1,\,\dotsc,\,\xi_k)=\left(\nabla f(0),\,\dotsc,\,\nabla f^k(0)\right)\,;
\end{equation}
the particular choice of the ``trivial connection'' $\nabla_0$ of $V_{|U}$ that turns $(e_\mu)_{1\,\leq\,\mu\,\leq\,r}$ into a parallel frame precisely yields the components $\xi_s\in V_{|U}\simeq\bC^r$ appearing in~\eqref{eq:Taylor}. We could of course also use a $C^\infty$ connection $\nabla=\nabla_0+\Gamma$ where
$\Gamma\in C^\infty(U,T^*_X\otimes\Hom(V,V))$, and in this case, the corresponding trivialization~\eqref{eq:triv} is just a
$C^\infty$ diffeomorphism; the advantage, though, is that we can always produce such a global $C^\infty$ connection $\nabla$ by using a partition of unity on $X$, and then~\eqref{eq:triv} becomes a global $C^\infty$ diffeomorphism. Now, there is a global holomorphic $\bC^*$ action on $J^kV$ given at the level of germs by $f\mapsto\alpha\cdot f$ where $\alpha\cdot f(t)\bydef f(\alpha t)$,
$\alpha\in\bC^*$. With respect to our trivializations~\eqref{eq:triv}, this is the weighted $\bC^*$ action defined by
\[
\alpha\cdot(\xi_1,\xi_2,\dotsc,\xi_k)=\left(\alpha\xi_1,
\alpha^2\xi_2,\dotsc,\alpha^k\xi_k\right),\quad \xi_s\in V.
\]
We see that $J^kV\to X$ is an algebraic fiber bundle with typical fiber $\bC^{rk}$, and that the projectivized $k$-jet bundle
\[
X_k(V)
\bydef \left(J^kV\setminus\{0\}\right)/\bC^*,\qquad \pi_k:\smash{X_k(V)}\to X
\]
is a $\bP(1^{[r]},2^{[r]},\,\dotsc,\,k^{[r]})$ weighted projective bundle over $X$, of total dimension
\[
\dim X_k(V)=n+kr-1.
\]

\begin{defi}
We define $\cO_X(E_{k,m} V^*)$ to be the sheaf over $X$ of holomorphic functions $P(z\,;\,\xi_1,\,\dotsc,\,\xi_k)$ on $J^kV$ that are weighted polynomials of degree $m$ in $(\xi_1,\,\dotsc,\,\xi_k)$.
\end{defi}

In coordinates and in multi-index notation, we can write
\[
P(z\,;\,\xi_1,\,\dotsc,\,\xi_k)=
\sum_{\substack{\alpha_1,\,\dotsc,\,\alpha_k\in\bN^r\\
\abs{\alpha_1}+2\abs{\alpha_2}+\cdots+k\abs{\alpha_k}=m}} a_{\alpha_1\ldots\alpha_k}(z)\,\xi_1^{\alpha_1}\ldots\xi_k^{\alpha_k}
\]
where the $a_{\alpha_1\ldots\alpha_k}(z)$ are holomorphic functions in $z=(z_1,\,\dotsc,\,z_n)$ and $\xi_s^{\alpha_s}$ actually means
\[
\xi_s^{\alpha_s}=\xi_{s,1}^{\alpha_{s,1}}\ldots\,\xi_{s,r}^{\alpha_{s,r}}\quad
\text{for}
\xi_s=(\xi_{s,1},\,\dotsc,\,\xi_{s,r})\in\bC^r,~~
\alpha_s=(\alpha_{s,1},\,\dotsc,\,\alpha_{s,r})\in\bN^r,
\]
and $\abs{\alpha_s}=\sum_{j=1}^r\alpha_{s,j}$. Such sections can be interpreted as algebraic differential operators acting on holomorphic curves $f\colon\bD(0,R)\to X$ tangent to $V$, by putting
$P(f)\bydef u$ where
\[
u(t)=
\sum_{\substack{\alpha_1,\,\dotsc,\,\alpha_k\in\bN^r\\
\abs{\alpha_1}+2\abs{\alpha_2}+\cdots+k\abs{\alpha_k}=m}} a_{\alpha_1\ldots\alpha_k}(f(t))\;f'(t)^{\alpha_1}\ldots\,f^{(k)}(t)^{\alpha_k}.
\]
Here $f^{(s)}(t)^{\alpha_s}$ is actually to be expanded as
\[
f^{(s)}(t)^{\alpha_s}=f_1^{(s)}(t)^{\alpha_{s,1}}\ldots\,f_r^{(s)}(t)^{\alpha_{s,r}}
\]
with respect to the components $f_j^{(s)}$ defined in~\eqref{eq:Taylor}. We also set
\[
u = P\left(f\,;\,f',f'',\dotsc,f^{(k)}\right)
\]
when we want to make more explicit the dependence of the expression in terms of the derivatives of $f$. We thus get a sheaf of graded algebras
\[
\bigoplus_{m\,\in\,\bN}\cO_X\left(E_{k,m}V^*\right).
\]
Locally in coordinates, the algebra is isomorphic to the weighted polynomial ring
\[
\cO_X\left[f_j^{(s)}\right]_{1\,\leq\,j\,\leq\,r,\,1\,\leq\,s\,\leq\,k},\quad
\deg f_j^{(s)}=s
\]
over $\cO_X$. An immediate consequence of these definitions is:

\begin{prop}\label{prop:Xk}
The projectivized bundle $\pi_k\colon X_k(V)\to X$ can be identified with
\[
\tag{a}
\Proj\left(\bigoplus_{m\,\in\,\bN}\cO_X\left(E_{k,m}V^*\right)\right)\to X,
\]
and, if $\cO_{X_k(V)}(m)$ denote the associated tautological sheaves, we have the direct image formula
\[
\tag{b} (\pi_k)_*\smash{\cO_{X_k(V)}(m)}=\cO_X\left(E_{k,m} V^*\right).
\]
\end{prop}

\begin{rema}
These objects were denoted $X_k^{\text{GG}}$ and $E_{k,m}^{\text{GG}} V^*$ in our previous paper~\cite{Dem97}, as a reference to the work of Green--Griffiths~\cite{GG80}, but we will avoid here the superscript GG to simplify the notation.
\end{rema}

Thanks to the Faà di Bruno formula, a change of coordinates $w=\psi(z)$ on
$X$ leads to a transformation rule
\[
(\psi\circ f)^{(k)}=\psi'\circ f\cdot f^{(k)}+Q_\psi\left(f',\,\dotsc,\,f^{(k-1}\right)
\]
where $Q_\psi$ is a polynomial of weighted degree $k$ in the lower order derivatives. This shows that the transformation rule of the top derivative is linear and, as a consequence, the partial degree in $f^{(k)}$ of the polynomial $P(f\,;\,f',\dotsc,f^{k)})$ is intrinsically defined. By taking the corresponding filtration and factorizing the monomials
$(f^{(k)})^{\alpha_k}$ with polynomials in $f',f'',\dotsc,f^{(k-1)}$, we get graded pieces
\[
G^\bullet\left(E_{k,m}V^*\right)=\bigoplus_{\ell_k\,\in\,\bN} E_{k-1,m-k\ell_k}V^*\otimes S^{\ell_k}V^*.
\]
By considering successively the partial degrees with respect to $f^{(k)}, f^{(k-1)},\,\dotsc,\,f'', f'$ and merging inductively the resulting filtrations, we get a multi-filtration such that
\begin{equation}\label{eq:graded_bundle}
G^\bullet\left(E_{k,m}V^*\right)=\bigoplus_{\ell_1,\,\dotsc,\,\ell_k\in\bN,\,
\ell_1+2\ell_2+\cdots+k\ell_k=m}S^{\ell_1}V^*\otimes S^{\ell_2}V^*\otimes\cdots
\otimes S^{\ell_k}V^*.
\end{equation}

\subsection{Logarithmic directed varieties}\label{sse:log_directed}
We now turn ourselves to the logarithmic case. Let $(X,V,D)$ be a non-singular logarithmic variety, where $D=\sum\Delta_j$ is a simple normal crossing divisor. Fix a point $x_0\in X$. By the assumption that
$D$ is transverse to $V$, we can then select holomorphic coordinates
$(z_1,\,\dotsc,\,z_n)$ centered at $x_0$ such that
$V_{x_0}=\Span(\tfrac{\partial}{\partial z_j})_{1\,\leq\,j\,\leq\,r}$ and $\Delta_j=\{z_j=0\}$, $1\leq j\leq p$, are the components of $D$ that contain $x_0$ (here $p\leq r$ and we can have $p=0$ if $x_0\notin\abs{D}$). What we want is to introduce an algebra of differential operators, defined locally near $x_0$ as the weighted polynomial ring
\begin{equation}\label{eq:diffs}
\cO_X\left[(\log f_j)^{(s)}_{1\,\leq\,j\,\leq \,p}\,,\left(f_j^{(s)}\right)_{p+1\,\leq\,j\,\leq\,r}
\right]_{1\,\leq\,s\,\leq \,k},\quad \deg f_j^{(s)}=\deg\left(\log f_j\right)^{(s)}=s,
\end{equation}
or equivalently
\begin{equation'}\label{eq:diffs'}
\cO_X\left[\left(f_j^{-1}f_j^{(s)}\right)_{1\,\leq\,j\,\leq\,p}\,,\left(f_j^{(s)}\right)_{p+1\,\leq\,j\,\leq\,r}
\right]_{1\,\leq\,s\,\leq\,k},\quad \deg f_j^{(s)}=s,~\deg f_j^{-1}=0.
\end{equation'}
For this we notice that
\begin{align*}
(\log f_1)''&=\left(f_1^{-1}f_1'\right)'=f_1^{-1}f_1''-\left(f_1^{-1}f_1'\right)^2,
\\
(\log f_1)'''&=f_1^{-1}f_1'''-3\left(f_1^{-1}f_1'\right)\left(f_1^{-1}f_1''\right)+2 \left(f_1^{-1}f_1'\right)^3\ldots\,.
\end{align*}
A similar argument easily shows that the above graded rings do not depend on the particular choice of coordinates made, as soon as they satisfy
$\Delta_j=\{z_j=0\}$.

Now (as is well known in the absolute case $V=T_X$), we have a corresponding logarithmic directed structure
$V\langle D\rangle$ and its dual $V^*\langle D\rangle$. If the coordinates $(z_1,\,\dotsc,\,z_n)$ are chosen so that $V_{x_0}=\{dz_{r+1}=\ldots=dz_n=0\}$, then the fiber $V\langle D\rangle_{x_0}$ is spanned by the derivations
\[
z_1\tfrac{\partial}{\partial z_1},\,\dotsc,\,z_p\tfrac{\partial}{\partial z_p},
\tfrac{\partial}{\partial z_{p+1}},\,\dotsc,\,\tfrac{\partial}{\partial z_r}.
\]
The dual sheaf $\cO_X(V^*\langle D\rangle)$ is the locally free sheaf generated by
\[
\tfrac{dz_1}{ z_1},\,\dotsc,\,\tfrac{dz_p}{ z_p}, dz_{p+1},\dotsc,dz_r
\]
[where the $1$-forms are considered in restriction to
$\cO_X(V\langle D\rangle)\subset\cO_X(V)\,$]. It follows from this that $\cO_X(V\langle D\rangle)$ and
$\cO_X(V^*\langle D\rangle)$ are locally free sheaves of rank $r$. By taking $\det(V^*\langle D\rangle)$ and using the above generators, we find
\[
\det\left(V^*\langle D\rangle\right)=\det(V^*)\otimes\cO_X(D)=K_V+ D
\]
in additive notation. Quite similarly to Props.~\ref{prop:Xk} and~\ref{eq:graded_bundle}, we have:
\begin{prop}\label{prop2.4}
Let
\[
\bigoplus_{m\,\in\,\bN}\cO_X\left(E_{k,m}V^*\langle D\rangle\right)
\]
be the graded algebra defined in coordinates by~\eqref{eq:diffs} or~\eqref{eq:diffs'}. We define the logarithmic $k$-jet bundle to be
\begin{equation}\label{prop2.4.a}
\tag{a} X_k(V\langle D\rangle)
\bydef
\Proj\left(\bigoplus_{m\,\in\,\bN}\cO_X(E_{k,m}V^*\langle D\rangle)\right)\to X.
\end{equation}
If $\cO_{X_k(V\langle D\rangle)}(m)$ denote the associated tautological sheaves, we get the direct image formula
\begin{equation}\label{prop2.4.b}
\tag{b} (\pi_k)_*\smash{\cO_{X_k(V\langle D\rangle)}(m)}=\cO_X\left(E_{k,m} V^*
\langle D\rangle\right).
\end{equation}
Moreover, the multi-filtration by the partial degrees in the derivatives
$f_j^{(s)}$ has graded pieces
\begin{equation}\label{prop2.4.c}
\tag{c} G^\bullet\big(E_{k,m}V^*\langle D\rangle\big) =
\bigoplus_{\substack{\ell_1,\,\dotsc,\,\ell_k\,\in\,\bN\\\ell_1+2\ell_2+\cdots+k\ell_k=m}} S^{\ell_1}V^*\langle D\rangle\otimes S^{\ell_2}V^*\langle D\rangle\otimes\cdots\otimes S^{\ell_k}V^*\langle D\rangle.
\end{equation}
\end{prop}

\subsection{Orbifold directed varieties}

We finally consider a non-singular directed orbifold $(X,V,D)$, where $D=\sum(1-\tfrac{1}{\rho_j})\Delta_j$ is a simple normal crossing divisor transverse to $V$. Let $\lceil D\rceil=\sum\Delta_j$ be the corresponding reduced divisor. By \S~\ref{sse:log_directed}, we have associated logarithmic sheaves
$\cO_X(E_{k,m}V^*\langle\lceil D\rceil\rangle)$. We want to introduce a graded subalgebra
\[
\bigoplus_{m\,\in\,\bN}\cO_X\left(E_{k,m}V^*\langle D\rangle\right)
\subseteq
\bigoplus_{m\,\in\,\bN}\cO_X\left(E_{k,m}V^*\langle\lceil D\rceil\rangle\right)
\]
in such a way that for every germ $P\in \cO_X(E_{k,m}V^*\langle D\rangle)$ and every germ of orbifold curve $f\colon(\bC,0)\to(X,V,D)$ the germ of meromorphic function $P(f)(t)$ is bounded at $t=0$ (hence holomorphic). Assume that $\Delta_1=\{z_1=0\}$ and that $f$ has multiplicity $q\geq \rho_1>1$ along $\Delta_1$ at $t=0$. Then $f_1^{(s)}$ still vanishes at order${}\geq(q-s)_+$, thus $(f_1)^{-\beta}f_1^{(s)}$ is bounded as soon as $\beta q\leq(q-s)_+$, i.e.\ $\beta\leq(1-\tfrac{s}{ q})_+$. Thus, it is sufficient to ask that $\beta\leq(1-\tfrac{s}{ \rho_1})_+$. At a point $x_0\in \abs{\Delta_1}\cap\ldots\cap\abs{\Delta_p}$, a sufficient condition for a monomial of the form
\begin{equation}\label{eq:jet_monomial}
f_1^{-\beta_1}\ldots\,f_p^{-\beta_p}\prod_{s=1}^k\prod_{j=1}^r \left(f_j^{(s)}\right)^{\alpha_{s,j}},
\quad
\alpha_s=(\alpha_{s,j})\in\bN^r,~\beta_1,\,\dotsc,\,\beta_p\in\bN
\end{equation}
to be bounded is to require that the multiplicities of poles satisfy
\begin{equation'}\label{eq:poles}
\beta_j
\leq
\sum_{s=1}^k
\alpha_{s,j}\Big(1-\tfrac{s}{\rho_j}\Big)_+,
\quad 1\leq j\leq p.
\end{equation'}
\begin{defi}
The subalgebra $\bigoplus_{m\,\in\,\bN}\cO_X(E_{k,m}V^*\langle D\rangle)$ is taken to be the graded ring generated by monomials~\eqref{eq:jet_monomial} of degree
$\sum s\abs{\alpha_s}=m$, satisfying the pole multiplicity conditions~\eqref{eq:poles}. These conditions do not depend on the choice of coordinates, hence we get a globally and intrinsically defined sheaf of algebras on $X$.
\end{defi}

\begin{proof}
We only have to prove the last assertion. Consider a change of variables
$w=\psi(z)$ such that $\Delta_j$ can still be expressed as
$\Delta_j=\{w_j=0\}$. Then, for $j=1,\,\dotsc,\,p$, we can write
$w_j=z_ju_j(z)$ with an invertible holomorphic factor $u_j$. We need to check that the monomials~\eqref{eq:jet_monomial} computed with $g=\psi\circ f$ are holomorphic combinations of those associated with $f$. However, we have
$g_j=f_ju_j(f)$, hence $g_j^{(s)}=\sum_{0\,\leq\,\ell\,\leq\,s}\binom{s}{\ell} f_j^{(\ell)}(u_j(f))^{(s-\ell)}$ by the Leibniz formula, and we see that
\[
g_1^{-\beta_1}\,\ldots\,g_p^{-\beta_p}\prod_{s=1}^k\prod_{j=1}^r \left(g_j^{(s)}\right)^{\alpha_{s,j}}
\]
expands as a linear combination of monomials
\[
f_1^{-\beta_1}\ldots\,f_p^{-\beta_p}\prod_{s=1}^k\prod_{j=1}^r
\prod_{m=1}^{\alpha_{s,j}}f_j^{(\ell_{s,j,m})},\quad \ell_{s,j,m}\leq s,
\]
multiplied by holomorphic factors of the form
\[
\prod_{j=1}^pu_j(f)^{-\beta_j}\times
\prod_{s=1}^k\prod_{j=1}^r\prod_{m=1}^{\alpha_{s,j}}\left(u_j(f)\right)^{\left(s-\ell_{j,s,m}\right)}.
\]
However, we have
\[
\beta_j\leq\sum_{s=1}^k\alpha_{s,j}\left(1-\tfrac{s}{ \rho_j}\right)_+
\leq\sum_{s=1}^k\sum_{m=1}^{\alpha_{s,j}}
\left(1-\tfrac{\ell_{s,j,m}}{ \rho_j}\right)_+,\quad
\]
so the $f$-monomials satisfy again the required multiplicity conditions for the poles $f_j^{-\beta_j}$.
\end{proof}

The above conditions~\eqref{eq:poles} suggest to introduce as in~\cite{CDR20} a sequence of ``differentiated'' orbifold divisors
\[
D^{(s)}=\sum_j\left(1-\tfrac{s}{\rho_j}\right)_{\kern-3pt+}\Delta_j.
\]
We say that $D^{(s)}$ is the order $s$ orbifold divisor associated to $D$; its ramification numbers are $\rho_j^{(s)}=\max(\rho_j/s,1)$. By definition, the logarithmic components ($\rho_j=\infty$) of $D$ remain logarithmic in $D^{(s)}$, while all others eventually disappear when $s$ is large.

Now, we introduce (in a purely formal way) a sheaf of rings $\smash{\widetilde\cO}_X=\cO_X[z_j^\bullet]$ by adjoining all positive real powers of coordinates $z_j$ such that \text{$\Delta_j=\{z_j=0\}$} is locally a component of $D$. Locally over $X$, this can be done by taking the universal cover $Y$ of a punctured polydisk
\[
\bD^*(0,r)
\bydef
\prod_{1\,\leq\,j\,\leq\,p}\bD^*(0,r_j)\times
\prod_{p+1\,\leq\,j\,\leq\,n}\bD(0,r_j)~~\subset~~
\bD(0,r)
\bydef
\prod_{1\,\leq\,j\,\leq\,n}\bD(0,r_j)
\]
in the local coordinates $z_j$ on $X$. If $\gamma\colon Y\to\bD^*(0,r) \hookrightarrow X$ is the covering map and $U\subset\bD(0,r)$ is an open subset, we can then consider the functions of $\smash{\widetilde\cO}_X(U)$ as being defined on $\gamma^{-1}(U\cap\bD^*(0,r))$. In case $X$ is projective, one can even achieve such a construction ``globally'', at least on a Zariski open set, by taking $Y$ to be the universal cover of a complement $X\setminus(\abs{D}\cup\abs{A})$, where
$A=\sum A_j$ is a very ample normal crossing divisor transverse to $D$, such that $\cO_X(\Delta_j)_{|X\setminus\abs{A}}$ is trivial for every $j\,$; then
$\smash{\widetilde\cO}_X$ is well defined as a genuine sheaf on $X\setminus\abs{A}$.

In this setting, the subalgebra $\bigoplus_m\cO_X(E_{k,m}V^*\langle D\rangle)$ still has a multi-filtration induced by the one on $\bigoplus_m\cO_X(E_{k,m}V^*\langle\lceil D\rceil\rangle)$, and by extending the structure sheaf $\cO_X$ into $\smash{\widetilde\cO}_X$, we get an inclusion
\[
\widetilde\cO_X\left(G^\bullet E_{k,m}V^*\langle D\rangle\right)\subset
\bigoplus_{\ell_1+2\ell_2+\cdots+k\ell_k=m}
\widetilde\cO_X\left(S^{\ell_1}V^*\left\langle D^{(1)}\right\rangle\right)\otimes \cdots
\otimes\widetilde\cO_X\left(S^{\ell_k}V^*\left\langle D^{(k)}\right\rangle\right),
\]
$\widetilde\cO_X(V^*\langle D^{(s)}\rangle)$ is the ``$s^{\rm th}$ orbifold (dual) directed structure'', generated by the order
$s$ differentials
\[
z_j^{-(1-s/\rho_j)_+}d^{(s)}z_j,~~1\leq j\leq p,~~~d^{(s)}z_j,~~p+1\leq j\leq r.
\]
By construction, we have
\begin{equation}\label{eq:canonical}
\det\left(\widetilde\cO_X\left(V^*\left\langle D^{(s)}\right\rangle\right)\right)=\widetilde\cO_X\left(K_V+ D^{(s)}\right).
\end{equation}

\begin{rema}\label{rema:q_coefficients}
When $\rho_j=a_j/b_j\in\bQ_+$, one can find a finite ramified Galois cover
\text{$g\colon Y\to X$} from a smooth projective variety $Y$ onto $X$, such that the compositions $(z_j\circ g)^{1/a_j}$ become single-valued functions
$w_j$ on $Y$. In this way, the pull-back $\cO_Y(g^*V^*\langle D^{(s)}\rangle)$ is actually a locally free $\cO_Y$-module. On can also introduce a sheaf of algebras which we will denote by $\bigoplus\cO_Y(E_{k,m}\widetilde V^*\langle D\rangle)$, generated, according to the notation of \S~\ref{sse:log_directed}, by the elements $g^*(z_j^{(1-s/\rho_j)_+}d^{(s)}z_j)$, $1\leq j\leq p$, and
$g^*(d^{(s)}z_j)$, $p+1\leq j\leq r$. Then, as already shown in~\cite{CDR20}, there is indeed a multifiltration on
$\cO_Y(E_{k,m}\widetilde V^*\langle D\rangle)$ whose graded pieces are
\[
\cO_Y\left(G^\bullet E_{k,m}\widetilde V^*\langle D\rangle\right)=
\bigoplus_{\ell_1+2\ell_2+\cdots+k\ell_k=m}
\cO_Y\left(S^{\ell_1}\widetilde V^*\left\langle D^{(1)}\right\rangle\right)\otimes \cdots
\otimes\cO_Y\left(S^{\ell_k}\widetilde V^*\left\langle D^{(k)}\right\rangle\right).
\]
However, we will adopt here an alternative viewpoint that avoids the introduction of finite or infinite covers, and suits better our approach. The general philosophy is to consider a ``jet orbifold directed structure'' $X_k(V\langle D\rangle)$ as the underlying ``jet logarithmic directed structure'' $X_k(V\langle\lceil D\rceil\rangle)$, equipped additionally with a submultiplicative sequence of ideal sheaves
$\cJ_m\langle D\rangle\subset \cO_{X_k(V\langle\lceil D\rceil\rangle)}$. These are precisely defined as the base loci ideals of the local sections defined by~\eqref{eq:jet_monomial} and~\eqref{eq:poles}, seen as sections of the logarithmic tautological sheaves
$\cO_{X_k(V\langle\lceil D\rceil\rangle)}(m)$. The corresponding analytic viewpoint is to consider ad hoc singular hermitian metrics on $\cO_{X_k(V\langle\lceil D\rceil\rangle)}(1)$ whose singularities are asymptotically described by the limit of the formal $m^{\rm th}$ root of $\cJ_m\langle D\rangle$, see \S~\ref{sse:orbi_jet_metrics}. It then becomes possible to deal without trouble with real coefficients $\rho_j\in{}]1,\infty]$, and since we no longer have to worry about the existence of Galois covers, the projectivity assumption on $X$ can be dropped as well.
\end{rema}

\section{Preliminaries on holomorphic Morse inequalities}\label{se:Morse}
\subsection{Basic results} We first recall the basic results concerning holomorphic Morse inequalities for smooth hermitian line bundles, first proved in~\cite{Dem85}.
\begin{theo}\label{theo:Morse}
Let $X$ be a compact complex manifolds, $E\to X$ a holomorphic vector bundle of rank $r$, and $(L,h)$ a hermitian line bundle. We denote by
$\Theta_{L,h}=\tfrac{\imath}{ 2\pi}\nabla_h^2=-\tfrac{\imath}{\pi}\ddbar\log h$ the curvature form of $(L,h)$ and introduce the open subsets of $X$
\[
\tag{$*$}
\begin{cases}
X(L,h,q)=\big\{x\in X\,;\;\Theta_{L,h}(x)~\text{has signature $(n-q,q)$}\big\},
\\
\displaystyle X(L,h,S)=\bigcup_{q\,\in\,S} X(L,h,q),\quad \forall S\subset\{0,1,\,\dotsc,\,n\}.
\end{cases}
\]
Then, for all $q=0,1,\,\dotsc,\,n$, the dimensions $h^q(X,E\otimes L^m)$ of cohomology groups of the tensor powers $E\otimes L^m$ satisfy the following ``Strong Morse inequalities'' as $m\to +\infty$:
\[
\tag*{$\SM(q)$:}
\sum_{0\,\leq\,j\,\leq\,q} (-1)^{q-j}h^j\left(X,E\otimes L^m\right) \leq r \tfrac{m^n}{ n!}
\int_{X(L,h,\,\leq\,q)}(-1)^q\Theta_{L,h}^n+o(m^n),
\]
with equality
$\chi(X,E\otimes L^m)= r\tfrac{m^n}{ n!}\int_X \Theta_{L,h}^n + o(m^n)$ for the Euler characteristic $(q=n)$.
\end{theo}

As a consequence, one gets upper and lower bounds for all cohomology groups, and especially a very useful criterion for the existence of sections of large multiples of $L$.

\begin{coro}
Under the above hypotheses, we have
\begin{enumerate}\alphenumi
\item\label{coro3.2.a} Upper bound for $h^q$ (Weak Morse inequalities):
\[
h^q(X,E\otimes L^m)\leq r \tfrac{m^n}{ n!}\int_{X(L,h,q)} (-1)^q \Theta_{L,h}^n + o(m^n).
\]
\item\label{coro3.2.b} Lower bound for $h^0$:
\[
h^0(X,E\otimes L^m)\geq h^0-h^1\ge r\tfrac{m^n}{ n!}\int_{X(L,h,\leq 1)}\Theta_{L,h}^n -o(m^n).
\]
Especially $L$ is big as soon as $\int_{X(L,h,\leq 1)}\Theta_{L,h}^n>0$ for some hermitian metric $h$ on $L$.
\item\label{coro3.2.c} Lower bound for $h^q$:
\[
h^q(X,E\otimes L^m)\geq h^q-h^{q-1}-h^{q+1}\ge r\tfrac{m^n}{ n!}\int_{X(L,h,\{q,q\pm 1\})} (-1)^q \Theta_{L,h}^n + o(m^n).
\]
\end{enumerate}
\end{coro}
\begin{proof}
\eqref{coro3.2.a} is obtained by taking $\SM(q)+\SM(q\,{-}\,1)$, \eqref{coro3.2.b} is equivalent to
$-\SM(1)$ and~\eqref{coro3.2.c} is equivalent to $-(\SM(q\,{+}\,1)+\SM(q\,{-}\,2))$.
\end{proof}

The following simple lemma is the key to derive algebraic Morse inequalities from their analytic form (cf.~\cite[Theorem~12.3]{Dem94}).
\begin{lemm}\label{lemm:forms_ineq}
Let $\eta=\alpha-\beta$ be a difference of semipositive
$(1,1)$-forms on an $n$-dimensional complex manifold $X$, and let $\bOne_{\eta,\leq\,q}$ be the characteristic function of the open set where $\eta$ is non-degenerate with a number of negative eigenvalues at most equal to $q$. Then
\[
(-1)^q\bOne_{\eta,\,\leq\,q}~\eta^n\leq \sum_{0\,\leq\,j\,\leq\,q}(-1)^{q-j}
\binom{n}{ j}\alpha^{n-j}\wedge\beta^j,
\]
in particular
\[
\bOne_{\eta,\,\leq\,1}~\eta^n\geq \alpha^n-n\alpha^{n-1}\wedge \beta\qquad\text{for $q=1$.}
\]
\end{lemm}

\begin{proof}
Without loss of generality, we can assume $\alpha>0$ positive definite, so that
$\alpha$ can be taken as the base hermitian metric on $X$. Let us denote by
\[
\lambda_1\geq\lambda_2\geq\ldots\geq\lambda_n\geq 0
\]
the eigenvalues of $\beta$ with respect to $\alpha$. The eigenvalues of $\eta=\alpha-\beta$ are then given by
\[
1-\lambda_1\leq\ldots\leq 1-\lambda_q\leq 1-\lambda_{q+1}\leq\ldots\leq 1-\lambda_n,
\]
hence the open set $\{\lambda_{q+1}<1\}$ coincides with the support of
$\bOne_{\eta,\,\leq\,q}$, except that it may also contain a part of the degeneration set $\eta^n=0$. On the other hand we have
\[
\binom{n}{ j}\alpha^{n-j}\wedge\beta^j =
\sigma_n^j(\lambda)\,\alpha^n,
\]
where $\sigma_n^j(\lambda)$ is the $j^{\rm th}$ elementary symmetric function in the $\lambda_j$'s. Thus, to prove the lemma, we only have to check that
\[
\sum_{0\,\leq\,j\,\leq\,q}(-1)^{q-j}\sigma_n^j(\lambda)-
\bOne_{\left\{\lambda_{q+1}<1\right\}}(-1)^q\prod_{1\,\leq\,j\,\leq\,n}(1-\lambda_j)\geq 0.
\]
This is easily done by induction on $n$ (just split apart the parameter
$\lambda_n$ and write
$\sigma_n^j(\lambda) =
\sigma_{n-1}^j(\lambda)+\sigma_{n-1}^{j-1}(\lambda)\,\lambda_n
$).
\end{proof}

\begin{coro}\label{coro:Morse}
Assume that $\eta=\Theta_{L,h}$ can be expressed as a difference $\eta=\alpha-\beta$ of smooth $(1,1)$-forms $\alpha,\beta\geq 0$. Then we have
\[
\tag*{$\SM(q)$:}
\sum_{0\,\leq\,j\,\leq\,q} (-1)^{q-j}h^j(X,E\otimes L^m) \leq r \tfrac{m^n}{ n!}
\int_X\sum_{0\,\leq\,j\,\leq\,q}(-1)^{q-j}\binom{n}{ j}\alpha^{n-j}\wedge\beta^j+o(m^n),
\]
and in particular, for $q=1$,
\[
h^0(X,E\otimes L^m)\geq h^0-h^1\ge r\tfrac{m^n}{ n!}\int_X\alpha^n-n\alpha^{n-1}\wedge\beta+o(m^n).
\]
\end{coro}

\begin{rema}
These estimates are consequences of Theorem~\ref{theo:Morse} and Lemma~\ref{lemm:forms_ineq}, by taking the integral over $X$. The estimate for $h^0$ was stated and studied by Trapani~\cite{Tra95}. In the special case $\alpha=\Theta_{A,h_A}>0$,
$\beta=\Theta_{B,h_B}>0$ where $A,B$ are ample line bundles, a direct proof can be obtained by purely algebraic means, via the Riemann--Roch formula. However, we will later have to use Corollary~\ref{coro:Morse} in case $\alpha$ and $\beta$ are not closed, a situation in which no algebraic proof seems to exist.
\end{rema}


\subsection{Singular holomorphic Morse inequalities}
The case of singular hermitian metrics has been considered in Bonavero's PhD thesis~\cite{Bon93} and will be important for us. We assume that $L$ is equipped with a singular hermitian metric $h=h_\infty e^{-\varphi}$ with analytic singularities, i.e., $h_\infty$ is a smooth metric, and on an neighborhood $V\ni x_0$ of an arbitrary point $x_0\in X$, the weight
$\varphi$ is of the form
\begin{equation}\label{eq:weight}
\varphi(z) = c\log\sum_{1\,\leq\,j\,\leq\,N}\abs{g_j}^2+u(z)
\end{equation}
where $g_j\in\cO_X(V)$ and $u\in C^\infty(V)$. We then have
$\Theta_{L,h}=\alpha+\tfrac{\imath}{ 2\pi}\ddbar\varphi$ where
$\alpha=\Theta_{L,h_\infty}$ is a smooth closed $(1,1)$-form on $X$. In this situation, the multiplier ideal sheaves
\[
\cI(h^m)=\cI(k\varphi)=\left\{f\in\cO_{X,x},\;\;\exists V\ni x,~
\int_V\abs{f(z)}^2e^{-m\varphi(z)}d\lambda(z)<+\infty\right\}
\]
play an important role. We define the singularity set of $h$ by
$\Sing(h)=\Sing(\varphi)=\varphi^{-1}(-\infty)$ which, by definition, is an analytic subset of $X$. The associated $q$-index sets are
\[
X(L,h,q) =
\big\{ x\in X\setminus\Sing(h)
\,;\;
\Theta_{L,h}(x)\text{ has signature }(n-q,q)
\big\}.
\]
We can then state:
\begin{theo}[\cite{Bon93}]\label{theo:Bonavero}
Morse inequalities still hold in the context of singular hermitian metric with analytic singularities, provided the cohomology groups under consideration are twisted by the appropriate multiplier ideal sheaves, i.e.\ replaced by $H^q(X,E\otimes L^m\otimes\cI(h^m))$.
\end{theo}

\begin{rema}\label{rema:integrability}
The assumption~\eqref{eq:weight} guarantees that the measure
\[
\bOne_{X\setminus\Sing(h)}\left(\Theta_{L,h}\right)^n
\]
is locally integrable on $X$, as is easily seen by using the Hironaka desingularization theorem and by taking a log resolution $\mu\colon\widetilde X\to X$ such that
$\mu^*(g_j)=(\gamma)\subset \cO_{\smash{\widetilde X}}$ becomes a principal ideal associated with a simple normal crossing divisor
$E=\operatorname{div}(\gamma)$. Then $\mu^*\Theta_{L,h}= c[E]+\beta$ where $\beta$ is a smooth closed
$(1,1)$-form on $\widetilde X$, hence
\[
\mu^*\left(\bOne_{X\setminus\Sing(h)}\Theta_{L,h}^n\right)=\beta^n~~\Rightarrow~~
\int_{X\setminus\Sing(h)}\Theta_{L,h}^n=\int_{\widetilde X}\beta^n.
\]
It should be observed that the multiplier ideal sheaves $\cI(h^m)$ and the integral $\int_{X\setminus\Sing(h)}\Theta_{L,h}^n$ only depend on the equivalence class of singularities of $h$: if we have two metrics with analytic singularities $h_j=h_\infty e^{-\varphi_j}$, $j=1,2$, such that
$\psi=\varphi_2-\varphi_1$ is bounded, then, with the above notation, we have $\mu^*\Theta_{L,h_j}=c[E]+\beta_j$ and $\beta_2=\beta_1+\tfrac{\imath}{ 2\pi}\ddbar\psi$, therefore $\int_{\widetilde X}\beta_2^n=\int_{\widetilde X}\beta_1^n$ by Stokes theorem. By using Monge--Ampère operators in the sense of Bedford--Taylor~\cite{BT76}, it is in fact enough to assume $u\in L^\infty_\loc(X)$ in~\eqref{eq:weight}, and $\psi\in L^\infty(X)$ here. In general, however, the Morse integrals
$\int_{X(L,h_j,q)}(-1)^q\Theta_{L,h_j}^n$, $j=1,2$, will differ.
\end{rema}


\section{Construction of jet metrics and orbifold jet metrics}\label{se:jet_metrics}
\subsection{Jet metrics and curvature tensor of jet bundles}\label{sse:jet_metrics}

Let $(X,V)$ be a non-singular directed variety and $h$ a hermitian metric on $V$. We assume that $h$ is smooth at this point (but will later relax a little bit this assumption and allow certain singularities). Near any given point $z_0\in X$, we can choose local coordinates
$z=(z_1,\,\dotsc,\,z_n)$ centered at $z_0$ and a local holomorphic coordinate frame
$(e_\lambda(z))_{1\,\leq\,\lambda\,\leq\,r}$ of $V$ on an open set $U\ni z_0$, such that
\begin{equation}\label{eq:frame}
\langle e_\lambda(z),e_\mu(z)\rangle_{h(z)} =\delta_{\lambda\mu}+
\sum_{1\,\leq\,i,\,j\,\leq\,n,\,1\,\leq\,\lambda,\,\mu\,\leq\,r}c_{ij\lambda\mu}z_i\bar z_j+ O(\abs{z}^3)
\end{equation}
for suitable complex coefficients $(c_{ij\lambda\mu})$. It is a standard fact that such a normalized coordinate system always exists, and that the Chern curvature tensor $\tfrac{\imath}{ 2\pi}\nabla^2_{V,h}$ of $(V,h)$ at $z_0$ is given by
\[
\Theta_{V,h}(z_0)=-\tfrac{\imath}{ 2\pi}
\sum_{i,j,\lambda,\mu} c_{ij\lambda\mu}\,dz_i\wedge d\bar z_j\otimes e_\lambda^*\otimes e_\mu.
\]
Therefore, $(\tfrac{\imath}{ 2\pi}\,c_{ij\lambda\mu})$ are the coefficients of
$-\Theta_{V,h}$. Up to taking the transposed tensor with respect to $\lambda,\mu$, these coefficients are also the components of the curvature tensor
$\Theta_{V^*,h^*}=-{}^t\Theta_{V,h}$ of the dual bundle $(V^*,h^*)$. By~\eqref{eq:triv}, the connection $\nabla=\nabla_h$ yields a $C^\infty$ isomorphism $J_kV\to V^{\oplus k}$. Let us fix an integer $b\in\bN^*$ that is a multiple of $\lcm(1,2,\,\dotsc,\,k)$, and positive numbers
$1=\varepsilon_1\gg\varepsilon_2\gg\cdots\gg \varepsilon_k>0$. Following~\cite{Dem11}, we define a global weighted Finsler metric on $J^kV$ by putting for any $k$-jet $f\in J^kV_z$
\[
\Psi_{h,b,\varepsilon}(f)
\bydef
\left(
\sum_{1\,\leq\,s\,\leq\,k}
\varepsilon_s^{2b}
\Abs{\nabla^s f(0)}_{h(z)}^{2b/s}
\right)^{1/b},
\]
where $\Abs{\cdot}_{h(z)}$ is the hermitian metric $h$ of $V$ evaluated on the fiber $V_z$, $z=f(0)$. The function $\Psi_{h,b,\varepsilon}$ satisfies the fundamental homogeneity property
\[
\Psi_{h,b,\varepsilon}(\alpha\cdot f) =
\abs{\alpha}^2\,\Psi_{h,b,\varepsilon}(f)
\]
with respect to the $\bC^*$ action on $J^kV$, in other words, it induces a hermitian metric on the dual $L_k^*$ of the tautological $\bQ$-line bundle
$L_k=\cO_{X_k(V)}(1)$ over $X_k(V)$. The curvature of $L_k$ is given by
\begin{equation}\label{eq:curvature}
\pi_k^*\Theta_{L_k,\Psi^*_{h,b,\varepsilon}}=\tfrac{\imath}{ 2\pi}\ddbar
\log\Psi_{h,b,\varepsilon}
\end{equation}

Our next goal is to compute precisely the curvature and to apply holomorphic Morse inequalities to $L\to X_k(V)$ with the above metric. This might look a priori like an untractable problem, since the definition of $\Psi_{h,b,\varepsilon}$ is a rather complicated one, involving the hermitian metric in an intricate manner. However, the ``miracle'' is that the asymptotic behavior of $\Psi_{h,b,\varepsilon}$ as $\varepsilon_s/\varepsilon_{s-1}\to 0$ is in some sense uniquely defined, and ``splits'' according to the natural multifiltration on jet differentials. This leads to a computable asymptotic formula, which is moreover simple enough to produce useful results.

\begin{lemm}\label{lemm:rescaling}
Let us consider the global $C^\infty$ bundle isomorphism
$J^kV\to V^{\oplus k}$ associated with an arbitrary global $C^\infty$ connection
$\nabla$ on $V\to X$, and let us introduce the rescaling transformation
\[
\rho_{\nabla,\varepsilon}(\xi_1,\xi_2,\dotsc,\xi_k)= \left(\varepsilon_1^1\xi_1,\varepsilon_2^2\xi_2,\dotsc,
\varepsilon_k^k\xi_k\right)\quad
\text{on fibers $J^kV_z$, $z\in X$}.
\]
Such a rescaling commutes with the $\bC^*$-action. Moreover, if $p$ is a multiple of $\lcm(1,2,\,\dotsc,\,k)$ and the ratios
$\varepsilon_s/\varepsilon_{s-1}$ tend to $0$ for all $s=2,\dotsc,k$, the rescaled Finsler metric
\text{$\Psi_{h,b,\varepsilon}\circ\rho_{\nabla,\varepsilon}^{-1} (\xi_1,\,\dotsc,\,\xi_k)$} converges towards the limit
\[
\left(\sum_{1\,\leq\,s\,\leq\,k}
\Abs{\xi_s}^{2b/s}_h\right)^{1/b}
\]
on every compact subset of $V^{\oplus k}\setminus\{0\}$, uniformly in $C^\infty$ topology, and the limit is independent of the connection $\nabla$. The error is measured by a multiplicative factor
$1\pm O(\max_{2\,\leq\,s\,\leq\,k}(\varepsilon_s/\varepsilon_{s-1})^s)$.
\end{lemm}

\begin{proof}
Let us pick another $C^\infty$ connection $\widetilde\nabla= \nabla+\Gamma$ where $\Gamma\in C^\infty(U,T^*_X\otimes \Hom(V,V))$. Then
$\widetilde\nabla^2f=\nabla^2f+\Gamma(f)(f')\cdot f'$, and inductively we get
\[
\widetilde\nabla^sf=\nabla^sf+P_s\left(f\,;\,\nabla^1f,\,\dotsc,\,\nabla^{s-1}f\right)
\]
where $P(z\,;\,\xi_1,\,\dotsc,\,\xi_{s-1})$ is a polynomial with $C^\infty$ coefficients in $z\in U$, which is of weighted homogeneous degree $s$ in
$(\xi_1,\,\dotsc,\,\xi_{s-1})$. In other words, the corresponding isomorphisms
$J^kV\simeq V^{\oplus k}$ correspond to each other by a
$\bC^*$-homogeneous transformation
$(\xi_1,\,\dotsc,\,\xi_k)\mapsto (\widetilde\xi_1,\,\dotsc,\,\widetilde\xi_k)$ such that
\[
\widetilde\xi_s=\xi_s+P_s\left(z\,;\,\xi_1,\dotsc,\xi_{s-1}\right).
\]
Let us introduce the corresponding rescaled components
\[
(\xi_{1,\varepsilon},\dotsc,\xi_{k,\varepsilon})= \left(\varepsilon_1^1\xi_1,\,\,\dotsc,\,\varepsilon_k^k\xi_k\right),\qquad (\widetilde\xi_{1,\varepsilon},\,\dotsc,\,\widetilde\xi_{k,\varepsilon})= \left(\varepsilon_1^1\widetilde\xi_1,\,\dotsc,\,\varepsilon_k^k\widetilde\xi_k\right).
\]
Then
\begin{align*}
\widetilde\xi_{s,\varepsilon} &=
\xi_{s,\varepsilon}+
\varepsilon_s^s\,P_s\left(x\,;\,\varepsilon_1^{-1}\xi_{1,\varepsilon},\,\dotsc,\,
\varepsilon_{s-1}^{-(s-1)}\xi_{s-1,\varepsilon}\right)
\\
&=
\xi_{s,\varepsilon}+O(\varepsilon_s/\varepsilon_{s-1})^s\, O\left(\Abs{\xi_{1,\varepsilon}} +\dotsb+
\Abs{\xi_{s-1,\varepsilon}}^{1/(s-1)}\right)^s
\end{align*}
and it is easily seen, as a simple consequence of the mean value inequality
$\abs{\Abs{x}^\gamma-\Abs{y}^\gamma}\leq\gamma\sup_{z\,\in\,[x,y]}
\Abs{z}^{\gamma-1}\Abs{x-y}$, that the ``error term'' in the difference
$\Abs{\widetilde\xi_{s,\varepsilon}}^{2b/s}-\Abs{\xi_{s,\varepsilon}}^{2b/s}$ is bounded by
\[
(\varepsilon_s/\varepsilon_{s-1})^s\,
\big(\Abs{\xi_{1,\varepsilon}}+\cdots+
\Abs{\xi_{s-1,\varepsilon}}^{1/(s-1)}+
\Abs{\xi_{s,\varepsilon}}^{1/s}\big)^{2b}.
\]
When $b/s$ is an integer, similar bounds hold for all derivatives $D_{z,\xi}^\beta(\Abs{\widetilde\xi_{s,\varepsilon}}^{2b/s}-
\Abs{\xi_{s,\varepsilon}}^{2b/s})$ and the Lemma~\ref{lemm:rescaling} follows.
\end{proof}

Now, we fix a point $z_0\in X$, a local holomorphic frame
$(e_\lambda(z))_{1\,\leq\,\lambda\,\leq\,r}$ satisfying~\eqref{eq:frame} on a neighborhood $U$ of $z_0$, and the \emph{holomorphic} connection $\nabla$ on $V_{|U}$ such that
$\nabla e_\lambda=0$. Since the uniform estimates of Lemma~\ref{lemm:rescaling} also apply locally (provided they are applied on a relatively compact open subset $U'\Subset U$), we can use the corresponding holomorphic trivialization
$J^kV_{|U}\simeq V_{|U}^{\oplus k}\simeq U\times(\bC^r)^{\oplus k}$ to make our calculations. We do this in terms of the rescaled components $\xi_s=\varepsilon_s^s\nabla^sf(0)$. Then, uniformly on compact subsets of $J^kV_{|U}\setminus\{0\}$, we have
\[
\Psi_{h,b,\varepsilon}\circ\rho_{\nabla,\varepsilon}^{-1}(z\,;\,\xi_1,\,\dotsc,\,\xi_k) =\left(\sum_{1\,\leq\,s\,\leq\,k}\Abs{\xi_s}^{2b/s}_{h(z)}\right)^{1/b} +O\left(\max(\varepsilon_s/\varepsilon_{s-1})^{1/b}\right),
\]
and the error term remains of the same magnitude when we take any derivative $D_{z,\xi}^\beta$. By~\eqref{eq:frame} we find
\[
\Abs{\xi_s}_{h(z)}^2=
\sum_\lambda\abs{\xi_{s,\lambda}}^2+
\sum_{i,j,\lambda,\mu}c_{ij\lambda\mu}\,z_i\bar z_j
\,\xi_{s,\lambda}\bar \xi_{s,\mu}+O\left(\abs{z}^3\abs{\xi}^2\right).
\]
The question is thus reduced to evaluating the curvature of the weighted Finsler metric on $V^{\oplus k}$ defined by
\begin{align*}
\Psi(z\,;\,\xi_1,\,\dotsc,\,\xi_k) &=\left(\sum_{1\,\leq\,s\,\leq\,k}\Abs{\xi_s}^{2b/s}_{h(z)}\right)^{1/b}
\\
&=\left(\sum_{1\,\leq\,s\,\leq\,k}\left(\sum_\lambda\abs{\xi_{s,\lambda}}^2+
\sum_{i,j,\lambda,\mu}c_{ij\lambda\mu}\,z_i\bar z_j\,
\xi_{s,\lambda}\bar\xi_{s,\mu}\right)^{b/s}\right)^{1/b}+O(\abs{z}^3).
\end{align*}
We set $\abs{\xi_s}^2=\sum_\lambda\abs{\xi_{s,\lambda}}^2$. A straightforward calculation yields the Taylor expansion
\begin{multline*}
\log\Psi(z\,;\,\xi_1,\,\dotsc,\,\xi_k)\\
=\tfrac{1}{ b}\log\sum_{1\,\leq\,s\,\leq\,k}\abs{\xi_s}^{2b/s}+
\sum_{1\,\leq\,s\,\leq\,k}\tfrac{1}{ s}\,\tfrac{\abs{\xi_s}^{2b/s}}{ \sum_t\abs{\xi_t}^{2b/t}}AAa
\sum_{i,j,\lambda,\mu}c_{ij\lambda\mu}z_i\bar z_j
\tfrac{\xi_{s,\lambda}\bar\xi_{s,\mu}}{\abs{\xi_s}^2}+O(\abs{z}^3).
\end{multline*}
By~\eqref{eq:curvature}, the curvature form of $L_k=\cO_{X_k(V)}(1)$ is given at the central point $z_0$ by the formula
\[
\Theta_{L_k,\Psi^*_{h,b,\varepsilon}}(z_0,[\xi])\simeq
\omega_{r,k,b}(\xi)+\tfrac{\imath}{ 2\pi}
\sum_{1\,\leq\,s\,\leq\,k}\tfrac{1}{ s}\,\tfrac{\abs{\xi_s}^{2b/s}}{ \sum_t\abs{\xi_t}^{2b/t}}
\sum_{i,j,\lambda,\mu}c_{ij\lambda\mu}
\tfrac{\xi_{s,\lambda}\bar\xi_{s,\mu}}{\abs{\xi_s}^2}\,dz_i\wedge d\bar z_j
\]
where
\[
[\xi] = [\xi_1,\,\dotsc,\,\xi_k]\in\bP\left(1^{[r]},2^{[r]},\,\dotsc,\,k^{[r]}\right)
\]
and
\[
\omega_{r,k,b}(\xi) =
\tfrac{\imath}{ 2\pi}\ddbar\left(\tfrac{1}{ b}\log\sum_{1\,\leq \,s\,\leq\,k} \abs{\xi_s}^{2b/s}\right).
\]
The fibers $\bP(1^{[r]},2^{[r]},\,\dotsc,\,k^{[r]})$ of
$X_k(V)\to X$ can be represented as a quotient of the ``weighted ellipsoid'' $\sum_{s=1}^k\abs{\xi_s}^{2b/s}=1$ by the $\bS^1$-action induced by the weighted $\bC^*$-action. This suggests to make use of polar coordinates and to set
\begin{align}\label{eq:coordinates}
x_s&=\abs{\xi_s}^{2b/s},\quad x=(x_1,\,\dotsc,\,x_k)\in\bR^k,
\\
\tag{\theequation\ensuremath{'}} u_s&=\tfrac{\xi_s}{ \abs{\xi_s}}\in \bS^{2r-1}\subset\bC^r,\quad u=(u_1,\,\dotsc,\,u_k)\in(\bS^{2r-1})^k,
\end{align}
so that
\[
\tag{\theequation\ensuremath{'}\ensuremath{'}}
\sum_{s=1}^kx_s=1\quad\text{and}\quad \xi_s=x_s^{s/2b}u_s.\kern102pt
\]
The Morse integrals will then have to be computed for
$(x,u)\in\bDelta^{k-1}\times(\bS^{2r-1})^k$, where
$\bDelta^{k-1}\subset\bR^k$ is the $(k-1)$-dimensional simplex.

\begin{prop}
With respect to the rescaled components $\xi_s=\varepsilon_s^s\nabla^sf(0)$ at $z=f(0)\in X$ and the above choice of coordinates (\eqref{eq:coordinates}, \eqref{eq:coordinates}\ensuremath{'}, \eqref{eq:coordinates}\ensuremath{'}\ensuremath{'}), the curvature of the tautological sheaf $L_k=\cO_{X_k(V)}(1)$ admits an approximate expression
\begin{flalign*}
\quad
\tag{a}
\Theta_{L_k,\Psi^*_{h,b,\varepsilon}}(z,[\xi])=
\omega_{r,k,b}(\xi)+g_{V,k}(z,x,u)+\text{\rm(error terms)}, &&
\end{flalign*}
where $(x,u)\in\bDelta^{k-1}\times (\bS^{2r-1})^k$, $\xi_s=x_s^{s/2b}u_s\in\bC^r$,
\begin{flalign*}
\quad
\tag{b}
\omega_{r,k,b}(\xi)=\tfrac{\imath}{ 2\pi}\ddbar\bigg(
\tfrac{1}{ b}\log \sum_{1\,\leq\,s\,\leq\,k}\abs{\xi_s}^{2b/s}\bigg) &&
\end{flalign*}
is a Fubini--Study type K\"ahler metric on
$\bP(1^{[r]},2^{[r]},\,\dotsc,\,k^{[r]})$, associated with the canonical
$\bC^*$ action on $J^kV$ of weight $a=(1^{[r]},2^{[r]},\,\dotsc,\,k^{[r]})$, and
\begin{flalign*}
\quad
\tag{c} g_{V,k}(z,x,u)=\tfrac{\imath}{ 2\pi}\sum_{1\,\leq\,s\,\leq\,k}\tfrac{x_s}{ s}
\sum_{i,j,\lambda,\mu}c_{ij\lambda\mu}(z)\, u_{s,\lambda}\bar u_{s,\mu}\,dz_i\wedge d\bar z_j. &&
\end{flalign*}
Here $(\tfrac{\imath}{ 2\pi}\,c_{ij\lambda\mu})$ are the coefficients of $-\Theta_{V,h}$, and the error terms admit an upper bound
\begin{flalign*}
\quad
\tag{d}
\text{\rm(error terms)}\le O\left(\max_{2\,\leq\,s\,\leq\,k}(\varepsilon_s/\varepsilon_{s-1})^s\right)\quad &&
\end{flalign*} uniformly on the compact variety $X_k(V)$.
\end{prop}
\begin{proof}
The error terms on $\Theta_{L_k}$ come from the differentiation of the error terms on the Finsler metric, found in Lemma~\ref{lemm:rescaling}. They can indeed be differentiated if $b$ is a multiple of $\lcm(1,2,\,\dotsc,\,k)$, since $2b/s$ is then an even integer.
\end{proof}

For the calculation of Morse integrals, it is useful to find the expression of the volume form $\omega_{r,k,b}^{kr-1}$ on
$\bP(1^{[r]},2^{[r]},\,\dotsc,\,k^{[r]}) = (\bDelta^{k-1}\times (\bS^{2r-1})^k)/\bS^1$ in terms of the coordinates $(x,u)$. We refer to~\cite[Prop.~1.13]{Dem11} for the proof.

\begin{prop}
\leavevmode\label{prop:volume_form}
\begin{enumerate}\alphenumi
\item\label{prop4.3.a} The volume form $\omega_{r,k,b}^{kr-1}$ is the quotient of the measure $\tfrac{1}{ k!^r}\nu_{k,r}\otimes\mu$ on $\bDelta^{k-1}\times(\bS^{2r-1})^k$, where
\[
d\nu_{k,r}(x)=(kr-1)!\tfrac{(x_1\ldots\,x_k) ^{r-1}}{ (r-1)!\,{}^k} dx_1\wedge\ldots\wedge dx_{k-1},\quad d\mu(u)=d\mu_1(u_1)\ldots d\mu_k(u_k)
\]
are probability measures on $\bDelta^{k-1}$ and $(\bS^{2r-1})^k$ respectively
$(\mu$ being the rotation inva\-riant one).
\item\label{prop4.3.b} We have the equality
$\int_{\bP(1^{[r]},2^{[r]},\,\dotsc,\,k^{[r]})}\omega_{r,k,b}^{kr-1}=\tfrac{1}{ k!^r}$~(independent of $b$).
\end{enumerate}
\end{prop}

\subsection{Logarithmic and orbifold jet metrics}\label{sse:orbi_jet_metrics}
Consider now an orbifold directed structure
$(X,V,D)$, where $V\subset T_X$ is a subbundle, $r=\rank(V)$, and
$D=\sum(1-\tfrac{1}{ \rho_j})\Delta_j$ is a normal crossing divisor that is assumed to intersect $V$ transversally everywhere. One then performs very similar calculations to what we did in \S~\ref{sse:jet_metrics}, but with adapted Finsler metrics. Fix a point $z_0$ at which $p$ components
$\Delta_j$ meet, and use coordinates $(z_1,\,\dotsc,\,z_n)$ such that
$V_{z_0}$ is spanned by $(\tfrac{\partial}{\partial z_1},\,\dotsc,\,\tfrac{\partial}{\partial z_r})$ and $\Delta_j$ is defined by
$z_j=0$, $1\leq j\leq p\leq r$. In the logarithmic case
$\rho_j=\infty$, the logarithmic dual bundle $\cO(V^*\langle D\rangle)$ is spanned by
\[
\tfrac{dz_1}{ z_1},\,\dotsc,\,\tfrac{dz_p}{ z_p},~dz_{p+1},\dotsc,dz_n.
\]
The logarithmic jet differentials are just polynomials in
\[
\tfrac{d^sz_1}{ z_1},\,\dotsc,\,\tfrac{d^sz_p}{ z_p},~d^sz_{p+1},\dotsc,d^sz_n,
\quad 1\leq s\leq k,
\]
and the corresponding $(\varepsilon_1,\,\dotsc,\,\varepsilon_k)$-rescaled Finsler metric is
\begin{equation}\label{eq:Finsler}
\left(
\sum_{s=1}^k\varepsilon_s^{2b}\left(\sum_{j=1}^p
\abs{f_j}^{-2}\left|f_j^{(s)}\right|^2+\sum_{j=p+1}^r\left|f_j^{(s)}\right|^2\right)^{2b/s}
\right)^{1/b}.
\end{equation}
Alternatively, we could replace $\abs{f_j}^{-2}\abs{f_j^{(s)}}^2$ by
$\abs{(\log f_j)^{(s)}}^2$ which has the same leading term and differs by a weighted degree $s$ polynomial in the $f_j^{-1}f_j^{(\ell)}$,
$\ell<s\,$; an argument very similar to the one used in the proof of Lemma~\ref{lemm:rescaling} then shows that the difference is negligible when $\varepsilon_1\gg \varepsilon_2\gg \cdots\gg\varepsilon_k$. However~\eqref{eq:Finsler} is just the case of the model metric, in fact we get
$r$-tuples $\xi_s=(\xi_{s,j})_{1\,\leq\,j\,\leq\,r}$ of components produced by the trivialization of the logarithmic bundle
$\cO(V\langle D\rangle)$, such that
\[
\xi_{s,j}=f_j^{-1}f_j^{(s)}\quad\text{for $1\leq s\leq p$ and}\quad
\xi_{s,j}=f_j^{(s)}\quad\text{for $p+1\leq s\leq r$}.
\]
In general, we are led to consider Finsler metrics of the form
\[
\left(\sum_{s=1}^k\varepsilon_s^{2b}\Abs{\xi_s}_{h(z)}^{2b/s}\right)^{1/b},
\quad\xi_s = (\xi_{s,j})_{1\,\leq\,j\,\leq\,r},
\]
where $h(z)$ is a variable hermitian metric on the logarithmic bundle
$V\langle D\rangle$. In the orbifold case, the appropriate ``model'' Finsler metric is
\[
\left(
\sum_{s=1}^k\varepsilon_s^{2b}
\left(
\sum_{j=1}^p
\abs{f_j}^{-2(1-s/\rho_j)_+}\left|f_j^{(s)}\right|^2 +
\sum_{j=p+1}^r
\left|f_j^{(s)}\right|^2
\right)^{2b/s}
\right)^{1/b}.
\]
As a consequence of Remark~\ref{rema:integrability}, we would get a metric with equivalent singularities on the dual $L_k^*$ of the tautological sheaf
$L_k=\cO_{X_k(V\langle D\rangle)}(1)$ by replacing
$\sum_{j=p+1}^r\abs{f_j^{(s)}}^2$ with $\sum_{j=1}^r\abs{f_j^{(s)}}^2$ (or by any smooth hermitian norm $h$ on $V$), since the extra terms
$\sum_{j=1}^p\abs{f_j^{(s)}}^2$ are anyway controlled by the ``orbifold part'' of the summation. Of course, we need to find a suitable~Finsler metric that is globally defined on $X$. This can be done by taking smooth metrics $h_{V,s}$ on $V$ and $h_j$ on $\cO_X(\Delta_j)$ respectively, as well as smooth connections $\nabla$ and $\nabla_j$. One can then consider the globally defined metric
\begin{equation}\label{eq:orbi_Finsler}
\left(\sum_{s=1}^k\varepsilon_s^{2b}\left(\left|\nabla^{(s)}f\right|_{h_{V,s}}^2 +\sum_j\Abs{\sigma_j(f)}_{h_j}^{-2(1-s/\rho_j)_+}
\left|\nabla_j^{(s)}(\sigma_j\circ f)\right|_{h_j}^2\right)^{2b/s}\right)^{1/b}
\end{equation}
where $D=\sum(1-\tfrac{1}{\rho_j})\Delta_j$ and $\sigma_j\in H^0(X,\cO_X(\Delta_j))$ are the tautological sections; here, we want the flexibility of not necessarily taking the same hermitian metrics on $V$ to evaluate the various norms $\Abs{\nabla^{(s)}f}_{h_{V,s}}$. We obtain Finsler metrics with equivalent singularities by just changing the $h_{V,s}$ and $h_j$ (and keeping $\nabla$, $\nabla_j$ unchanged). If we also change the connections, then an argument very similar to the one used in the proof of Lemma~\ref{lemm:rescaling} shows that the ratio of the corresponding metrics is
$1\pm O(\max(\varepsilon_s/\varepsilon_{s-1}))$, and therefore arbitrary close to $1$ whenever $\varepsilon_1\gg\varepsilon_2\gg \cdots\gg\varepsilon_k$; in any case, we get metrics with equivalent singularities. Fix $z_0\in X$ and use coordinates $(z_1,\dotsc,z_n)$ as described at the beginning of \S~\ref{sse:orbi_jet_metrics}, so that $\sigma_j(z)=z_j$, $1\leq j\leq p$, in a suitable trivialization of $\cO_X(\Delta_j)$. Let $f$ be a $k$-jet of curve such that $f(0)=z\in X\setminus\abs{D}$ is in a sufficiently small neighborhood of $z_0$. By employing the trivial connections associated with the above coordinates, the derivative $f^{(s)}$ is described by components
\[
\xi_{s,j}=f_j^{(s)},~~1\leq j\leq r,\quad
\xi^{\log}_{s,j}=f_j^{-1}f_j^{(s)},\quad
\xi^\orb_{s,j}=f_j^{-(1-s/\rho_j)_+}f_j^{(s)},\quad 1\leq j\leq p,
\]
and $\xi^\orb_{s,j}=\xi^{\log}_{s,j}=\xi_{s,j}$ for $p+1\leq j\leq r$. Here $\xi^\orb_{s,j}$ are to be thought of as the components of $f^{(s)}$ in the ``virtual'' vector bundle $V\langle D^{(s)}\rangle$, and the fact that the argument of these complex numbers is not uniquely defined is irrelevant, because the only thing we need to compute the norms is $\abs{\xi^\orb_{s,j}}$. Accordingly, for
$v\in V_z$, $v\simeq(v_j)_{1\,\leq\,j\,\leq\,r}\in\bC^r$, we put
\[
v^{\log}_j=z_j^{-1}v_j=\sigma_j(z)^{-1}\nabla_j\sigma_j(v)\quad\text{and}\quad v^\orb_j=z_j^{-(1-s/\rho_j)_+}v_j,~~1\leq j\leq p,
\]
and define the orbifold hermitian norm on $V\langle D^{(s)}\rangle$ associated with $h_{V,s}$ and $h_j$ by

\begin{align}\label{eq:orbi_norm}
\Abs{v^\orb}_{\widetilde h_s}{\kern-3pt}^2&=\Abs{v}_{h_{V,s}}^2 +\sum_{j=1}^p\Abs{\sigma_j(z)}_{h_j}^{-2\left(1-s/\rho_j)_+\right)}
\Abs{\nabla_j\sigma_j(v)}_{h_j}^2
\\\label{eq:orbi_norm'}
\tag{\theequation\ensuremath{'}} &=\Abs{v}_{h_{V,s}}^2+\sum_{j=1}^p
\Abs{\sigma_j(z)}_{h_j}^{2\left(1-(1-s/\rho_j)_+\right)}\abs{v_j^{\log}}^2
\\\label{eq:orbi_norm''}
\tag{\theequation\ensuremath{'}\ensuremath{'}} &=\Abs{v}_{h_{V,s}}^2+
\sum_{j=1}^p\Abs{v_j^\orb}_{h_j^{1-\left(1-s/\rho_j\right)_+}}^2.
\end{align}
With this notation, the orbifold Finsler metric~\eqref{eq:orbi_Finsler} on $k$-jets is reduced to an expression
\begin{equation}\label{eq:orbi_Finsler'}
\begin{split}
\left\|\xi^\orb\right\|_{\Psi_{h,b,\varepsilon}}^{\,2\phantom{\big|}}=
\left(
\sum_{s=1}^k
\varepsilon_s^{2b}
\left\|\xi_s^\orb\right\|_{\widetilde h_s}^{2b/s}
\right)^{1/b},
\quad\xi_s^\orb &=\left(\xi_{s,j}^\orb\right)_{1\,\leq\,j\,\leq\,r}\,,\\
\xi^\orb &=\left(\xi_s^\orb\right)_{1\,\leq\,s\,\leq\,k}\,,
\end{split}
\end{equation}
formally identical to what we had in the compact or logarithmic cases. If $v$ is a local holomorphic section of $\cO_X(V)$, formula~\eqref{eq:orbi_norm} shows that the norm
$\Abs{v^\orb}_{\widetilde h_s}$ can take infinite values when $z\in\abs{D}$, while, by~\eqref{eq:orbi_norm'}, the norm is always bounded (but slightly degenerate along
$\abs{D}$) if $v$ is a section of the logarithmic sheaf
$\cO_X(V\langle\lceil D\rceil\rangle)$; we think intuitively of the orbifold total space $V\langle D^{(s)}\rangle$ as the subspace of $V$ in which the tubular neighborhoods of the zero section are defined by $\Abs{v^\orb}_{\widetilde h_s}<\varepsilon$ for $\varepsilon>0$.

\begin{rema}
When $\rho_j\in\bQ$, we can take an adapted Galois cover $g\colon Y\to X$ such that $(z_j\circ g)^{1-(1-s/\rho_j)_+}$ is univalent on $Y$ for all components $\Delta_j$ involved, and we then get a well defined locally free sheaf $\cO_Y\big(g^*V\langle D^{(s)}\big)$ such that
\[
g^*\big(\cO_X(V\langle\lceil D\rceil\rangle)\big)\subset
\cO_Y\big(g^*V\left\langle D^{(s)}\right\rangle\big)\subset g^*\big(\cO_X(V)\big).
\]
However, as already stressed in Remark~\ref{rema:q_coefficients}, this viewpoint is not needed in our analytic approach.
\end{rema}

\subsection{Orbifold tautological sheaves and their curvature} In this context, we define the orbifold tautological sheaves
\[
\cO_{X_k(V\langle D\rangle)}(m)
\bydef
\cO_{X_k(V\langle\lceil D\rceil\rangle)}(m)\otimes
\cI((\Psi_{k,b,\varepsilon}^*)^m)
\]
to be the logarithmic tautological sheaves $\cO_{X_k(V\langle\lceil D\rceil\rangle)}(m)$ twisted by the multiplier ideal sheaves associated with the dual metric $\Psi_{k,b,\varepsilon}^*$ (cf.~\eqref{eq:orbi_Finsler'}), when these are viewed as singular hermitian metrics over the logarithmic
$k$-jet bundle $X_k(V\langle\lceil D\rceil\rangle)$. In accordance with this viewpoint, we simply define the orbifold $k$-jet bundle to be
$X_k(V\langle D\rangle)=X_k(V\langle\lceil D\rceil\rangle)$. The calculation of the curvature tensor is formally the same as in the case
$D=0$, and we obtain:

\begin{prop}\label{prop:orbi_curvature}
With respect to the (rescaled) orbifold $k$-jet components
\[
\xi_{s,\lambda}=\varepsilon_s^sf_\lambda^{(1-(1-\rho_\lambda/s)_+)} f_\lambda^{(s)}(0),~~1\leq\lambda\leq p,\quad\text{and}\quad
\xi_{s,\lambda}=\varepsilon_s^sf_\lambda^{(s)}(0),~~p+1\leq\lambda\leq r,
\]
and of the dual metric $\Psi^*_{h,b,\varepsilon}$, the curvature form of the tautological sheaf $L_k=\cO_{X_k(V\langle D\rangle)}(1)$ admits at any point $(z,[\xi])\in X_k(V\langle D\rangle)$ an approximate expression
\begin{flalign}\label{prop4.5.a}
\tag{a}
\Theta_{L_k,\Psi^*_{h,b,\varepsilon}}(z,[\xi])
\simeq
\omega_{r,k,b}(\xi)+g_{V,D,k}(z,x,u), &&
\end{flalign}
where $x_s=\abs{\xi_s}^{2b/s}$, $u_s=\tfrac{\xi_s}{\abs{\xi_s}}\in \bS^{2r-1}$ are polar coordinates associated with $\xi_s=(\xi_{s,\lambda})_{1\,\leq\,\lambda\,\leq\,k}$ in $\bC^r$, $x=(x_1,\,\dotsc,\,x_k)\in\bDelta^{k-1}$,
$[\xi]=[\xi_1,\,\dotsc,\,\xi_k]\in\bP(1^{[r]},2^{[r]},\dotsc,k^{[r]})$ and
\begin{flalign}\label{prop4.5.b}
\tag{b} g_{V,D,k}(z,x,u)=\tfrac{\imath}{ 2\pi}
\sum_{1\,\leq\,s\,\leq\,k}\tfrac{x_s}{ s}\sum_{i,j,\lambda,\mu}c^{(s)}_{ij\lambda\mu}(z)\, u_{s,\lambda}\bar u_{s,\mu}\,dz_i\wedge d\bar z_j. &&
\end{flalign}
Here $(\tfrac{\imath}{ 2\pi}\,c^{(s)}_{ij\lambda\mu})$ are the coefficients of the curvature tensor $-\Theta_{V\langle D^{(s)}\rangle,\widetilde h_s}$, and the error terms are
$O(\max_{2\,\leq\,s\,\leq\,k}(\varepsilon_s/\varepsilon_{s-1})^s)$, uniformly on the projectivized orbifold variety $X_k(V\langle D\rangle)$.
\end{prop}

Notice, as is clear from the expressions~\eqref{eq:orbi_norm''}, \eqref{eq:orbi_Finsler'} and the fact that
$v_j=z_jv^\orb_j$, that our orbifold Finsler metrics always have fiberwise positive curvature, equal to $\omega_{k,r,b}(\xi)$, along the fibers of $X_k(V\langle D\rangle)\to X$ (even after taking into account the so-called error terms, because fiberwise, the functions under consideration are just sums of even powers $\abs{\widetilde\xi_s^\orb}^{2b/s}$ in suitable $k$-jet components, and are therefore plurisubharmonic.)

\section{Existence theorems for jet differentials}\label{se:Morse_integrals}
\subsection{Expression of the Morse integral}

Thanks to the uniform approximation provided by Proposition~\ref{prop:orbi_curvature}, we can (and will) neglect the $O(\varepsilon_s/\varepsilon_{s-1})$ error terms in our calculations. Since $\omega_{r,k,b}$ is positive definite on the fibers of $X_k(V\langle D\rangle)\to X$ (at least outside of the axes $\xi_s=0$), the index of the $(1,1)$ curvature form
$\Theta_{L_k,\Psi^*_{h,b,\varepsilon}}(z,[\xi])$ is equal to the index of the $(1,1)$-form $g_{V,D,k}(z,x,u)$. By the binomial formula, the $q$-index integral of $(L_k,\Psi_{h,b,\varepsilon}^*)$ on
$X_k(V\langle D\rangle)$ is therefore equal to
\begin{multline}
\int_{X_k(V\langle D\rangle)(L_k,q)}
\Theta_{L_k,\Psi^*_{h,b,\varepsilon}}^{n+kr-1}
\\
=\tfrac{(n+kr-1)!}{ n!(kr-1)!}
\int_{z\,\in\,X}\int_{\xi\,\in\,\bP\left(1^{[r]},\,\dotsc,\,k^{[r]}\right)}\omega_{r,k,b}^{kr-1}(\xi)
\wedge\bOne_{g_{V,D,k},q}(z,x,u)\,g_{V,D,k}(z,x,u)^n
\end{multline}
where $\bOne_{g_{V,D,k},q}(z,x,u)$ is the characteristic function of the open set of points where $g_{V,D,k}(z,x,u)$ has signature $(n-q,q)$ in terms of the $dz_j$'s.

Notice that since $g_{V,D,k}(z,x,u)^n$ is a determinant, the product
\[
\bOne_{g_{V,D,k},q}(z,x,u)\,g_{V,D,k}(z,x,u)^n
\]
gives rise to a continuous function on $X_k(V\langle D\rangle)$. By Proposition~\ref{prop:volume_form}$\MK$\eqref{prop4.3.b}, we get

\begin{multline}\label{eq:volume}
\int_{X_k(V\langle D\rangle)(L_k,q)}
\Theta_{L_k,\Psi^*_{h,b,\varepsilon}}^{n+kr-1} =\tfrac{(n+kr-1)!}{n!\,k!^r(kr-1)!}~~\times
\\
\int_{z\in X}
\int_{(x,u)\,\in\,\bDelta^{k-1}\times\left(\bS^{2r-1}\right)^k}\bOne_{g_{V,D,k},q}(z,x,u)\, g_{V,D,k}(z,x,u)^n\,d\nu_{k,r}(x)\,d\mu(u).
\end{multline}

\subsection{Probabilistic estimate of cohomology groups}\label{sse:proba_estimate}
We assume here that we are either in the ``compact'' case $(D=0)$, or in the logarithmic case $(\rho_j=\infty)$. Then the curvature coefficients $\smash{c^{(s)}_{ij\lambda\mu}}=c_{ij\lambda\mu}$ do not depend on $s$ and are those of the dual bundle $V^*$ (resp.\ $V^*\langle D\rangle$). In this situation, Proposition~\ref{prop:orbi_curvature}$\MK$\eqref{prop4.5.b} for $g_{V,D,k}(z,x,u)$ can be thought of as a ``Monte Carlo'' evaluation of the curvature tensor, obtained by averaging the curvature at random points $u_s\in \bS^{2r-1}$ with certain positive weights $x_s/s\,$; we then think of the \text{$k$-jet} $f$ as some sort of random variable such that the derivatives $\nabla^kf(0)$ (resp.\ logarithmic derivatives) are uniformly distributed in all directions. Let us compute the expected value of $(x,u)\mapsto g_{V,D,k}(z,x,u)$ with respect to the probability measure $d\nu_{k,r}(x)\,d\mu(u)$. Since
\[
\int_{\bS^{2r-1}}u_{s,\lambda}\bar u_{s,\mu}d\mu(u_s)=\tfrac{1}{ r}
\delta_{\lambda\mu}\quad
\text{and}\quad\int_{\bDelta^{k-1}}x_s\,d\nu_{k,r}(x)=\tfrac{1}{ k},
\]
we find the expected value
\[
\bE(g_{V,D,k}(z,\bullet,\bullet))=\tfrac{1}{ kr}
\sum_{1\,\leq\,s\,\leq\,k}\tfrac{1}{ s}\cdot\tfrac{\imath}{ 2\pi}\sum_{i,j,\lambda} c_{ij\lambda\lambda}(z)\,dz_i\wedge d\bar z_j.
\]
In other words, we get the normalized trace of the curvature, i.e.
\begin{equation}\label{4.3}
\bE(g_{V,D,k}(z,\bullet,\bullet))=\tfrac{1}{ kr}\Big(1+\tfrac{1}{ 2}+\cdots+\tfrac{1}{ k}\Big)
\Theta_{\det(V^*\langle D\rangle),\det h^*},
\end{equation}
where $\Theta_{\det(V^*\langle D\rangle),\det h^*}$ is the
$(1,1)$-curvature form of $\det(V^*\langle D\rangle)$ with the metric induced by $h$. It is natural to guess that
$g_{V,D,k}(z,x,u)$ behaves asymptotically as its expected value
$\bE(g_{V,D,k}(z,\bullet,\bullet))$ when $k$ tends to infinity. If we replace brutally
$g_{V,D,k}$ by its expected value in~\eqref{eq:volume}, we get the integral
\[
\tfrac{(n+kr-1)!}{ n!\,k!^r(kr-1)!}\tfrac{1}{ (kr)^n}
\left(1+\tfrac{1}{ 2}+\cdots+\tfrac{1}{ k}\right)^n\int_X\bOne_{\eta,q}\eta^n,
\]
where $\eta\bydef\Theta_{\det(V^*\langle D\rangle),\det h^*}$ and
$\bOne_{\eta,q}$ is the characteristic function of its $q$-index set in $X$. The leading constant is equivalent to $(\log k)^n/n!\,k!^r$ modulo a multiplicative factor \text{$1+O(1/\log k)$}. By working out a more precise analysis of the deviation, the following result has been proved in [Dem11] in the compact case; the more general logarithmic case can be treated without any change, so we state the result in this situation by just transposing the results of~\cite{Dem11}.

\begin{enonce}
{Probabilistic estimate}\label{pe}
Let $(X,V,D)$ be a non-singular logarithmic directed variety. Fix smooth hermitian metrics $\omega$ on $T_X$, $h$ on $V\langle D\rangle$, and write
$\omega=\tfrac{\imath}{ 2\pi} \sum\omega_{ij}dz_i\wedge d\bar z_j$ on $X$. Denote by $\Theta_{V\langle D\rangle,h}=-\tfrac{\imath}{ 2\pi}\sum c_{ij\lambda\mu}dz_i\wedge d\bar z_j\otimes e_\lambda^*\otimes e_\mu$ the curvature tensor of $V\langle D\rangle$ with respect to an $h$-orthonormal frame $(e_\lambda)$, and put
\[
\eta(z)
\bydef
\Theta_{\det(V^*\langle D\rangle),\det h^*}=
\tfrac{\imath}{ 2\pi}\sum_{1\,\leq\,i,j\,\leq\,n}\eta_{ij} dz_i\wedge d\bar z_j,\qquad
\eta_{ij}
\bydef
\sum_{1\,\leq\,\lambda\,\leq\,r}c_{ij\lambda\lambda}.
\]
Finally consider the $k$-jet line bundle
$L_k=\smash{\cO_{X_k(V\langle D\rangle)}(1)}\to X_k(V\langle D\rangle)$ equipped with the induced metric
$\Psi^*_{h,b,\varepsilon}$ (as defined above, with
$1=\varepsilon_1\gg\varepsilon_2\gg\ldots\gg \varepsilon_k>0)$. When $k$ tends to infinity, the integral of the top power of the curvature of
$L_k$ on its $q$-index set $X_k(V\langle D\rangle)(L_k,q)$ is given by
\[
\int_{X_k(V\langle D\rangle)(L_k,q)}
\Theta_{L_k,\Psi^*_{h,b,\varepsilon}}^{n+kr-1}=
\tfrac{(\log k)^n}{ n!\,k!^r}\bigg(
\int_X\bOne_{\eta,q}\eta^n+O\left((\log k)^{-1}\right)\bigg)
\]
for all $q=0,1,\,\dotsc,\,n$, and the error term $O((\log k)^{-1})$ can be bounded explicitly in terms of $\Theta_{V\langle D\rangle}$, $\eta$ and
$\omega$. Moreover, the left hand side is identically zero for $q>n$.
\end{enonce}

The final statement follows from the observation that the curvature of $L_k$ is positive along the fibers of $X_k(V\langle D\rangle)\to X$, by the plurisubharmonicity of the weight (this is true even when the error terms are taken into account, since they depend only on the base); therefore the
$q$-index sets are empty for $q>n$. It will be useful to extend the above estimates to the case of sections of
\[
L_{F,k} =
\cO_{X_k(V\langle D\rangle)}(1)
\otimes
\pi_k^*\cO\Big(-\tfrac{1}{ kr}\Big(1+\tfrac{1}{ 2}+\cdots+\tfrac{1}{ k}\Big)F\Big)
\]
where $F\in\Pic_\bQ(X)$ is an arbitrary $\bQ$-line bundle on $X$ and
$\pi_k\colon X_k(V\langle D\rangle)\to X$ is the natural projection. We assume here that $F$ is also equipped with a smooth hermitian metric $h_F$. In formulas~\eqref{eq:volume}--\eqref{pe}, the curvature $\Theta_{L_{F,k}}$ of $L_{F,k}$ takes the form $\Theta_{L_{F,k}}=\omega_{r,k,b}(\xi)+g_{V,D,F,k}(z,x,u)$ where
\[
g_{V,D,F,k}(z,x,u)=g_{V,D,k}(z,x,u)-
\tfrac{1}{ kr}\Big(1+\tfrac{1}{ 2}+\cdots+\tfrac{1}{ k}\Big)\Theta_{F,h_F}(z),
\]
and by the same calculations its normalized expected value is
\[
\eta_F(z)
\bydef
\tfrac{1}{\tfrac{1}{ kr}(1+\tfrac{1}{ 2}+\cdots+\tfrac{1}{ k})}
\bE(g_{V,D,F,k}(z,\bullet,\bullet))=
\Theta_{\det V^*\langle D\rangle,\det h^*}(z)-\Theta_{F,h_F}(z).
\]
Then the variance estimate for $g_{V,D,F,k}$ is the same as the variance estimate for $g_{V,D,k}$, and the recentered
$L^p$ bounds are still valid, since our forms are just shifted by adding the constant smooth term $\Theta_{F,h_F}(z)$. The probabilistic estimate 4.4 is therefore still true in exactly the same form for $L_{F,k}$, provided we use $g_{V,D,F,k}$ and $\eta_F$ instead of $g_{V,D,k}$ and $\eta$. An application of holomorphic Morse inequalities gives the desired cohomology estimates for
\begin{multline*}
h^q\Big(X,E_{k,m}V^*\langle D\rangle\otimes
\cO\Big(-\tfrac{m}{ kr}\Big(1+\tfrac{1}{ 2}+\cdots+\tfrac{1}{ k}\Big)F\Big)\Big)\\
= h^q\Big(X_k(V\langle D\rangle),\cO_{X_k(V\langle D\rangle)}(m)\otimes
\pi_k^*\cO\Big(-\tfrac{m}{ kr}\Big(1+\tfrac{1}{ 2}+\cdots+\tfrac{1}{ k}\Big)F\Big)\Big),
\end{multline*}
provided $m$ is sufficiently divisible to give a multiple of $F$ which is a $\bZ$-line bundle.
\begin{theo}\label{theo:Morse_Euler}
Let $(X,V\langle D\rangle)$ be a non-singular logarithmic directed variety,
$F\to X$ a $\bQ$-line bundle,
$(V\langle D\rangle,h)$ and $(F,h_F)$ smooth hermitian structures on $V\langle D\rangle$ and on $F$ respectively. We define
\begin{align*}
L_{F,k}&=\cO_{X_k(V\langle D\rangle)}(1)\otimes
\pi_k^*\cO\Big(-\tfrac{1}{ kr}\Big(1+\tfrac{1}{ 2}+\cdots+\tfrac{1}{ k}\Big)F\Big),\\
\eta_F&=\Theta_{\det V^*\langle D\rangle,\det h^*}-\Theta_{F,h_F} =
\Theta_{\det V^*\langle D\rangle\otimes F^{-1},\det h^*}.
\end{align*}
Then for all $q\geq 0$ and all $m\gg k\gg 1$ such that
$m$ is sufficiently divisible, we have
\begin{align*}
\tag{a} h^q(X_k(V\langle D\rangle),\cO(L_{F,k}^{\otimes m})) &\leq \tfrac{m^{n+kr-1}}{ (n+kr-1)!}\tfrac{(\log k)^n}{ n!\,k!^r}\left(
\int_{X(\eta_F,q)}(-1)^q\eta_F^n+O\left((\log k)^{-1}\right)\right),
\\
\tag{b} h^0(X_k(V\langle D\rangle),\cO(L_{F,k}^{\otimes m})) &\geq \tfrac{m^{n+kr-1}}{ (n+kr-1)!}
\tfrac{(\log k)^n}{ n!\,k!^r}\left(
\int_{X(\eta_F,\leq 1)}\eta_F^n-O\left((\log k)^{-1}\right)\right),
\\
\tag{c}
\chi(X_k(V\langle D\rangle),\cO(L_{F,k}^{\otimes m}))&=\tfrac{m^{n+kr-1}}{ (n+kr-1)!}
\tfrac{(\log k)^n}{ n!\,k!^r}\big(c_1(V^*\langle D\rangle\otimes F)^n+O\left((\log k)^{-1}\right)\big).
\end{align*}
\end{theo}

Green and Griffiths~\cite{GG80} already checked the Riemann--Roch calculation (Theorem~\ref{theo:Morse_Euler}$\,$c) in the special case $D=0$, $V=T_X^*$ and $F=\cO_X$ and prove the existence of jet differentials for surfaces of general type. Their proof is much simpler since it relies only on Chern class calculations, but it cannot provide any information on the individual cohomology groups, except in very special cases where vanishing theorems can be applied; in fact in dimension~2, the Euler characteristic satisfies $\chi=h^0-h^1+h^2\le h^0+h^2$, hence it is enough to get the vanishing of the top cohomology group
$H^2$ to infer $h^0\geq\chi\,$; this works for surfaces by means of a well-known vanishing theorem of Bogomolov which implies in general
\[
H^n\bigg(X, E_{k,m} T_X^*
\otimes
\cO\Big(-\tfrac{m}{kr}\Big(1+\tfrac{1}{2}+\dotsb+\tfrac{1}{k}\Big)F\Big)
\bigg) = 0
\]
as soon as $K_X\otimes F^{-1}$ is big and $m\gg 1$.

In fact, thanks to Bonavero's singular holomorphic Morse inequalities (Theorem~\ref{theo:Bonavero}, cf.~\cite{Bon93}), everything works almost unchanged in the case where the metric $h$ on $V$ is taken to a product
$h=h_\infty e^\varphi$ of a smooth metric $h_\infty$ by the exponential of a quasi-plurisubharmonic weight $\varphi$ with analytic singularities (so that $\det(h^*)=\det(h_\infty^*)e^{-r\varphi}$). Then $\eta$ is a $(1,1)$-current with logarithmic poles, and we just have to twist our cohomology groups by the appropriate multiplier ideal sheaves
$\cI_{k,m}$ associated with the weight $\tfrac{1}{ k}(1+\tfrac{1}{ 2}+\cdots+\tfrac{1}{ k})m\,\varphi$, since this is the multiple of $\det V^*$ that occurs in the calculation, up to the factor $\tfrac{1}{ r}\times r\varphi$. The corresponding Morse integrals need only be evaluated in the complement of the poles, i.e., on $X(\eta,q)\setminus S$ where $S=\Sing(\varphi)$. Since
\[
(\pi_k)_*\big(\cO(L_{F,k}^{\otimes m})\otimes\cI_{k,m}\big)\subset E_{k,m} V^*\otimes
\cO\Big(-\tfrac{m}{ kr}\Big(1+\tfrac{1}{ 2}+\cdots+\tfrac{1}{ k}\Big)F\Big)\Big)
\]
we still get a lower bound for the $H^0$ of the latter sheaf (or for the $H^0$ of the un-twisted line bundle $\cO(L_k^{\otimes m})$ on $\smash{X_k(V)}$). If we assume that $K_V\otimes F^{-1}$ is big, these considerations also allow us to obtain a strong estimate in terms of the volume, by using an approximate Zariski decomposition on a suitable blow-up of $X$.

\begin{coro}\label{coro5.3}
If $F$ is an arbitrary $\bQ$-line bundle over $X$, one has
\begin{multline*}
h^0\bigg(X_k(V),\cO_{X_k(V)}(m)\otimes\pi_k^*\cO
\Big(-\tfrac{m}{ kr}\Big(1+\tfrac{1}{ 2}+\cdots+\tfrac{1}{ k}\Big)F\Big)\bigg)
\\
\geq \tfrac{m^{n+kr-1}}{(n+kr-1)!}
\tfrac{(\log k)^n}{ n!\,k!^r}\left(
\Vol(K_V\otimes F^{-1})-O\left((\log k)^{-1}\right)\right)-o\left(m^{n+kr-1}\right),
\end{multline*}
$k$-jet differentials of degree $m$ twisted by the appropriate power of $F$ if $K_V\otimes F^{-1}$ is big.
\end{coro}
\begin{proof}
The volume is computed here as usual, i.e.\ after performing a suitable modifi\-cation $\mu:\smash{\widetilde X}\to X$ which converts $K_V$ into an invertible sheaf. There is of course nothing to prove if
$K_V\otimes F^{-1}$ is not big, so we can assume $\Vol(K_V\otimes F^{-1})>0$. Let us fix smooth hermitian metrics $h_0$ on $T_X$ and
$h_F$ on $F$. They induce a metric $\mu^*(\det h_0^{-1}\otimes h_F^{-1})$ on $\mu^*(K_V\otimes F^{-1})$ which, by our definition of $K_V$, is a smooth metric. By the result of Fujita [Fuj94] on approximate Zariski decomposition, for every $\delta>0$, one can find a modification $\mu_\delta:\smash{\widetilde X_\delta}\to X$ dominating
$\mu$ such that
\[
\mu_\delta^*\left(K_V\otimes F^{-1}\right) =\cO_{\widetilde X_\delta}(A+E)
\]
where $A$ and $E$ are $\bQ$-divisors, $A$ ample and $E$ effective, with
\[
\Vol(A)=A^n\geq \Vol\left(K_V\otimes F^{-1}\right)-\delta.
\]
If we take a smooth metric $h_A$ with positive definite curvature form
$\Theta_{A,h_A}$, then we get a singular hermitian metric $h_Ah_E$ on
$\mu_\delta^*(K_V\otimes F)$ with poles along $E$, i.e.\ the quotient
$h_Ah_E/\mu^*(\det h_0^{-1}\otimes h_F)$ is of the form $e^{-\varphi}$ where
$\varphi$ is quasi-psh with log poles $\log\abs{\sigma_E}^2$ (mod $C^\infty(\smash{\widetilde X_\delta}))$ precisely given by the divisor $E$. We then only need to take the singular metric $h$ on $T_X$ defined by
\[
h=h_0e^{\tfrac{1}{ r}(\mu_\delta)^*\varphi}
\]
(the choice of the factor $\tfrac{1}{ r}$ is there to correct adequately the metric on $\det V$). By construction $h$ induces an admissible metric on $V$ and the resulting curvature current $\eta_F=\Theta_{K_V,\det h^*}-\Theta_{F,h_F}$ is such that
\[
\mu_\delta^*\eta_F = \Theta_{A,h_A} +[E],\qquad
\text{$[E]={}$current of integration on $E$.}
\]
Then the $0$-index Morse integral in the complement of the poles is given by
\[
\int_{X(\eta,0)\setminus S}\eta_F^n=\int_{\widetilde X_\delta}\Theta_{A,h_A}^n=A^n\ge
\Vol\left(K_V\otimes F^{-1}\right)-\delta
\]
and Corollary~\ref{coro5.3} follows from the fact that $\delta$ can be taken arbitrary small.
\end{proof}

\begin{rema}
Since the probability estimate requires $k$ to be very large, and since all non-logarithmic components disappear from $D^{(s)}$ when $s$ is large, the above lower bound does not work in the general orbifold case. In that case, one can only hope to get an interesting result when $k$ is fixed and not too large. This is what we will do in \S~\ref{se:curvature_tangent}.
\end{rema}

\section{Positivity concepts for vector bundles and Chern inequalities}\label{se:Chern_inequalities}
\subsection{Griffiths, Nakano and strong (semi-)positivity} Let $E\to X$ be a holomorphic vector bundle equipped with a hermitian metric. Then $E$ possesses a uniquely defined Chern connection $\nabla_h$ compatible with $h$ and such that $\nabla_h^ {0,1}=\dbar$. The curvature tensor of $(E,h)$ is defined to be
\[
\Theta_{E,h}
\bydef
\tfrac{\imath}{ 2\pi}
\nabla_h^2\in C^\infty\left(X,\Lambda^{1,1}T^*_X\otimes\Hom(E,E)\right).
\]
One can then associate bijectively to $\Theta_{E,h}$ a hermitian form
$\widetilde\Theta_{E,h}$ on $TX\otimes E$, such that
\[
\widetilde\Theta_{E,h}(\xi\otimes u,\xi\otimes u)=
\langle\Theta_{E,h}(\xi,\xi)\cdot u,u\rangle_h.
\]
and can be written
\[
\Theta_{E,h}=
\tfrac{\imath}{ 2\pi}
\sum_{i,j,\lambda,\mu} c_{ij\lambda\mu}\,dz_i\wedge d\bar z_j\otimes e_\lambda^*\otimes e_\mu
\]
Let $(z_1,\,\dotsc,\,z_n)$ be a holomorphic coordinate system and let $(e_\lambda)_{1\,\leq\,\lambda\,\leq\,r}$ be a smooth frame of $e$. If $(e_\lambda)$ is chosen to be orthonormal, then we can write
\begin{align*}
\Theta_{E,h} &=
\tfrac{\imath}{ 2\pi}
\sum_{i,j,\lambda,\mu} c_{ij\lambda\mu}\,dz_i\wedge d\bar z_j\otimes e_\lambda^*\otimes e_\mu,
\\
\widetilde\Theta_{E,h}(\xi\otimes u,\xi\otimes u)&=\tfrac{1}{ 2\pi}
\sum_{i,j,\lambda,\mu} c_{ij\lambda\mu}\,\xi_i\bar\xi_j\,u_\lambda\bar u_\mu,
\end{align*}
and more generally
$\widetilde\Theta_{E,h}(\tau,\tau)=\tfrac{1}{ 2\pi} \sum_{i,j,\lambda,\mu}c_{ij\lambda\mu}\,\tau_{i\lambda}\bar \tau_{j\mu}$ for every tensor $\tau\in T_X\otimes E$. We now consider three concepts of (semi-)positivity, the first two being very classical.

\begin{defi}
Let $\theta$ be a hermitian form on a tensor product $T\otimes E$ of complex vector spaces. We say that
\begin{enumerate}
\item
$\theta$ is Griffiths semi-positive if
$\theta(\xi\otimes u,\xi\otimes u)\geq 0$ for every $\xi\in T$ and every $v\in E;$
\item
$\theta$ is Nakano semi-positive if
$\theta(\tau,\tau)\geq 0$ for every $\tau\in T\otimes E\,;$
\item
$\theta$ is strongly semi-positive if there exist a finite collection of linear forms $\alpha_j\in T^*$,
$\psi_j\in E^*$ such that $\theta=\sum_j \abs{\alpha_j\otimes\psi_j}^2$, i.e.
\[
\theta(\tau,\tau)=\sum_j\abs{(\alpha_j\otimes\psi_j)\cdot \tau}^2,\quad
\forall\tau\in T\otimes E.
\]
\end{enumerate}

Semi-negativity concepts are introduced in a similar way.
\begin{enumerate}
\item We say that the hermitian bundle $(E,h)$ is Griffiths semi-positive, resp.\ Nakano semi-positive, resp.\ strongly semi-positive, if
$\widetilde\Theta_{E,h}(x) \in \Herm(T_{X,x}\otimes E_x)$ satisfies the corresponding property for every point $x\in X$.
\item (Strict) Griffiths positivity means that
$\widetilde\Theta_{E,h}(\xi\otimes u,\xi\otimes u)>0$ for every non-zero vectors
$\xi\in T_{X,x}$, $v\in E_x$.
\item (Strict) strong positivity means that at every point $x\in X$ we can decompose $\widetilde\Theta_{E,h}$ as
$\widetilde\Theta_{E,h}=\sum_j\abs{\alpha_j\otimes\psi_j}^2$ where
$\Span(\alpha_j\otimes\psi_j)=T^*_{X,x}\otimes E^*_x$.
\end{enumerate}
\end{defi}

We will denote respectively by $\geq_G$, $\geq_N$, $\geq_S$ (and $>_G$, $>_N$, $>_S$) the Griffiths, Nakano, strong (semi-)positivity relations. It is obvious that
\[
\theta\geq_S 0~~\Rightarrow~~\theta\geq_N 0~~\Rightarrow~~\theta\geq_G 0,
\]
and one can show that the reverse implications do not hold when
$\dim T>1$ and $\dim E>1$. The following result from [Dem80] will be useful.

\begin{prop}
Let $\theta\in\Herm(T\otimes E)$, where
$(E,h)$ is a hermitian vector space. We define $\Tr_E(\theta)\in\Herm(T)$ to be the hermitian form such that
\[
\Tr_E(\theta)(\xi,\xi)=\sum_{1\,\leq\,\lambda\,\leq\,r}
\theta(\xi\otimes e_\lambda,\xi\otimes e_\lambda)
\]
where $(e_\lambda)_{1\,\leq\,\lambda\,\leq\,r}$ is an arbitrary orthonormal basis of $E$. Then
\[
\theta\geq_G 0~~\Longrightarrow~~
\theta+\Tr_E(\theta)\otimes h\geq_S 0.
\]
As a consequence, if $(E,h)$ is a Griffiths (semi-)positive vector bundle, then the tensor product
$(E\otimes \det E,h\otimes\det(h))$ is strongly (semi-)positive.
\end{prop}
\begin{proof}
Since~\cite{Dem80} is written in French and perhaps not so easy to find, we repeat here briefly the arguments. They are based on a Fourier inversion formula for discrete Fourier transforms.
\begin{lemm}\label{5.6}
Let $q$ be an integer $\geq 3$, and
$x_\alpha,~y_\beta,~1\leq\alpha,\beta\leq r$, be complex numbers. Let $\chi$ describe the set $U^r_q$ of $r$-tuples of $q^{\rm th}$ roots of unity and put
\[
\widehat x(\chi) = \sum_{1\,\leq\,\alpha\,\leq\,r} x_\alpha\overline \chi_\alpha,~~~~
\widehat y(\chi) = \sum_{1\,\leq\,\beta\,\leq\,r} y_\beta\overline\chi_\beta,~~~~\chi\in U^r_q.
\]
Then for every pair $(\lambda,\mu),~1\leq\lambda,\mu\leq r$, the following identity holds:
\[
q^{-r} \sum_{\chi\,\in\,U^r_q}
\widehat x(\chi)\,\overline{\widehat y(\chi)}\,\chi_\lambda\overline \chi_\mu =
\begin{cases}
x_\lambda\overline y_\mu &\text{if~~}\lambda\ne\mu,\\
\sum\limits_{1\,\leq\,\alpha\,\leq\,r} x_\alpha\overline y_\alpha &\text{if~~}\lambda=\mu.
\end{cases}
\]
\end{lemm}


\begin{proof}
In fact, the coefficient of $x_\alpha\overline y_\beta$ in the summation
\begin{align*}
&q^{-r} \sum_{\chi\,\in\,U^r_q}\widehat x(\chi)\,\overline{\widehat y(\chi)}\, \chi_\lambda\overline \chi_\mu
\\
\intertext{is given by}
&q^{-r} \sum_{\chi\,\in\,U^r_q}\chi_\alpha \overline \chi_\beta\overline
\chi_\lambda \chi_\mu,
\end{align*}
so it is equal to $1$ when the pairs $\{ \alpha,\mu\}$ and
$\{\beta,\lambda\}$ coincide, and is equal to $0$ otherwise. The identity stated in Lemma~\ref{5.6} follows immediately.
\end{proof}

Now, let $(t_j)_{1\,\leq\,j\,\leq\,n}$ be a basis of $T$, $(e_\lambda)_{11\,\leq1\,\lambda1\,\leq1\,r}$ an orthonormal basis of $E$ and $\xi=\sum_j\xi_jt_j\in T$, $w=\sum_{j,\lambda}w_{j\lambda}\,t_j\otimes e_\lambda \in T\otimes E$. The coefficients $c_{jk\lambda\mu}$ of $\theta$ with respect to the basis $t_j\otimes e_\lambda$ satisfy the symmetry relation $\overline c_{jk\lambda\mu}=c_{kj\mu\lambda}$, and we have the formulas
\begin{align*}
\theta(w,w) &=\sum_{j,k,\lambda,\mu}c_{jk\lambda\mu}w_{j\lambda}\overline w_{k\mu},\quad
\Tr_E\theta(\xi,\xi)=\sum_{j,k,\lambda}c_{jk\lambda\lambda}\xi_j\overline\xi_k,\\
(\theta+\Tr_E\theta\otimes h)(w,w)&=\sum_{j,k,\lambda,\mu}c_{jk\lambda\mu} w_{j\lambda} \overline w_{k\mu}+ c_{jk\lambda\lambda}w_{j\mu}\overline w_{k\mu}.
\end{align*}
For every $\chi\in U^r_q$, let us put
\[
\widehat w_j(\chi)=\sum_\alpha w_{j\alpha}\overline\chi_\alpha,\quad
\widehat w(\chi)=\sum_j\widehat w_j(\chi)\,t_j\in T\,,\quad
\widehat e_\chi=\sum_\lambda\chi_\lambda e_\lambda\in E.
\]
Lemma~\ref{5.6} implies
\begin{align*}
q^{-r} \sum_{\chi\,\in\,U^r_q}\theta(\widehat w(\chi)\otimes\widehat e_\chi,
\widehat w(\chi)\otimes\widehat e_\chi) &=q^{-r}\sum_{\chi\,\in\,U^r_q}~\sum_{j,k,\lambda,\mu}c_{jk\lambda\mu}\,
\widehat w_j(\chi)\overline{\widehat w_k(\chi)}\,\chi_\lambda \overline \chi_\mu\\
&= \sum_{j,k,\lambda\ne \mu} c_{jk\lambda\mu} w_{j\lambda} \overline w_{k\mu} +
\sum_{j,k,\lambda,\mu} c_{jk\lambda\lambda} w_{j\mu} \overline w_{k\mu}.
\end{align*}
The Griffiths positivity assumption $\theta_G\geq 0$ shows that $\xi\mapsto q^{-r}\,\theta(\xi\otimes\widehat e_\chi,\xi\otimes\widehat e_\chi)$ is a semi-positive hermitian form on $T$, hence there are linear forms $\ell_{\chi,j}\in T^*$ such that $q^{-r}\,\theta(\xi\otimes\widehat e_\chi, \xi\otimes\widehat e_\chi)=\sum_j\abs{\ell_{\chi,j}(\xi)}^2$ for all $\xi\in T$. Similarly, there are $\ell'_{\lambda,j}\in T^*$ such that
\[
\sum_{j,k} c_{jk\lambda\lambda}\,\xi_j\bar\xi_k =\sum_j\left|\ell'_{\lambda,j}(\xi)\right|^2,\quad
\text{for all }\lambda=1,\,\dotsc,\,r.
\]
Our final Fourier identity can be rewritten
\begin{align*}
(\theta + \Tr_E \theta \otimes h)(w,w) &=
\sum_{j,k,\lambda,\mu}c_{jk\lambda\mu}w_{j\lambda} \overline w_{k\mu}+
\sum_{j,k,\lambda,\mu}c_{jk\lambda\lambda}w_{j\mu} \overline w_{k\mu}
\\
&= q^{-r} \sum_{\chi\,\in\,U^r_q}\theta\big(\widehat w(\chi)\otimes\widehat e_\chi,
\widehat w(\chi)\otimes\widehat e_\chi\big)+\sum_{j,k,\lambda} c_{jk\lambda\lambda}\, w_{j\lambda}\bar w_{j\lambda}
\\
&=
\sum_{\chi\,\in\,U^r_q}\sum_j\left|\ell_{\chi,j}(\widehat w(\chi))\right|^2+\sum_{j,\lambda}
\left|\ell'_{\lambda,j}(w_{{\scriptscriptstyle\bullet},\lambda})\right|^2
\\
&=
\sum_{\chi\,\in\,U^r_q}\sum_j\left|\ell_{\chi,j}\otimes\chi^*(w)\right|^2+
\sum_{j,\lambda}\left|\ell'_{\lambda,j}\otimes e_\lambda^*(w)\right|^2
\end{align*}
where $\chi^*=\langle\bullet,\chi\rangle\in E^*$, thus $\theta + \Tr_E \theta \otimes h\geq_S 0$.
\end{proof}

\begin{coro}
Let $r=\dim E$ and $\Theta\in\Herm(T\otimes E)$.
\begin{enumerate}
\item If $\theta\geq_G 0$, then\kern27pt
$-\Tr_E\theta\otimes h~\leq_S~\theta~\leq_S~r\,\Tr_E\theta\otimes h$.
\item If $\theta\leq_G 0$, then~
$-r\,\Tr_E(-\theta)\otimes h~\leq_S~\theta~\leq_S~\Tr_E(-\theta)\otimes h$.
\item If $\pm\theta\leq_G\tau\otimes h$ where $\tau\in\Herm(T)$ is semi-positive, then
\[
-(2r+1)\,\tau\otimes h~\leq_S~\theta~\leq_S~(2r+1)\,\tau\otimes h.
\]
\end{enumerate}
\end{coro}
\begin{proof}
\begin{enumerate}
\item It is easy to check that $\theta'=\Tr_E\theta\otimes h-\theta$ satisfies $\theta'\geq_G 0$ and that we have $\Tr_E\theta'=(r-1)\Tr_E\theta$. Lemma~\ref{5.6} implies
\[
\theta'+\Tr_E\theta'\otimes h=r\,\Tr_E\theta\otimes h-\theta\geq_S 0.
\]
\item follows from (a), after replacing $\theta$ with $-\theta$.
\item also follows from Lemma~\ref{5.6} by taking $\theta'=\tau\otimes h+\theta$ (resp.\ $\theta'=\tau\otimes h-\theta$), since $\Tr_E\theta\leq r\,\tau$ and we have e.g.
\[
0\leq_S\theta'+\Tr_E\theta'\otimes h=
\theta+\Tr_E\theta\otimes h+(r+1)\tau\otimes h
\leq_S\theta+(2r+1)\tau\otimes h.
\qedhere
\]
\end{enumerate}
\end{proof}

\subsection{Chern form inequalities} In view of the estimates developed in Section~\ref{se:curvature_tangent}, we will have to evaluate integrals involving powers of curvature tensors, and the following basic inequalities will be useful.
\begin{lemm}\label{lemm:Chern_form_inequalities}
Let $\ell_j\in(\bC^r)^*$, $1\leq j\leq p$, be non-zero complex linear forms on $\bC^r$, where $(\bC^r)^*\simeq\bC^r$ is equipped with its standard hermitian form, and let $\mu$ the rotation invariant probability measure on $\bS^{2r-1}\subset\bC^r$. Then
\[
I(\ell_1,\,\dotsc,\,\ell_p)=\int_{\bS^{2r-1}}\abs{\ell_1(u)}^2\ldots\,\abs{\ell_p(u)}^2\, d\mu(u)
\]
satisfies the following inequalities:\\
\begin{eqlist}
\[
\tag{a}\label{item:Chern_form_inequality1}
I(\ell_1,\,\dotsc,\,\ell_p)
\le
\tfrac{p!\,(r-1)!}{ (p+r-1)!}\,\prod_{j=1}^p\abs{\ell_j}^2,
\]
and the equality occurs if and only if the $\ell_j$ are proportional;
\[
\tag{b}\label{item:Chern_form_inequality2}
I(\ell_1,\,\dotsc,\,\ell_p)
\ge
\tfrac{(r-1)!}{ (p+r-1)!}\,\prod_{j=1}^p\abs{\ell_j}^2,
\]
and the equality occurs if and only if $p\leq r$ and the $\ell_j$ are pairwise orthogonal.
\end{eqlist}
\end{lemm}

\begin{proof}
Denote by $d\lambda$ the Lebesgue measure on Euclidean space and by $d\sigma$ the area measure of the sphere. One can easily check that the projection
\[
\bS^{2r-1}\to \bB^{2r-2},\quad u=\left(u_1,\,\dotsc,\,u_r\right)\mapsto v=\left(u_1,\,\dotsc,\,u_{r-1}\right),
\]
yields $d\sigma(u)=d\theta\wedge d\lambda(v)$ where $u_r=\abs{u_r}\,e^{i\theta}$ [just check that the wedge products of both sides with $\tfrac{1}{ 2}d\abs{u}^2$ are equal to $d\lambda(u)$, and use the fact that $d\theta=\tfrac{1}{ 2i}(du_r/u_r -d\bar u_r/\bar u_r)$], thus, in terms of polar coordinates $v=t\,u'$, $u'\in\bS^{2r-1}$, we have
$d\sigma(u)=d\theta\wedge t^{2r-3}\,dt\wedge d\sigma'(u')$, and going back to the invariant probability measures $\mu$ on $\bS^{2r-1}$ and $\mu'$ on $\bS^{2r-3}$, we get
$\abs{u_r}^2=1-\abs{v}^2=1-t^2$ and an equality
\begin{equation}\label{eq:measure}
d\mu(u)=\tfrac{2r-2}{ 2\pi}\,d\theta\wedge t^{2r-3}\,dt
\wedge d\mu'(u').
\end{equation}
If $\ell_1,\,\dotsc,\,\ell_p$ are independent of $u_r$, \eqref{eq:measure} and the Fubini theorem imply by homogeneity
\begin{align}\label{5.10}
\int_{\bS^{2r-1}}\abs{\ell_1(u')}^2\ldots\,\abs{\ell_p(u')}^2\,d\mu(u) =
\tfrac{r-1}{p+r-1}
\int_{\bS^{2r-3}}\abs{\ell_1(u')}^2\ldots\,\abs{\ell_p(u')}^2\,d\mu'(u'),
\end{align}
\begin{multline}\label{5.10'}
\tag{\theequation\ensuremath{'}}
\int_{\bS^{2r-1}}\abs{\ell_1(u')}^2\ldots\,\abs{\ell_{p-1}(u')}^2\,\abs{u_r}^2\,d\mu(u) =
\\
\tfrac{r-1}{(p+r-2)(p+r-1)} \int_{\bS^{2r-3}}\abs{\ell_1(u')}^2\ldots\,\abs{\ell_{p-1}(u')}^2\,d\mu'(u')
\end{multline}
(for instance, in case~\eqref{5.10'}, we have to integrate $t^{2p-2}(1-t^2)\times t^{2r-3}\,dt$). For $p\leq r$, the formulas
\[
\int_{\bS^{2r-1}}\abs{u_1}^{2p}\,d\mu(u)=\tfrac{p!\,(r-1)!}{ (p+r-1)!},\quad
\int_{\bS^{2r-1}}\abs{u_1}^2\ldots\,\abs{u_p}^2\,d\mu(u)=\tfrac{(r-1)!}{ (p+r-1)!},
\]
are then obtained by induction on $r$ and $p$.
\begin{enumerate}
\item For any $\ell\in(\bC^r)^*$, we can find orthonormal coordinates on $\bC^r$ such that $\ell(u)=\abs{\ell}\,u_1$ in the new coordinates. Hence
\[
\int_{\bS^{2r-1}}\abs{\ell(u)}^{2p}\,d\mu(u)=m_{r,p}\,\abs{\ell}^{2p}\quad
\text{where}~~m_{r,p}=\int_{\bS^{2r-1}}\abs{u_1}^{2p}\,d\mu(u)=
\tfrac{p!\,(r-1)!}{ (p+r-1)!}.
\]
It follows from H\"older's inequality that
\[
I(\ell_1,\,\dotsc,\,\ell_p)\leq \prod_{j=1}^p
\left(\int_{\bS^{2r-1}}\abs{\ell_j}^{2p}\,d\mu(u)\right)^{1/p}=m_{r,p}
\prod_{j=1}^p\abs{\ell_j}^2,
\]
and that the equality occurs if and only if all $\ell_j$ are proportional.

\item We prove the inequality
\[
I(\ell_1,\,\dotsc,\,\ell_p)\geq \tfrac{(r-1)!}{(p+r-1)!}\prod_{j=1}^p\abs{\ell_j}^2
\]
by induction on $p$, the result being clear for $p=0$ or $p=1$. If we choose an orthonormal basis $(e_1,\,\dotsc,\,e_r)\in\bC^r$ such that
$\ell_j(e_r)\ne 0$ for all $j$ and replace $\ell_j$ by
$(\ell_j(e_r))^{-1}\ell_j$, we can assume $\ell_j(e_r)=1$. We then write $u=u'+u_re_r$ with $u'\in e_r^\perp\simeq\bC^{r-1}$ and
\[
\ell_j(u)=\ell_j'(u')+u_r,\quad 1\leq j\leq p,\quad\ell'_j
\bydef
\ell_{j|e_r^\perp}.
\]
Let $s_k(\ell'_\bullet(u'))$ be the elementary symmetric functions in $\ell'_j(u')$, $1\leq j\leq p$, with $s_0\bydef1$. We have
\[
I(\ell_1,\,\dotsc,\,\ell_p) =\int_{\bS^{2r-1}}\prod_{j=1}^p\left|\ell'_j(u')+u_r\right|^2\,d\mu(u) =\int_{\bS^{2r-1}}
\abs*{\sum_{k=0}^ps_k(\ell'_\bullet(u'))\,u_r^{p-k}}^2 d\mu(u).
\]
We make a change of variable $u_r\mapsto u_r\,e^{i\theta}$ and take the average over $\theta\in[0,2\pi]$. Parseval's formula gives
\[
I(\ell_1,\,\dotsc,\,\ell_p) =\int_{\bS^{2r-1}}\sum_{k=0}^p\abs*{s_k(\ell'_\bullet(u'))}^2\,\abs{u_r}^{2(p-k)} d\mu(u),
\]
and since
\[
(2r-2)\int_0^1t^{2k}(1-t^2)^{p-k}\,t^{2r-3}dt=
\tfrac{(r-1)\,(k+r-2)!\,(p-k)!}{(p+r-1)!},
\]
formula~\eqref{eq:measure} implies
\[
I(\ell_1,\,\dotsc,\,\ell_p) =\int_{\bS^{2r-3}}\sum_{k=0}^p\tfrac{(r-1)\,(k+r-2)!\,(p-k)!}{(p+r-1)!}\,
\abs*{s_k(\ell'_\bullet(u'))}^2\,d\mu'(u').
\]
As $\abs{\ell_j}^2=1+\abs{\ell'_j}^2$, our inequality~\ref{lemm:Chern_form_inequalities}\eqref{item:Chern_curv_inequality2} is equivalent to
\begin{equation}\label{eq:basic_inequality}
\int_{\bS^{2r-3}}\sum_{k=0}^p\tfrac{(k+r-2)!\,(p-k)!}{(r-2)!}\,
\abs*{s_k(\ell'_\bullet(u'))}^2\,d\mu'(u')\ge
\prod_{j=1}^p\left(1+\abs{\ell'_j}^2\right)
\end{equation}
for all linear forms $\ell'_j\in(\bC^{r-1})^*$. We actually prove~\eqref{eq:basic_inequality} by induction on $p$ (observing that the inequality is a trivial equality for $p=0,1$). Assume that~\eqref{eq:basic_inequality} (and hence~\ref{lemm:Chern_form_inequalities}\eqref{item:Chern_curv_inequality2}) is known for any $(p-1)$-tuple of linear forms $(\ell'_1,\,\dotsc,\,\ell'_{p-1})$. As~\ref{lemm:Chern_form_inequalities}\eqref{item:Chern_curv_inequality2} is invariant under the action of $U(r)$, it is sufficient to consider the case when $\ell_p(u)=u_r$, i.e.\ $\ell'_p=0$. The induction hypothesis tells us that
\[
\int_{\bS^{2r-3}}\sum_{k=0}^{p-1}\tfrac{(k+r-2)!\,(p-1-k)!}{(r-2)!}\,
\abs*{s_k(\ell'_\bullet(u'))}^2\,d\mu'(u')\ge
\prod_{j=1}^{p-1}(1+\abs{\ell'_j}^2).
\]
However, when we add the factor $\ell_p$, the elementary symmetric functions $s_k(\ell'_\bullet(u'))$ are left unchanged for $k\leq p-1$, while
$s_p(\ell'_\bullet(u'))=0$ and $1+\abs{\ell'_p}^2=1$. Therefore~\eqref{eq:basic_inequality} holds true for $p$, since $(p-k)!\geq (p-1-k)!$ for all $k=0,1,\dotsc,p-1$. We have proved the inequality at order $p$ whenever
$\ell_p=\alpha_p\langle\bullet,e_r\rangle$ and $\ell_j(e_r)\ne 0$ for $j\leq p-1$. Since those $(\ell_1,\,\dotsc,\,\ell_p)$ are dense in the space $((\bC^r)^*)^p$ of $p$-tuples of linear forms, the proof of the lower bound is complete.

\item (b, equality case) We argue by induction on $r$. For $r=1$, we have in fact $\ell_j(u)=\alpha_ju_1$, $\alpha_j\in\bC^*$, and
$I(\ell_1,\,\dotsc,\,\ell_r)=\prod\abs{\ell_j}^2$, thus the coefficient $\tfrac{1}{ (p+r-1)!}=\tfrac{1}{ p!}$ is reached if and only if $p\leq 1$. Now, assume $r\geq 2$ and the equality case solved for dimension $r-1$. By rescaling and reordering the $\ell_j$, we can always assume that $\ell_j(e_r)\ne 0$ (and hence
$\ell_j(e_r)=1$) for $q+1\leq j\leq p$, while $\ell_j(e_r)=0$ for
$1\leq j\leq q$ (we can possibly have $q=0$ here). Then we write
$\ell_j(u)=\ell'_j(u')$ for $1\leq j\leq q$ and
$\ell_j(u)=\ell'_j(u')+u_r$ for $q+1\leq j\leq p$. Therefore, if
$s_k(\ell'(u'))$ denotes the $k^{\rm th}$ elementary symmetric function in
$(\ell'_j(u')_{q+1\leq j\leq p}$, we find
\begin{align*}
I(\ell_1,\,\dotsc,\,\ell_p) &=\int_{\bS^{2r-1}}\prod_{j=1}^q\left|\ell'_j(u')\right|^2\,
\prod_{j=q+1}^p\left|\ell'_j(u')+u_r\right|^2\,d\mu(u)\\
&=\int_{\bS^{2r-1}}\prod_{j=1}^q\left|\ell'_j(u')\right|^2
\left|\sum_{k=0}^{p-q}s_k(\ell'(u'))\,u_r^{p-q-k}\right|^2 d\mu(u)
\\
&=\int_{\bS^{2r-1}}\prod_{j=1}^q\left|\ell'_j(u')\right|^2
\sum_{k=0}^{p-q}
\left|s_k(\ell'(u'))\right|^2\,\abs{u_r}^{2(p-q-k)}\,d\mu(u)\\
&=\int_{\bS^{2r-3}}\prod_{j=1}^q\left|\ell'_j(u')\right|^2
\sum_{k=0}^{p-q}\tfrac{(r-1)\,(k+r-2)!\,(p-q-k)!}{(p-q+r-1)!}\,
\left|s_k(\ell'(u'))\right|^2\,d\mu'(u')\\
&\geq\tfrac{(r-1)!}{(p+r-1)!}\,\prod_{j=1}^q\abs{\ell'_j}^2
\prod_{j=q+1}^p\left(1+\abs{\ell'_j}^2\right)
\end{align*}
by what we have just proved. In an equivalent way, we get
\begin{multline*}
\int_{\bS^{2r-3}}\prod_{j=1}^q\left|\ell'_j(u')\right|^2
\sum_{k=0}^{p-q}\tfrac{(k+r-2)!\,(p-q-k)!\,(p+r-1)!}{(r-2)!\,(p-q+r-1)!}\,
\left|s_k(\ell'(u'))\right|^2\,d\mu'(u')
\\
\geq\prod_{j=1}^q\abs{\ell'_j}^2
\prod_{j=q+1}^p\left(1+\abs{\ell'_j}^2\right)
\end{multline*}
In general, we can rotate coordinates in such a way that $\ell_p(u)=u_r$ and $\ell'_p=0$, and we see that the above inequality holds when $p$ is replaced by $p-1$, as soon as $q\leq p-2$. Then the corresponding coefficients $k=0$ for $p$, $p-1$ are
\[
\tfrac{(p-q)!\,(p+r-1)!}{(p-q+r-1)!}>\tfrac{(p-1-q)!\,(p-1+r-1)!}{(p-1-q+r-1)!},
\]
and since $s_0=1$, we infer that the inequality is strict. The only possibility for the equality case is $q=p-1$, but then
\[
I(\ell_1,\,\dotsc,\,\ell_p)=
\int_{\bS^{2r-1}}\prod_{j=1}^{p-1}\left|\ell'_j(u')\right|^2\,\abs{u_r}^2\,d\mu(u)=
\tfrac{r-1}{ p+r-1}\int_{\bS^{2r-3}}\prod_{j=1}^{p-1}\left|\ell'_j(u')\right|^2\,d\mu'(u'),
\]
and we see that we must have equality in the case $(r-1,p-1)$. By induction, we conclude that $p-1\leq r-1$ and that the
$\ell_j(u)=\ell'_j(u')$ are orthogonal for $j\leq p-1$, as desired.\qedhere
\end{enumerate}
\end{proof}

\begin{rema}
When $r=2$, our inequality~\eqref{eq:basic_inequality} is equivalent to the ``elementary'' inequality
\[
\tag{*}
\prod_{j=1}^p\left(1+\abs{a_j}^2\right)\leq \sum_{k=0}^p k!\,(p-k)!\,\abs{s_k}^2,
\]
relating a polynomial $X^p-s_1X^{p-1}+\cdots+(-1)^ps_p$ and its complex roots $a_j$ (just consider $\ell'_j(u')=a_ju_1$ and $\ell_j(u)=a_ju_1+u_2$ on $\bC^2$ to get this). It should be observed that $(*)$ is not optimal asymptotically when $p\to+\infty\,$; in fact, Landau's inequa\-lity~\cite{Land05} gives
$\prod\max(1,\abs{a_j})\leq(\sum\abs{s_k}^2)^{1/2}$, from which one can easily derive that $\prod(1+\abs{a_j}^2)\leq 2^p\sum\abs{s_k}^2$, which improves
$(*)$ as soon as $p\geq 7$ (observe that $2^7=128$ and
$k!(7-k)!\geq 3!\,4!=144$). Our discussion of the equality case shows that inequality~\eqref{item:Chern_curv_inequality2} from Lemma~\ref{lemm:Chern_form_inequalities} is never sharp when $p>r$. It would be interesting, but probably challenging, if~not impossible, to compute the optimal constant for all pairs $(r,p)$, $p>r$, since this is an optimization problem involving the distribution of a large number of points in projective space.
\end{rema}

We finally state one of the main consequences of these estimates concerning the Chern curvature form of a hermitian holomorphic vector bundle.

\begin{prop}\label{prop:Chern_curv_inequalities}
Let $T$, $E$ be complex vector spaces of respective dimensions $\dim T=n$, $\dim E=r$. Assume that $E$ is equipped with a hermitian structure $h$, and denote by $\mu$ the unitary invariant probability measure $\mu$ on the unit sphere bundle $S(E)=\{u\in E\,;\abs{u}_h=1\}$ of $E$.
\begin{enumerate}\alphenumi
\item\label{item:Chern_curv_inequality1}
If $\ell_1,\,\dotsc,\,\ell_k\in E^*$ and $\theta_1,\dotsc,\theta_{p-k}\geq_S 0$ are strongly semi-positive hermitian tensors in
$\Herm(T\otimes E)\simeq\Lambda_\bR^{1,1}T^*\otimes_\bR\Herm(E,E)$, then
\begin{multline*}
\int_{u\,\in\,S(E)}\abs{\ell_1(u)}^2\ldots\abs{\ell_k(u)}^2\,
\langle\theta_1(u),u\rangle_h\wedge\ldots
\wedge\langle\theta_{p-k}(u),u\rangle_h\,d\mu(u)
\\
\begin{cases}
\displaystyle
\geq \tfrac{(r-1)!}{ (p+r-1)!}\,\bigg(\prod_{j=1}^k\abs{\ell_j}^2\bigg)
\Tr_h\theta_1\wedge\ldots\wedge\Tr_h\theta_{p-k},\\
\displaystyle
\leq \tfrac{p!\,(r-1)!}{ (p+r-1)!}\,\bigg(\prod_{j=1}^k\abs{\ell_j}^2\bigg)
\Tr_h\theta_1\wedge\ldots\wedge\Tr_h\theta_{p-k},
\end{cases}
\end{multline*}
as pointwise strong inequalities of $(p-k,p-k)$-forms.
\item\label{item:Chern_curv_inequality2}
If $\theta\geq_G 0$ in $\Lambda_\bR^{1,1}T^*\otimes_\bR\Herm(E,E)$ and $\ell_j\in E^*$, then
\[
\int_{u\,\in\,S(E)}\abs{\ell_1(u)}^2\ldots\abs{\ell_k(u)}^2\,
\langle\theta(u),u\rangle_h^{p-k}\,d\mu(u)\le
\tfrac{p!\,(r-1)!}{ (p+r-1)!}\bigg(\prod_{j=1}^k\abs{\ell_j}^2\bigg)(\Tr_h\theta)^{p-k}
\]
as a pointwise weak inequality of $(p-k,p-k)$-forms.
\end{enumerate}
In particular, the above inequalities apply when
$(E,h)$ is a hermitian holomorphic vector bundle of rank $r$ on a complex $n$-dimensional manifold $X$, and one takes
$\theta_j=\Theta_{E,h}$ to be the curvature tensor of $E$, so that
$\Tr_h\theta_j=c_1(E,h)$ is the first Chern form of $(E,h)$.
\end{prop}
\begin{proof}
\eqref{item:Chern_curv_inequality1} The assumption $\theta_q\geq_S 0$ means that at every point
$x\in X$ we can write $\theta$ as
\[
\theta_q=\sum_{1\,\leq\,j\,\leq\,N_q}\abs{\beta_{qj}\otimes\ell_{qj}}^2
\simeq \sum_{1\,\leq\,j\,\leq\,N_q}\imath\beta_{qj}\wedge\bar\beta_{qj}
\otimes\ell_{qj}\otimes\ell_{qj}^*,\quad\beta_{qj}\in T^*,~~\ell_{qj}\in E^*
\]
as an element of $\Lambda_\bR^{1,1}T^*\otimes_\bR\Herm(E,E)$, hence
\[
\left\langle\theta_q(u),u\right\rangle_h=\sum_{1\,\leq\,j\,\leq\, N_q}
\imath\beta_{qj}\wedge\bar\beta_{qj}\,\abs{\ell_{qj}(u)}^2.
\]
Without loss of generality, we can assume $\abs{\ell_{qj}}_{h^*}=1$. Then
\begin{multline*}
\abs{\ell_1(u)}^2\ldots\abs{\ell_k(u)}^2\,\langle\theta_1(u),u\rangle_h\wedge
\ldots\wedge\langle\theta_{p-k}(u),u\rangle_h
\\
=\sum_{j_1,\,\dotsc,\,j_{p-k}}
\imath\beta_{1j_1}\wedge\bar\beta_{1j_1}\wedge\ldots\wedge
\imath\beta_{p-k\,j_{p-k}}\wedge\bar\beta_{p-k\,j_{p-k}}\,
\prod_{1\,\leq\,s\,\leq\,k}\abs{\ell_s(u)}^2\prod_{1\,\leq\,s\,\leq\,p-k}\abs{\ell_{sj_s}(u)}^2,
\end{multline*}
and since $\abs{\ell_{qj}}_{h^*}=1$, Lemma~\ref{lemm:Chern_form_inequalities}\eqref{item:Chern_curv_inequality2} implies
\begin{align*}
\int_{u\,\in\,S(E)}&\abs{\ell_1(u)}^2\ldots\abs{\ell_k(u)}^2\,
\langle\theta_1(u),u\rangle_h\wedge\ldots\wedge
\langle\theta_{p-k}(u),u\rangle_h\,d\mu(u)\\
&\geq \tfrac{(r-1)!}{(p+r-1)!}~
\sum_{j_1,\,\dotsc,\,j_{p-k}}\imath\beta_{1j_1}\wedge\bar\beta_{1j_1}
\wedge\ldots\wedge\imath\beta_{p-k\,j_{p-k}}\wedge\bar\beta_{p-k\,j_{p-k}}\,
\prod_{1\leq s\leq k}\abs{\ell_s}^2\\
&=\tfrac{(r-1)!}{(p+r-1)!}\,\left(\prod_{1\,\leq\,j\,\leq\,k}\abs{\ell_j}^2\right)
\Tr_h\theta_1\wedge\ldots\wedge\Tr_h\theta_p,
\end{align*}
where $\geq$ is in the sense of the strong positivity of $(p,p)$-forms. The upper bound is obtained by the same argument, via Lemma~\ref{lemm:Chern_form_inequalities}\eqref{item:Chern_curv_inequality1}.


\eqref{item:Chern_curv_inequality2} By the definition of weak positivity of forms, it is enough to show the inequality in restriction to every $(p-k)$-dimensional subspace
$T'\subset T$. Without loss of generality, we can assume that $\dim T=p-k$ (and then take $T'=T$), that $\abs{\ell_j}=1$, and also that
$\theta>_G0$ (otherwise take a positive definite form
$\eta\in\Lambda_\bR^{1,1}T^*$, replace $\theta$ with $\theta_\varepsilon=\theta+\varepsilon\,\eta\otimes h$, and let $\varepsilon$ tend to $0$). For any $u\in S(E)$, let
\[
0\leq\lambda_1(u)\leq\cdots\leq\lambda_{p-k}(u)
\]
be the eigenvalues of the hermitian form
$q_u(\bullet)=\langle \theta(u),u\rangle$ on $T$ with respect to
\[
\omega=\Tr_h\theta=\sum_{j=1}^r\langle \theta(e_j),e_j\rangle\in
\Herm(T),\quad\omega>0,
\]
$(e_j)$ being any orthonormal frame of $E$. We have to show that
\[
\int_{u\,\in\,S(E)}
\abs{\ell_1(u)}^2\ldots\abs{\ell_k(u)}^2\,\lambda_1(u)\cdots\lambda_{p-k}(u)\,d\mu(u)\le
\tfrac{p!\,(r-1)!}{(p+r-1)!}.
\]
However, the inequality between geometric and arithmetic means implies
\[
\lambda_1(u)\cdots\lambda_p(u)\le
\bigg(\tfrac{1}{ p-k}\sum_{j=1}^{p-k}\lambda_j(u)\bigg)^p,
\]
thus, putting $Q(u)=\tfrac{1}{ p-k}\langle\Tr_\omega\theta(u),u\rangle$,
$Q\in\Herm(E)$, it is enough to prove that
\begin{equation}\label{5.14}
\int_{u\,\in\,S(E)}\abs{\ell_1(u)}^2\ldots\abs{\ell_k(u)}^2\,Q(u)^{p-k}\,d\mu(u)
\leq \tfrac{p!\,(r-1)!}{(p+r-1)!}.
\end{equation}
Our assumption $\theta>_G0$ implies
$Q(u)=\sum_{1\,\leq\,j\,\leq\,r} c_j\abs{\ell'_{qj}(u)}^2$ for some $c_j>0$ and some orthonormal basis $(\ell'_{qj})_{1\,\leq\,j\,\leq\,r}$ of
$E^*$, and
\[
\sum_{j=1}^r c_j=\Tr_hQ=\tfrac{1}{ p-k}\Tr_h(\Tr_\omega\theta) =\tfrac{1}{ p-k}\Tr_\omega(\Tr_h\theta)=\tfrac{1}{ p-k}\Tr_\omega(\omega)=1.
\]
Inequality~\eqref{5.14} is a consequence of Lemma~\ref{lemm:Chern_form_inequalities}\eqref{item:Chern_curv_inequality1}, by Newton's multinomial expansion.
\end{proof}

\begin{rema}\label{5.15}
For $p=1$, the inequalities of Proposition~\ref{prop:Chern_curv_inequalities} are identities, and no semi-positivity assumption is needed in that case. This can be seen directly from the fact that we have
\[
\int_{u\,\in\,S(E)}Q(u)\,d\mu(u)=\tfrac{1}{ r}\,\Tr Q
\]
for every hermitian quadratic form $Q$ on $E$. However, when $p\geq 2$, inequality~\ref{prop:Chern_curv_inequalities}~\eqref{item:Chern_curv_inequality1} does not hold under the assumption that
$E\geq_G0$ (or even that $E$ is dual Nakano semi-positive, i.e.\ $E^*$ Nakano semi-negative). Let us take for instance $E=T_{\bP^n}\otimes\cO(-1)$. It is well known that
$E$ is isomorphic to the tautological quotient vector bundle
$\bC^{n+1}/\cO(-1)$ over $\bP^n$, and that its curvature tensor form for the Fubini-Study metric is given by
\[
\Theta_{E,h}(\xi\otimes u,\xi\otimes u)=
\abs{\langle\xi,u\rangle}^2\geq 0
\]
(where $v$ is identified which a tangent vector via the choice of a unit element $e\in\cO(-1)$). Then $\det E=\cO(1)$ and thus
$c_1(E,h)=\omega_\FS>0$, although $\langle\Theta_{E,h}(u),u\rangle_h^p=0$ for all $p\geq 2$, as one can easily check.
\end{rema}

\section{On the curvature of orbifold tangent bundles}\label{se:curvature_tangent}
\subsection{Evaluation of the orbifold curvature tensor}\label{sse:orbi_metric}
The main qualitative result is summarized in the following statement.
\begin{prop}\label{prop:orbi_metric}
Let $X$ be a projective variety, $A$ an ample line bundle, and $(X,V,D)$ an orbifold directed structure where
$D=\sum_{1\,\leq\,j\,\leq\,N}(1-\tfrac{1}{\rho_j})\Delta_j$ is a normal crossing divisor transverse to $V$ in $X$. Let $d_j$ be the infimum of numbers
$\lambda\in\bR_+$ such that $\lambda A-\Delta_j$ is~nef, and $\gamma_V$ be the infimum of numbers $\gamma\geq 0$ such that $\gamma\,
\Theta_{A,h_A}\otimes\Id_V-\Theta_{V,h_V}\geq_G0$ for suitable smooth hermitian metrics $h_V$ on $V$. Then for every
$\gamma>\gamma_{V,D}\bydef\max(\max_j(d_j/\rho_j),\gamma_V)$, the~orbifold vector bundle $V\langle D\rangle$ possesses a hermitian metric $h_{V\langle D\rangle,\gamma,\varepsilon}$ such that
\begin{enumerate}
\item
$h_{V\langle D\rangle,\gamma,\varepsilon}$ is smooth on $X\setminus\abs{D},$
\item
$h_{V\langle D\rangle,\gamma,\varepsilon}$ has the appropriate orbifold singularities along $D,$
\item we have $\gamma\,\Theta_{A,h_A}\otimes\Id-
\Theta_{V\langle D\rangle,h_{V\langle D\rangle,\gamma,\varepsilon}}\geq_G0$ on $X\setminus\abs{D}$.
\end{enumerate}
\end{prop}
\begin{proof}
Let $h_A$ be a metric on $A$ such that $\Theta_{A,h_A}>0$, written locally as $h_A=e^{-\psi}$, and take
\text{$\gamma>\max(\max_j(d_j/\rho_j),\gamma_V)$}. Consider the tautological sections
$\sigma_j\in H^0(X,\cO_X(\Delta_j))$ defining $\Delta_j=\sigma_j^{-1}(0)$, and let $h_V$, $h_j$ be smooth hermitian metrics on
$V$ and $\cO_X(\Delta_j)$ such that
\begin{align}
\stepcounter{equation}
\tag{$\theequation_0$} &
\gamma\,\Theta_{A,h_A}\otimes\Id_V-\Theta_{V,h_V}>_G0,
\\
\tag{$\theequation_j$}\label{6.2_j}
&
\gamma\,\Theta_{A,h_A}-\tfrac{1}{\rho_j}\Theta_{\cO_X(\Delta_j),h_j}>0,\quad
\forall j=1,\dotsc,N,
\end{align}
as is possible by our choice of the constants $d_j$ and $\gamma$. Finally, denote by $\nabla_j$ the associated Chern connection on
$\cO_X(\Delta_j)$. If we write $h_j=e^{-\varphi_j}$ in some local trivialization, then $\nabla_j\sigma_j=\nabla_j^{1,0}\sigma_j=
\partial\sigma_j-\sigma_j\partial\varphi_j$. Take $\omega_A=\Theta_{A,h_A}$ as the K\"ahler metric on $X$. We have
\[
\imath\ddbar\abs{\sigma_j}_{h_j}^{2/\rho_j}=
\tfrac{1}{\rho_j^2}\,\abs{\sigma_j}_{h_j}^{-2+2/\rho_j}\, i\left\langle\nabla_j\sigma_j,\nabla_j\sigma_j\right\rangle_{h_j} -\tfrac{1}{\rho_j}\abs{\sigma_j}_{h_j}^{2/\rho_j}\,\imath\ddbar\varphi_j,
\]
hence there exists $\delta>0$ small such that the metric
$h_{A,\delta}=h_A\exp(-\delta\sum_j\abs{\sigma_j}_{h_j}^{2/\rho_j})$ of weight $\psi_\delta=\psi+\delta\sum_j\abs{\sigma_j}_{h_j}^{2/\rho_j}$ satisfies
\begin{multline*}
\imath\ddbar\psi_\delta(\xi,\xi) =
\\
\abs{\xi}^2_{\omega_A} +
\delta\imath\ddbar
\sum_j\abs{\sigma_j}_{h_j}^{2/\rho_j}(\xi,\xi)
\ge (1-C\delta)\abs{\xi}^2_{\omega_A} +
\delta\sum_j\tfrac{1}{\rho_j^2}\abs{\sigma_j}_{h_j}^{-2+2/\rho_j}
\left|\nabla_j\sigma_j(\xi)\right|_{h_j}^2.
\end{multline*}
We can consider $\omega_{A,\delta}=\Theta_{A,h_{A,\delta}}=\imath\ddbar\psi_\delta$ as an orbifold K\"ahler metric, that is ``smooth'' from the point of view of the orbifold structure. Let us explain the more precise meaning of this ``orbifold smoothness'' assumption. In fact, there exists a ramified cover $g_Y\colon Y\to X$ such that $g^*\sigma_j=w_j^{m_j}$ for some local coordinate $w_j$ on $Y$, with arbitrary high multiplicity $m_j\in\bN^*$ along $\smash{g_Y^{-1}}(\Delta_j)=\{w_j=0\}$. Then $g_Y^*h_{A,\delta}=g_Y^*h_A\exp(-\delta\sum_j \smash{\abs{w_j}^{2m_j/\rho_j})}$ can be taken in any regularity class $C^p$, $p\in \bN^*$, by taking $m_j\geq p\,\rho_j$. Therefore, by pulling-back our calculations to $Y$, we would actually get forms of high regularity on $Y$. Of course, if we compute an integral over $X$, pulling-back forms to $Y$ multiplies the integral by the degree of $g_Y$, and it suffices to divide by that degree to recover the integral over $X$. For $\delta>0$ sufficiently small, our positivity conditions~\eqref{6.2_j} can be turned into the stronger form
\begin{align}
\stepcounter{equation}
\tag{$\theequation_0$}
\gamma\,\imath\ddbar\psi_\delta(\xi,\xi)\,\abs{u}^2-\widetilde\Theta_{V,h_V}(\xi\otimes u) &\ge c\left(\abs{\xi}_{\omega_A}^2+\sum_j\abs{\sigma_j}_{h_j}^{-2+2/\rho_j}\,\left|\nabla_j\sigma_j(\xi)\right|_{h_j}^2
\right),
\\
\tag{$\theequation_j$}\label{6.3_j}
\gamma\,\imath\ddbar\psi_\delta(\xi,\xi)-\tfrac{1}{\rho_j}\,\imath\ddbar\varphi_j(\xi,\xi) &\ge c\left(\abs{\xi}_{\omega_A}^2+\sum_j\abs{\sigma_j}_{h_j}^{-2+2/\rho_j}\,\left|\nabla_j\sigma_j(\xi)\right|_{h_j}^2
\right),
\end{align}
for some constant $c>0$ and all $\xi\in T_X$, $u\in V$ (observe that the right hand side can in fact be seen as a positive definite hermitian form with respect to the orbifold coordinates, we just exploit the fact that
$A$ remains ample when viewed as a line bundle on the orbifold structure). We are going to estimate the curvature of the orbifold metric
$h_{V\langle D\rangle,\varepsilon}$ on $V\langle D\rangle$ defined by
\begin{equation}\label{eq:orbi_metric}
\Abs{u}_{h_{V\langle D\rangle,\varepsilon}}^2 =
\abs{u}_{h_V}^2+\sum_j\varepsilon_j\,
\abs{\sigma_j}_{h_j}^{-2(1-1/\rho_j)}\,\left|\nabla_j\sigma_j(u)\right|_{h_j}^2,
\quad \varepsilon_j\ll 1.
\end{equation}
Again, this metric can be seen as orbifold smooth (in the sense that the metric $g_Y^*h_{V\langle D\rangle,\varepsilon}$ on $g_Y^*(V\langle D\rangle)$ may be taken of arbitrary high regularity; in case $\rho_j=\infty$, it is actually a smooth metric on the logarithmic bundle). Since
\[
\imath\ddbar\Abs{u}_{h_{V\langle D\rangle},\varepsilon}^2=
\imath\langle\nabla u,\nabla u\rangle_{h_{V\langle D\rangle},\varepsilon}-2\pi\,
\left\langle\Theta_{V\langle D\rangle,h_{V\langle D\rangle,\varepsilon}}(u),u
\right\rangle_{h_{V\langle D\rangle,\varepsilon}}
\]
where $\nabla u=du+\Gamma(dz)\cdot u$ is the Chern connection of
$(V\langle D\rangle,h_{V\langle D\rangle,\varepsilon})$, what we need to prove is that on the total space of $V$ over $X\setminus\abs{D}$, the $(1,1)$-form
\[
V\ni (z,u)\mapsto
\imath\ddbar\Abs{u}_{h_{V\langle D\rangle,\varepsilon}}^2+\gamma\,
\imath\ddbar\psi_\delta\,\Abs{u}_{h_{V\langle D\rangle,\varepsilon}},
\]
is non-negative. For this, we calculate the associated hermitian quadratic form on~$T_V$
\[
Q_{V\langle D\rangle,\gamma,\varepsilon}(z,u)(\xi,\eta),\quad (\xi,\eta)\in T_{V,(z,u)},\quad
\xi=\sum_{\ell=1}^n\xi_\ell\tfrac{\partial}{\partial z_\ell},\quad
\eta=\sum_{\lambda=1}^r\eta_\lambda\tfrac{\partial}{\partial u_\lambda},
\]
and observe that the curvature tensor is obtained by taking the restriction to the ``parallel'' directions $\nabla u=0$, that is, by substituting $du=-\Gamma(dz)\cdot u$, i.e.\
$\eta=-\Gamma(\xi)\cdot u$. Let us fix an arbitrary point
$z_0\in X\setminus\abs{D}$. We take local holomorphic coordinates $(z_1,\,\dotsc,\,z_n)$ centered at $z_0$, and let $(e_1,\,\dotsc,\,e_r)$ be a local holomorphic frame of $V$ such that
\[
\langle e_\lambda,e_\mu\rangle_{h_V}=\delta_{\lambda\mu}+
\sum_{\ell,m,\lambda,\mu}c_{\ell m\lambda\mu}\,z_\ell\bar z_m+O(\abs{z}^3),
\]
where the $\tfrac{\imath}{ 2\pi}c_{\ell m\lambda\mu}$ are the coefficients of
$-\Theta_{V,h_V}$. Let us write $u=\sum_{\lambda=1}^ru_\lambda e_\lambda$ and denote by $\langle u,v\rangle=\sum_{1\,\leq\,\lambda\,\leq\,r}u_\lambda\bar v_\lambda$ the standard hermitian form, $\abs{u}$ the associated norm. We find
\begin{multline*}
\tag{$\ref{eq:orbi_metric}_{0}$}\label{6.4_0}
\Abs{u}_{h_{V\langle D\rangle,\varepsilon}}^2 =
\abs{u}^2+\sum_{\ell,m,\lambda,\mu}c_{\ell m\lambda\mu}\,z_\ell\bar z_m u_\lambda\bar u_\mu+O(\abs{z}^3)
\\
+\sum_j\varepsilon_j\,\big(\abs{\sigma_j}^2e^{-\varphi_j}\big)^{-1+1/\rho_j}\,
\left|\partial\sigma_j(u)-\sigma_j\partial\varphi_j(u)\right|^2 e^{-\varphi_j},
\end{multline*}
since $\dbar\sigma_j=0$. In order to simplify the calculation, we set formally
\[
\left\{
\begin{aligned}
\tilde\sigma_j&=\sigma_j^{1/\rho_j}, &
\tilde\varepsilon_j&=\rho_j^2\varepsilon_j, &
\tilde\varphi_j&=\rho_j^{-1}\varphi_j, &&
\text{if }\rho_j<\infty, \\
\tilde\sigma_j&=\log\sigma_j, &
\tilde\varepsilon_j&=\varepsilon_j, &
\tilde\varphi_j&=\varphi_j,\hfill &&
\text{if }\rho_j=\infty.
\end{aligned}
\right.
\]
Respectively to the non-logarithmic and logarithmic situations, we then get the more tractable expression
\begin{multline}
\Abs{u}_{h_{V\langle D\rangle,\varepsilon}}^2 \\
=\abs{u}^2+{}\kern-4pt
\sum_{\ell,m,\lambda,\mu}\kern-4pt c_{\ell m\lambda\mu}\,z_\ell\bar z_m u_\lambda\bar u_\mu\,{+}\,O(\abs{z}^3)
\,{+}\sum_j\tilde\varepsilon_j\,
\left|\partial\tilde\sigma_j(u)\,{-}\,\sigma_j\partial\tilde\varphi_j(u)\right|^2 e^{-\tilde\varphi_j},
\end{multline}
\begin{multline}
\tag{$\theequation_{\infty}$}
\Abs{u}_{h_{V\langle D\rangle,\varepsilon}}^2 \\
=
\abs{u}^2+{}\kern-4pt
\sum_{\ell,m,\lambda,\mu}\kern-4pt c_{\ell m\lambda\mu}\,z_\ell\bar z_m u_\lambda\bar u_\mu\,{+}\,O(\abs{z}^3)
\,{+}\sum_j\tilde\varepsilon_j\,
\left|\partial\tilde\sigma_j(u)-\partial\tilde\varphi_j(u)\right|^2.
\end{multline}
More importantly, the poles have disappeared -- a fact reflecting the orbifold smoothness of the metric. In what follows, for the sake of simplicity, we remove the tildes in the notation, and conduct the calculation only in the non-logarithmic situation $(\rho_j<\infty)$, since the logarithmic case can be recovered by taking $\rho_j$ very large; this actually amounts to using a ramified change of variable $\tilde z'_\ell=\smash{z_\ell^{1/\rho_\ell}}$ in suitable coordinates, allowing us in this way to take $\rho_j=1$ in~\eqref{6.4_0}. Also, our later calculations will be done by adding the orbifold divisor components one by one. This essentially reduces the situation to the case where $D=(1-\tfrac{1}{\rho})\Delta$ only has one component, and the notation becomes much lighter. Therefore, we drop the indices $j$ and the summations $\sum_j$, and consider the simple situation where the metric is given by
\begin{align}\label{6.8}
\Abs{u}_{h_{V\langle D\rangle,\varepsilon}}^2 &=
\abs{u}^2+\sum_{\ell,m,\lambda,\mu}c_{\ell m\lambda\mu}\,z_\ell\bar z_m u_\lambda\bar u_\mu+O(\abs{z}^3)+\varepsilon\,|\partial\sigma(u) -\sigma\,\partial\varphi(u)|^2
\\
\tag{\theequation\ensuremath{'}}\label{6.8'}
\langle\!\langle u,v\rangle\!\rangle_{h_{V\langle D\rangle,\varepsilon}}^2 &=
\langle u,v\rangle^2+
\sum_{\ell,m,\lambda,\mu}c_{\ell m\lambda\mu}\,z_\ell\bar z_m u_\lambda\bar v_\mu+O(\abs{z}^3)
\\
\nonumber&
\qquad+\varepsilon\,\big(\partial\sigma(u)-\sigma\,\partial\varphi(u)\big)
\big(\,\overline{\partial\sigma(v)-\sigma\,\partial\varphi(v)}\,\big)
\,e^{-\varphi}.
\end{align}
We also take a holomorphic trivialization of the line bundle $\cO_X(\Delta)$ so that the associated weight $\varphi$ satisfies
$\varphi(z)=\sum_{\ell,m}\alpha_{\ell m}\,z_\ell\bar z_m+O(\abs{z}^3)$ near
$z_0=0$. Then
\[
\partial\varphi=\sum_{\ell,m}\alpha_{\ell m}\,\bar z_m dz_\ell+O(\abs{z}^2),
\quad
\dbar\varphi=\sum_{\ell,m}\alpha_{\ell m}\,z_\ell\,d\bar z_m+O(\abs{z}^2).
\]
At the point $z=z_0$, we have $\partial\varphi(z_0)=\partial\varphi(z_0)=0$,
$\nabla\sigma=\partial\sigma$, and our metric admits the expression
\[
\Abs{u}_{h_{V\langle D\rangle,\varepsilon}}^2=
\abs{u}^2+\varepsilon\,\abs{\partial\sigma(u)}^2,\quad
\langle\!\langle u,v\rangle\!\rangle_{h_{V\langle D\rangle,\varepsilon}}=
\langle u,v\rangle+\varepsilon\,\partial\sigma(u)\,
\overline{\partial\sigma(v)}.
\]
Let $u,v$ be arbitrary local holomorphic sections of $V$, and denote by
$\nabla_\xi$ the Chern covariant differentiation of $(V\langle D\rangle,h_{V\langle D\rangle,\varepsilon})$ in the direction
$\xi\in T_X$. By polarizing the quadratic form
$\Abs{u}_{h_{V\langle D\rangle,\varepsilon}}^2$ into a hermitian inner product
$\partial_\xi\langle\!\langle u,v
\rangle\!\rangle_{h_{V\langle D\rangle,\varepsilon}}$ and setting $\nabla_\xi u=
\nabla^{1,0}_\xi u=\partial_\xi u+\Gamma(\xi)\cdot u$, a differentiation of~\eqref{6.8'} at $z=z_0$ yields
\begin{align*}
\partial_\xi\langle\!\langle u,v
\rangle\!\rangle_{h_{V\langle D\rangle,\varepsilon}} &=
\langle\nabla_\xi u,v\rangle+\varepsilon\,
\partial\sigma(\nabla_\xi u)\,\overline{\partial\sigma(v)}
\\
&=
\langle\partial_\xi u\,,\,v\rangle+\varepsilon\;
\partial\sigma(\partial_\xi u)\;\overline{\partial\sigma(v)}
\\
&\phantom{=}\hspace{3\parindent} +\varepsilon\,\partial^2\sigma(\xi,u)\,\overline{\partial\sigma(v)} -\varepsilon\,
\partial\sigma(u)\,\bar \sigma\,\ddbar\varphi(\xi,v),
\end{align*}
where $\partial^2\sigma(\xi,u)\bydef\sum_\lambda\partial_\xi(\partial\sigma (e_\lambda)\,u_\lambda$ is viewed as an element of
$(T_X^*\otimes V^*)_{z_0}$ and $\ddbar\varphi$ as a hermitian form on $T_X$, operating on $T_X\otimes\bar V\subset T_X\otimes \bar T_X$. In fact,
$u\mapsto\partial\sigma(u)$ and $(\xi,u)\mapsto\partial^2\sigma(\xi,u)$ can be intrinsically defined as $\nabla^{1,0}\sigma_{|V}$ and
$\smash{\nabla_{V^*\otimes \cO(\Delta)}^{1,0}}(\nabla^{1,0}\sigma_{|V})$ at $z_0$, and we will denote them by $\nabla\sigma$ and $\nabla^2\sigma$. In this setting, a subtraction of the last two lines in our equalities shows that the $(1,0)$-form $\Gamma$ of the connection of $(V\langle D\rangle,h_{V\langle D\rangle})$ is given at $z_0$ by the formula
\begin{equation}\label{6.10}
\langle \Gamma(\xi)\cdot u,v\rangle +\varepsilon\,
\nabla\sigma(\Gamma(\xi)\cdot u)\,\overline{\nabla\sigma(v)} =\varepsilon\,
\nabla^2\sigma(\xi,u)\,\overline{\nabla\sigma(v)} -\varepsilon\,\nabla\sigma(u)\;\bar\sigma\,\ddbar\varphi(\xi,v).
\end{equation}
This equality is valid pointwise for any $u,v\in V_{z_0}$. As a consequence
\begin{equation}\label{6.11}
\Gamma(\xi)\cdot u+\varepsilon\,
\nabla\sigma(\Gamma(\xi)\cdot u)\,(\nabla\sigma)^* =\varepsilon\,
\nabla^2\sigma(\xi,u)\,(\nabla\sigma)^* -\varepsilon\,\nabla\sigma(u)\;\bar\sigma\, (\ddbar\varphi(\bullet,\xi))^*
\end{equation}

\break
\noindent where $\alpha^*\in V$ is the dual vector to a $1$-form $\alpha\in V^*$, such that $\langle\alpha^*,\cdot\rangle_{h_V}=\bar\alpha$. The special choice $v=\Gamma(\xi)\cdot u$ yields a (non-negative) real value in the left hand side of~\eqref{6.10}, and by taking the real part of the right hand side, we obtain
\stepcounter{equation}
\begin{multline}
\tag{$\theequation_{0}$}\label{6.12_0}
\abs{\Gamma(\xi)\cdot u}^2 +\varepsilon\,
\abs*{\nabla\sigma(\Gamma(\xi)\cdot u)}^2
\\
=\varepsilon\,\Re\big(
\nabla^2\sigma(\xi,u)\,
\overline{\nabla\sigma(\Gamma(\xi)\cdot u)}\,\big)-\varepsilon\,
\Re\big(
\nabla\sigma(u)\,\bar\sigma\;\ddbar\varphi(\xi,\Gamma(\xi)
\cdot u)\big).
\end{multline}
Also, by applying $\nabla\sigma$ to~\eqref{6.11}, we obtain
\begin{multline*}
\nabla\sigma(\Gamma(\xi)\cdot u)+\varepsilon\,
\nabla\sigma(\Gamma(\xi)\cdot u)\,
\langle\nabla\sigma,\nabla\sigma\rangle\\
=\varepsilon\,\nabla^2\sigma(\xi,u)\,
\left\langle\nabla\sigma,\nabla\sigma\right\rangle -\varepsilon\,\nabla\sigma(u)\;\bar\sigma\,
\left\langle\nabla\sigma,\ddbar\varphi(\bullet,\xi)\right\rangle,
\end{multline*}
hence
\begin{equation_}{1}\label{6.12_1}
\nabla\sigma(\Gamma(\xi)\cdot u)=
\tfrac{\varepsilon}{ 1+\varepsilon\,\abs{\nabla\sigma}^2}\big(
\nabla^2\sigma(\xi,u)\,\abs{\nabla\sigma}^2-\nabla\sigma(u)\;\bar\sigma\,
\left\langle\nabla\sigma,\ddbar\varphi(\bullet,\xi)\right\rangle\big).
\end{equation_}
As $2\pi\,\Theta_{A,h_A}=\imath\ddbar\psi_\delta$, we infer by a brute force calculation from~\eqref{6.8} that
\begin{align}
\nonumber\stepcounter{equation} Q_{V\langle D\rangle,\gamma,\varepsilon}(z,u)(\xi,\eta) &=
\ddbar\Abs{u}_{h_{V\langle D\rangle,\varepsilon}}^2\cdot(\xi,\eta) +\gamma\,\ddbar\psi_\delta(\xi,\xi)\,
\Abs{u}_{h_{V\langle D\rangle,\varepsilon}}^2
\\
\tag{$\theequation_1$}\label{6.13_1}
&=
\gamma\,\ddbar\psi_\delta(\xi,\xi)\,\abs{u}^2 +\sum_{\ell,m,\lambda,\mu}c_{\ell m\lambda\mu}\,\xi_\ell\bar\xi_m\, u_\lambda\bar u_\mu
\\
\tag{$\theequation_2$}\label{6.13_2}
&\qquad +\varepsilon\,\big(\gamma\,\ddbar\psi_\delta(\xi,\xi)-
\ddbar\varphi(\xi,\xi)\big)\,\abs{\nabla\sigma(u)}^2\kern80pt
\\
\tag{$\theequation_3$}\label{6.13_3}
&\qquad +\abs{\eta}^2+\varepsilon\,
\abs*{\nabla\sigma(\eta)+\nabla^2\sigma(\xi,u)}^2
\\
\tag{$\theequation_4$}\label{6.13_4}
&\qquad -2\varepsilon\,\Re\big(
\nabla\sigma(u)\,\bar\sigma\,\ddbar\varphi(\xi,\eta)\big)
\\
\tag{$\theequation_5$}\label{6.13_5}
&\qquad -2\varepsilon\,\Re\big(
\nabla\sigma(u)\,\ddbar\varphi(\xi,u)\,
\overline{\nabla\sigma(\xi)}\,\big)
\\
\tag{$\theequation_6$}\label{6.13_6}
&\qquad -2\varepsilon\,\Re\big(
\nabla\sigma(u)\,\bar\sigma\,\ddbar{\kern1pt}^2
\varphi(\xi,\xi,u)\big)
\\
\tag{$\theequation_7$}\label{6.13_7}
&\qquad +\varepsilon\,\abs{\sigma}^2\,\left|\ddbar\varphi(u,\xi)\right|^2,
\end{align}
where we identify a $(1,1)$-form such as $\ddbar\varphi$ with a hermitian form, and take $\eta=-\Gamma(\xi)\cdot u$. The second term in~\eqref{6.13_2} is obtained by differentiating
$\varepsilon\,\abs{\nabla\sigma(u)}^2$, while~\eqref{6.13_3}, \eqref{6.13_4} and~\eqref{6.13_5} actually come from the differentiation of the term $\dots\Re(\dots)$ in~\eqref{6.8}. By our assumptions~\eqref{6.3_j}, the first two terms~\eqref{6.13_1}, \eqref{6.13_2} are positive in the sense of Griffiths, and such that
\begin{align*}
\eqref{6.13_1}&\geq c\,\left(\abs{\xi}^2+\abs{\nabla\sigma(\xi)}^2\right)\,\abs{u}^2, \eqref{6.13_2}&\geq c\,\varepsilon\,\left(\abs{\xi}^2+\abs{\nabla\sigma(\xi)}^2\right)\,
\abs{\nabla\sigma(u)}^2,\quad c>0.
\end{align*}
(Here the term $\abs{\nabla\sigma(\xi)}^2$ is significant, because we will later replace $\sigma$ by $\sigma^{1/\rho}$ in the orbifold case, and then $\nabla\sigma^{1/\rho}(\xi)$ is unbounded with respect to $\abs{\xi})$. The third term~\eqref{6.13_3} is semi-positive. We claim that the terms (\ref{6.13_4}, \ref{6.13_5}, \ref{6.13_6}, \ref{6.13_7} are negligible for
$\varepsilon\ll 1$, in the sense that
$Q_{V\langle D\rangle,\gamma,\varepsilon}(z,u)(\xi,\eta)$ is comprised between $(1\pm\delta)$(\eqref{6.13_1}+\eqref{6.13_2}+\eqref{6.13_3}), with $\delta>0$ as small as we want when $\varepsilon\leq\varepsilon_0(\delta)$. In fact, since $\ddbar\varphi$ is smooth, there exists $C>0$ such that
\begin{align*}
\abs*{\eqref{6.13_4}} &\leq C\,\varepsilon\abs{\sigma}\,\abs{\nabla\sigma(u)}\,\abs{\xi}\,\abs{\eta}
\\
&\leq
\varepsilon^{3/2}\abs{\xi}^2
\abs{\nabla\sigma(u)}^2 + C^2\,\varepsilon^{1/2}
\abs{\sigma}^2
\abs{\eta}^2
\ll~\eqref{6.13_2}+\eqref{6.13_3}.
\end{align*}
Similarly
\begin{align*}
\abs*{\eqref{6.13_5}} &\leq C\,\varepsilon\, \abs{\xi}\,\abs{u}\,\abs{\nabla\sigma(\xi)}\,\abs{\nabla\sigma(u)}
\\
&\leq C\,\varepsilon^{3/2}\,\abs{\xi}^2\,\abs{\nabla\sigma(u)}^2 +C\,\varepsilon^{1/2}\,\abs{\nabla\sigma(\xi)}^2\,\abs{u}^2
\ll~\eqref{6.13_1}+\eqref{6.13_2}.
\end{align*}
The last two terms (\eqref{6.13_6}, \eqref{6.13_7}) are even easier, since
\begin{align*}
\abs*{\eqref{6.13_6}} &\leq C\,\varepsilon\,\abs{\sigma}\,\abs{\xi}^2\,\abs{u}\,\abs{\nabla\sigma(u)}
\leq \varepsilon^{1/2}\abs{\xi}^2\,\abs{u}^2+ C^2\,\varepsilon^{3/2}\,\abs{\sigma}^2\,\abs{\xi}^2\,\abs{\nabla\sigma(u)}^2
\\
&\ll~\eqref{6.13_1} +~\eqref{6.13_2},
\\[4pt]
\abs*{\eqref{6.13_7}} &\leq C\,\varepsilon\,\abs{\xi}^2\abs{u}^2
\ll~\eqref{6.13_1}.
\end{align*}
Finally, by replacing $\eta$ with $-\Gamma(\xi)\cdot u$ and using (\eqref{6.12_0}, \eqref{6.12_1}), we find
\begin{align}
\nonumber\stepcounter{equation}~\eqref{6.13_3}+\eqref{6.13_4} &=
\abs*{\Gamma(\xi)\cdot u}^2\\ &\quad{}+\varepsilon\,
\big|\nabla\sigma(\Gamma(\xi)\,{\cdot}\,u) -\nabla^2\sigma(\xi,u)\big|^2 +2\varepsilon\,\Re\big(
\nabla\sigma(u)\,\bar\sigma\,
\ddbar\varphi(\xi,\Gamma(\xi)\cdot u)\big)\\ &=\eqref{6.12_0}+\varepsilon\,\big|
\nabla^2\sigma(\xi,u)\big|^2-2\varepsilon\,\Re\big(
\,\overline{\nabla^2\sigma(\xi,u)}\,
\nabla\sigma(\Gamma(\xi)\cdot u)\big)\\ &\quad{}+2\varepsilon\,\Re\big(
\nabla\sigma(u)\,\bar\sigma\,
\ddbar\varphi(\xi,\Gamma(\xi)\cdot u)\big)\\ &=
\varepsilon
\abs*{\nabla^2\sigma(\xi,u)}^2 -\varepsilon\,\Re\big(\,
\overline{\nabla^2\sigma(\xi,u)}\,
\nabla\sigma(\Gamma(\xi)\cdot u)\big)\\ &\quad{}+\varepsilon\,\Re\big(
\nabla\sigma(u)\,\bar\sigma\,
\ddbar\varphi(\xi,\Gamma(\xi)\cdot u)\big).\\
\tag{$\theequation_1$}\label{6.14_1}
&=\tfrac{\varepsilon}{ 1+\varepsilon\abs{\nabla\sigma}^2}\,
\abs*{\nabla^2\sigma(\xi,u)}^2
\\
\tag{$\theequation_2$}\label{6.14_2}
&\quad{}+\tfrac{\varepsilon}{ 1+\varepsilon\,\abs{\nabla\sigma}^2}\,\Re\Big(
\overline{\nabla^2\sigma(\xi,u)}\,
\varepsilon\,\nabla\sigma(u)\,\bar\sigma\,
\left\langle\nabla\sigma,\ddbar\varphi(\bullet,\xi)\right\rangle\Big)
\\
\tag{$\theequation_3$}\label{6.14_3}
&\quad{}+\varepsilon\,\Re\big(
\nabla\sigma(u)\,\bar\sigma\,
\ddbar\varphi(\xi,\Gamma(\xi)\cdot u)\big).
\end{align}
The term~\eqref{6.14_3} equals $\tfrac{1}{2}\eqref{6.13_4}$, thus it is negligible, and the term~\eqref{6.14_2} admits an obvious bound
\begin{align*}
~\eqref{6.14_2} &\leq
\tfrac{\varepsilon}{ 1+\varepsilon\abs{\nabla\sigma}^2}
\Big(\varepsilon^{1/2}\,\left|\nabla^2\sigma(\xi,u)\right|^2+\varepsilon^{3/2}\,
\abs{\sigma}^2\,\abs{\nabla\sigma}^2\,\abs{\nabla\sigma(u)}^2\,\abs{\xi}^2\Big)\\
&\leq
\varepsilon^{1/2}\eqref{6.14_1}+
\varepsilon^{3/2}\abs{\sigma}^2\,\abs{\nabla\sigma(u)}^2\,\abs{\xi}^2\ll~\eqref{6.14_1}+\eqref{6.13_2}.
\end{align*}
By collecting all non-negligible terms~\eqref{6.13_1}, \eqref{6.13_2} and~\eqref{6.14_1}, we obtain a curvature form
\begin{multline*}
Q_{V\langle D\rangle,\gamma,\varepsilon}(z)(\xi\otimes u)
\simeq{}\gamma\,\ddbar\psi_\delta(\xi,\xi)\,\abs{u}^2 +\sum_{\ell,m,\lambda,\mu}c_{\ell m\lambda\mu}\,\xi_\ell\bar\xi_m\, u_\lambda\bar u_\mu\\
+\varepsilon\,
\big(\gamma\,\ddbar\psi_\delta(\xi,\xi)-\ddbar\varphi(\xi,\xi)\big)
\,\abs{\nabla\sigma(u)}^2 +\tfrac{\varepsilon}{ 1+\varepsilon\abs{\nabla\sigma}^2}\,
\abs*{\nabla^2\sigma(\xi,u)}^2.
\end{multline*}
At this point, we come back to the orbifold situation, and thus replace $\sigma$ by $\sigma^{1/\rho}$,
$\varphi$ by $\rho^{-1}\varphi$ and $\varepsilon$ by
$\rho^2\,\varepsilon$. This gives the curvature estimate
\begin{multline}
Q_{V\langle D\rangle,\gamma,\varepsilon}(z)(\xi\otimes u)\simeq{}
\gamma\,\ddbar\psi_\delta(\xi,\xi)\,\abs{u}^2 +\sum_{\ell,m,\lambda,\mu}c_{\ell m\lambda\mu}\,\xi_\ell\bar\xi_m\, u_\lambda\bar u_\mu\\
\begin{aligned}
&+\varepsilon\,\abs{\sigma}^{-2+2/\rho}\,
\big(\gamma\,\ddbar\psi_\delta(\xi,\xi)-\rho^{-1}\,\ddbar\varphi(\xi,\xi)\big)
\,\abs{\nabla\sigma(u)}^2
\\
&+\tfrac{\varepsilon\,\abs{\sigma}^{-2+2/\rho}}{ 1+\varepsilon\,\abs{\sigma}^{-2+2/\rho}\,
\abs{\nabla\sigma}^2}\,\big|\,\nabla^2\sigma(\xi,u) -(1-1/\rho)\,\sigma^{-1}\,\nabla\sigma(\xi)\nabla\sigma(u)\big|^2,
\end{aligned}
\end{multline}
In the general situation $D=\sum_{1\,\leq\,j\,\leq\,N}(1-1/\rho_j)\Delta_j$ of a multi-component orbifold divisor, we add the components $\Delta_j$ one by one, and obtain inductively the following quantitative estimate, which is a rephrasing of Theorem~\ref{theo:orbi_curv_estimate}.
\end{proof}

\begin{coro}\label{coro:curv_estimate}
With a choice of
$\gamma>\gamma_{V,D}\bydef\max(\max(d_j/\rho_j),\gamma_V)\geq 0$ determined by the curvature assumptions of Proposition~\ref{prop:orbi_metric}, and of hermitian metrics on $A$, $V$, $\cO_X(D)$ as prescribed by conditions~\eqref{6.3_j}, the orbifold metric
\begin{flalign}\label{coro7.2.a}
\tag{a}
\abs{u}_{h_{V\langle D\rangle,\varepsilon}}^2
\bydef
\abs{u}^2_{h_V}+\sum_{1\,\leq\,j\,\leq\,N}\varepsilon_j\,\abs{\sigma_j}^{-2+2/\rho_j}\,
\left|\nabla_j\sigma_j(u)\right|^2_{h_j} &&
\end{flalign}

yields a curvature tensor
$\theta_{V\langle D\rangle,\gamma,\varepsilon}
\bydef
\gamma\,\Theta_{A,h_{A,\delta}}\otimes\Id-
\Theta_{V\langle D\rangle,h_{V\langle D\rangle,\varepsilon}}$ such that the associated quadratic form $Q_{V\langle D\rangle,\gamma,\varepsilon}$ on
$T_X\otimes V$ satisfies for
$\varepsilon_N\ll\varepsilon_{N-1}\ll\cdots\ll\varepsilon_1\ll 1$ the curvature estimate
\begin{multline}\label{coro7.2.b}
\tag{b} Q_{V\langle D\rangle,\gamma,\varepsilon} (z)(\xi\otimes u)\simeq\gamma\,\ddbar\psi_\delta(\xi,\xi)\,\abs{u}^2 +\sum_{\ell,m,\lambda,\mu}c_{\ell m\lambda\mu}\,\xi_\ell\bar\xi_m\, u_\lambda\bar u_\mu
\\
+\sum_j\varepsilon_j\,\abs{\sigma_j}^{-2+2/\rho_j}\,
\big(\gamma\,\ddbar\psi_\delta(\xi,\xi)-\rho_j^{-1}\,\ddbar\varphi_j(\xi,\xi)\big)
\,\abs{\nabla_j\sigma_j(u)}^2\\ +\sum_j\tfrac{\varepsilon_j\,\abs{\sigma_j}^{-2+2/\rho_j}}{ 1+\varepsilon_j\,\abs{\sigma_j}^{-2+2/\rho_j}\,
\abs{\nabla_j\sigma_j}^2}\,\left|\,\nabla_j^2\sigma_j(\xi,u)-(1-1/\rho_j)\,
\sigma_j^{-1}\,\nabla_j\sigma_j(\xi)\nabla_j\sigma_j(u)\right|^2,
\end{multline}
where
\[
\nabla_{A,h_{A,\delta}}^2=\ddbar\psi_\delta,\quad
\nabla_{\Delta_j,h_j}^2=\ddbar\varphi_j,\quad
\text{$(c_{\ell m\lambda\mu})={}$coefficients of $-2\pi\,\Theta_{V,h_V}$}.
\]
Here, the symbol $\simeq$ means that the ratio of the left and right hand sides can be chosen in $[1-\alpha,1+\alpha]$ for any
$\alpha>0$ prescribed in advance.
\end{coro}

\subsection{Evaluation of some Chern form integrals and their limits}\label{sse:Chern_integrals}
Our aim is to apply Lemma~\ref{lemm:Chern_form_inequalities} and Corollary~\ref{coro:curv_estimate} to compute Morse integrals of the curvature tensor of a directed orbifold $(X,V,D)$, where
$D=\sum_j(1-1/\rho_j)\Delta_j$ is transverse to $V$. Let $A\in \Pic(X)$ be an ample line bundle, and $d_j$, $\gamma_V$,
$\gamma>\gamma_{V,D}$ be defined as in Corollary~\ref{coro:curv_estimate}. We get hermitian metrics
$h_{V\langle D\rangle,\varepsilon}$ on $V\langle D\rangle$ and corresponding curvature tensors
$\theta_{V\langle D\rangle,\gamma,\varepsilon}$ in $C^\infty(X\setminus\abs{D},
\Lambda^{1,1}T^*_X\otimes\Hom(V,V))$ that are ``orbifold smooth'', and such that $\theta_{V\langle D\rangle,\gamma,\varepsilon}\geq_G 0$. Given a smooth strongly positive $(n-p,n-p)$-form $\beta\geq_S 0$ on $X$, we want to evaluate the integrals

\begin{align}\label{6.17}
I_{p,\varepsilon}(\beta)&=\int_{S_\varepsilon(V\langle D\rangle)}
\left\langle\theta_{V\langle D\rangle,\gamma,\varepsilon}\cdot u,u\right\rangle^p
\wedge\beta\,d\mu_\varepsilon(u)
\\
\tag{\theequation\ensuremath{'}} &=\int_{z\,\in\,X}\int_{u\,\in\,S_\varepsilon(V\langle D\rangle)_z}
\left\langle\theta_{V\langle D\rangle,\gamma,\varepsilon}\cdot u,u\right\rangle^p
\wedge\beta(z)\,d\mu_\varepsilon(u),
\end{align}
where $S_\varepsilon(V\langle D\rangle)$ denotes the unit sphere bundle of $V\langle D\rangle$ with respect to $h_\varepsilon$, and
$\mu_\varepsilon$ the unitary invariant probability measure on the sphere. Proposition~\ref{prop:Chern_curv_inequalities}$\MK$\eqref{item:Chern_curv_inequality2} and the Fubini theorem imply the upper bound
\[
I_{p,\varepsilon}(\beta)\leq \tfrac{p!\,(r-1)!}{(p+r-1)!}
\int_X \left(\Tr\theta_{V\langle D\rangle,\gamma,\varepsilon}\right)^p\wedge \beta.
\]
When $\beta$ is closed, the upper bound can be evaluated by a cohomology class calculation, thanks to the following lemma.

\begin{lemm}
The $(1,1)$-form
$\Tr\theta_{V\langle D\rangle,\gamma,\varepsilon}\geq 0$ is closed and belongs to the cohomology class
\[
r\gamma\,c_1(A)-c_1(V)+\sum_j\left(1-1/\rho_j\right)c_1(\Delta_j).
\]
\end{lemm}
\begin{proof}
The trace can be seen as the curvature of
\begin{align*}
\det(\cO_X(\gamma A)\otimes V\langle D\rangle^*) &=\cO_X(r\gamma\,A)\otimes \det(V\langle D\rangle^*)\\
&=\cO_X(r\gamma\,A)\otimes \det(V^*)\otimes\cO_X(D)
\end{align*}
with the determinant metric. Since all metrics have equivalent behaviour along $\abs{D}$ (and can be seen as orbifold smooth), Stokes' theorem shows that the cohomology class is independent of $\varepsilon$. Formally, the result follows from~\eqref{eq:canonical}. One can also consider the intersection product
\[
\{\Tr\theta_{V\langle D\rangle,\gamma,\varepsilon}\}\cdot \{\beta\} =\int_X\Tr\theta_{V\langle D\rangle,\gamma,\varepsilon}\wedge\beta =r\int_{u\,\in\,S_\varepsilon(V\langle D\rangle)}
\langle\theta_{V\langle D\rangle,\gamma,\varepsilon}\cdot u,u\rangle
\wedge\beta\,d\mu_\varepsilon(u)
\]
for all smooth closed
$(n-1,n-1)$-forms $\beta$ on $X$, and apply Corollary~\ref{coro:curv_estimate}~(b) to evaluate the limit as $\varepsilon\to 0$. This will be checked later as the special case $p=1$ of~\eqref{6.17}.
\end{proof}

We actually need even more general estimates. The proof follows again from the Fubini theorem.

\begin{prop}\label{6.20}
Consider orbifold directed structures
$(X,V,D_s)$, $1\leq s\leq k$, with
$D_s=\sum_{1\,\leq\,j\,\leq\,N}(1-\tfrac{1}{\rho_{s,j}})\Delta_j$. We assume that the divisors $D_s$ are simple normal crossing divisors transverse to $V$, sharing the same components $\Delta_j$. Let $d_j$ be the infimum of numbers
$\lambda\in\bR_+$ such that $\lambda\,A-\Delta_j$ is~nef, and let
$\gamma_V$ be the infimum of numbers
$\gamma\geq 0$ such that
$\theta_{V,\gamma}\bydef\gamma\,\Theta_{A,h_A}\otimes\Id_V-\Theta_{V,h_V}\geq_G0$ for suitable hermitian metrics $h_V$ on $V$. Take $p=(p_1,\,\dotsc,\,p_k)\in\bN^k$ such that $p'=n-(p_1+\dotsb+p_k)\geq 0$ and a smooth, closed, strongly positive $(p',p')$ form $\beta\geq_S0$ on $X$. Then for every
\[
\gamma_s>\gamma_{V,D_s}\bydef\max\left(\max_j(d_j/\rho_{s,j}),\gamma_V\right)
\]
there exist hermitian metrics
$h_{V\langle D_s\rangle,\varepsilon_s}$ on the orbifold vector bundles
$V\langle D_s\rangle$ such that
\begin{multline*}
\theta_{V\langle D_s\rangle,\gamma_s,\varepsilon_s}\bydef\gamma_s
\,\Theta_{A,h_A}\otimes\Id_V-\Theta_{V\langle D_s\rangle, h_{V\langle D_s\rangle,\varepsilon_s}}>_G0,\\
\varepsilon_s=(\varepsilon_{s1},\,\dotsc,\,\varepsilon_{sN}),\quad 1\leq s\leq k
\end{multline*}
in the sense of Griffiths. Moreover, the integrals
\begin{multline}
\tag{a} I_{k,p,\varepsilon}(\beta)\\
=
\int_{z\,\in\,X}\int_{(u_s)\,\in\,\Pi_s S(V\langle D_s\rangle)_z}~\bigwedge_{1\,\leq\,s\,\leq\,k}
\left\langle\theta_{V\langle D_s\rangle,\gamma_s,\varepsilon_s}(u_s),u_s\right\rangle^{p_s}
\wedge\beta(z)\,\prod_{1\,\leq\,s\,\leq\,k}d\mu_{\varepsilon_s}(u_s)
\end{multline}
admit upper bounds
\begin{multline}
\tag{b} I_{k,p,\varepsilon}(\beta)\\
\le
\int_X\bigwedge_{1\,\leq\,s\,\leq\,k}\tfrac{p_s!\,(r-1)!}{(p_s+r-1)!}
\left(r\gamma_s\,\Theta_{A,h_{A,\delta}}-\Tr\Theta_{V,h_V}+
\sum_j(1-1/\rho_{s,j})\Theta_{\Delta_j,h_j}\right)^{p_s}\!\wedge\beta.
\end{multline}
When $\beta$ is closed, we get a purely cohomological upper bound
\begin{multline}
\tag{c} I_{k,p,\varepsilon}(\beta)\\
\le
\int_X\prod_{1\,\leq\,s\,\leq\,k}\tfrac{p_s!\,(r-1)!}{(p_s+r-1)!}\,
\left(r\gamma_s\,c_1(A)-c_1(V)+\sum_j(1-1/\rho_{s,j})c_1(\Delta_j)\right)^{p_s}
\cdot\{\beta\}.
\end{multline}
\end{prop}

\begin{enonce}{Complement}\label{6.21}
When $p_1=\ldots=p_k=1$, formulas~\ref{6.20}~(b) and~\ref{6.20}~(c) are equalities.
\end{enonce}
\begin{proof}
This follows from Remark~\ref{5.15}.
\end{proof}

In general, getting a lower bound for $I_{p,\varepsilon}(\beta)$ and
$I_{k,p,\varepsilon}(\beta)$ is substantially harder. We start with
$I_{p,\varepsilon}(\beta)$ and content ourselves to evaluate the iterated limit
\[
\lim_{\varepsilon\,\to\,0}I_{p,\varepsilon}(\beta)
\bydef
\lim_{\varepsilon_1\,\to\,0}\lim_{\varepsilon_2\to 0}\ldots\lim_{\varepsilon_N\,\to\,0} I_{p,\varepsilon}(\beta),
\quad~~
\varepsilon_N
\ll
\varepsilon_{N-1}
\ll \dotsb \ll
\varepsilon_1
\ll 1.
\]
For this, we consider the expression of the curvature form in a neigborhood of an arbitrary point
$z_0\in\Delta_{j_1}\cap\ldots\cap\Delta_{j_m}$ (if $z_0\in X\setminus\abs{\Delta}$, we have $m=0$). We take trivializations of the line bundles
$\cO_X(\Delta_j)$ so that the hermitian metrics have weights
$e^{-\varphi_j}$ with $\varphi_j(z_0)=d\varphi_j(z_0)=0$, and introduce the corresponding ``orbifold'' coordinates
\[
t_{j,\varepsilon}=\varepsilon_j^{1/2}\,\sigma_j(z)^{-(1-1/\rho_j)}\,\abs{\nabla_j\sigma_j(z_0)},
\quad j=j_1,\dotsc,j_m,
\]
We complete these coordinates with $n-m$ variables $z_\ell$ that define coordinates along $\Delta_{j_1}\cap\ldots\cap\Delta_{j_m}$. In this way, we get a $n$-tuple $(t_{j,\varepsilon},z_\ell)$ of complex numbers that provide local coordinates on the universal cover of $\Omega_{z_0}\setminus\abs{D}$, where $\Omega_{z_0}$ is a small neighborhood of $z_0$. Viewed on $X$, the coordinates $t_{j,\varepsilon}$ are multivalued near $z_0$, but we can make a ``cut'' in $X$ along $\Delta_j$ to exclude the negligible set of points where
$\sigma_j(z)\in\bR_-$, and take the argument in $]-\pi,\pi[$, so that
$\Arg(t_{j,\varepsilon})\in{}]-(1-1/\rho_j)\pi,(1-1/\rho_j)\pi[\,$. If we integrate over complex numbers $t_{j,\varepsilon}$ without such a restriction on the argument, the integral will have to be multiplied by the factor $(1-1/\rho_j)$ to get the correct value. Since $\abs{\sigma_j}$ is bounded, the range of the absolute value $\abs{t_{j,\varepsilon}}$ is an interval
$]O(\varepsilon_j^{1/2}),+\infty[\,$, thus $t_{j,\varepsilon}$ will cover asymptotically an entire angular sector in $\bC$ as $\varepsilon_j\to 0$. In the above coordinates, we have
\[
\tfrac{dt_{j,\varepsilon}}{ t_{j,\varepsilon}}=-(1-1/\rho_j)\tfrac{d\sigma_j}{\sigma_j}= -(1-1/\rho_j)\bigg(\tfrac{\nabla_j\sigma_j}{\sigma_j}+\partial\varphi_j\bigg)= -(1-1/\rho_j)\tfrac{\nabla_j\sigma_j}{\sigma_j}+O(1),
\]
since $\nabla_j\sigma_j=d\sigma_j-\sigma_j\partial\varphi_j$ and the weight
$\varphi_j$ of the metric of $\cO_X(\Delta_j)$ is smooth. Denote
\begin{align}
\stepcounter{equation}
\tag{$\theequation_1$}
\theta_{V,\gamma}&=
\gamma\,\Theta_{A,h_{A,\delta}}\otimes\Id_V-\Theta_{V,h_V},
\\ \tag{$\theequation_2$}
\theta_{V\langle D\rangle,\gamma,\varepsilon}&\bydef
\gamma\,\Theta_{A,h_{A,\delta}}\otimes\Id-
\Theta_{V\langle D\rangle,h_{V\langle D\rangle,\varepsilon}},
\\ \tag{$\theequation_3$} e_j^*&=
\tfrac{\nabla_j\sigma_j}{\abs{\nabla_j\sigma_j}}\in S(V^*).
\end{align}
By Corollary~\ref{coro:curv_estimate}, we have
\begin{multline}\label{6.26}
\langle\theta_{V\langle D\rangle,\gamma,\varepsilon}\cdot u,u\rangle
\simeq\gamma\,\Theta_{A,h_{A,\delta}}-\langle\Theta_{V,h_V}\cdot u,u\rangle
\\
\begin{aligned}
&+\sum_j\varepsilon_j\,\abs{\sigma_j}^{-2+2/\rho_j}\,
\big(\gamma\,\Theta_{A,h_{A,\delta}}-\rho_j^{-1}\,\Theta_{\Delta_j,h_j}\big)
\,\left|\nabla_j\sigma_j(u)\right|^2
\\
&+\tfrac{1}{ 2\pi}\sum_j\tfrac{\varepsilon_j\,\abs{\sigma_j}^{-2+2/\rho_j}}{ 1+\varepsilon_j\,\abs{\sigma_j}^{-2+2/\rho_j}\,
\abs{\nabla_j\sigma_j}^2}\,\left|\,\nabla_j^2\sigma_j(\xi,u)-(1-1/\rho_j)\,
\sigma_j^{-1}\,\nabla_j\sigma_j(\xi)\nabla_j\sigma_j(u)\right|^2,
\end{aligned}
\end{multline}
therefore
\stepcounter{equation}
\begin{align}
\tag{$\theequation_1$} \label{6.27_1}
\langle\theta_{V\langle D\rangle,\gamma,\varepsilon}\cdot u,u\rangle &\simeq
\langle\theta_{V,\gamma}\cdot u,u\rangle+\sum_j
\big(\gamma\,\Theta_{A,h_{A,\delta}}-\rho_j^{-1}\,\Theta_{\Delta_j,h_j}\big)
\,\abs{t_{j,\varepsilon}}^2\,\abs{e_j^*(u)}^2
\\ \tag{$\theequation_2$} \label{6.27_2}
&\kern10pt{} +\tfrac{i}{ 2\pi}
\sum_j\tfrac{\abs{t_{j,\varepsilon}}^2}{ 1+\abs{t_{j,\varepsilon}}^2}\,\Big\langle
\tfrac{dt_{j,\varepsilon}}{ t_{j,\varepsilon}}\,e_j^*(u)+b_j(u)\,,\,
\tfrac{dt_{j,\varepsilon}}{ t_{j,\varepsilon}}\,e_j^*(u)+b_j(u)\Big\rangle_{h_j},
\intertext{where}
\tag{$\theequation_3$} b_j&=\tfrac{1}{\abs{\nabla_j}}\nabla^2\sigma_j\in C^\infty\left(\Omega_{z_0},\Lambda^{1,0}T^*_X\otimes V^*\otimes\cO_X(\Delta_j)\right)
\end{align}
is a smooth $(1,0)$-form near $z_0$. The approximate equality $\simeq$ in formula (\eqref{6.27_1}, \eqref{6.27_2}) involves the approximation $\abs{\nabla_j\sigma_j(z)}/\abs{\nabla_j\sigma_j(z_0)}\simeq 1$, which holds in a sufficiently small neighborhood of $z_0\,$; if we apply the Fubini theorem and consider the fiber integral over $z_0\in X$, there is actually no error coming from this approximation. Now, we want to integrate the volume form
$\langle\theta_{V\langle D\rangle,\gamma,\varepsilon}\cdot u,u\rangle^p \wedge\beta\,d\mu_\varepsilon(u)$ along the fibers of
$S_\varepsilon(V\langle D\rangle)\to X$. The sphere bundle $S_\varepsilon(V\langle D\rangle)$ is defined by
$\abs{u}_{h_{V\langle D\rangle,\varepsilon}}^2=1$ where
\begin{equation}\label{6.28}
\abs{u}_{h_{V\langle D\rangle,\varepsilon}}^2=
\abs{u}^2+\sum_j\varepsilon_j\abs{\sigma_j}^{-2+2/\rho_j}\,\left|\nabla_j\sigma_j(u)\right|^2
\simeq
\abs{u}^2+\sum_j\left|t_{j,\varepsilon}\right|^2\left|e_j^*(u)\right|^2=1.
\end{equation}
For the sake of simplicity, we first deal with the case where the divisor $D=(1-1/\rho_j)\Delta_j$ has a single component. Along
$\Delta_j$, we then get an orthogonal decomposition
$V=(V\cap T_{\Delta_j})\oplus\bC e_j$, and by~\eqref{6.28} we can write
\[
\tag{$\theequation^{\perp}$} u=u'_j+e_j^*(u)\,e_j\in S(V),\quad
\abs{u}^2=\abs{u'_j}^2+\left|e_j^*(u)\right|^2,\quad u'_j\in V\cap T_{\Delta_j}.
\]
We reparametrize the integration in $u\in S_\varepsilon(V\langle D\rangle)$ on the sphere $S(V)$ by introducing the change of variables
\begin{align*}
&\tau=\tau_{j,\varepsilon}=\tfrac{\abs{t_{j,\varepsilon}}^2}{ 1+\abs{t_{j,\varepsilon}}^2}
\in[0,1],\quad 1-\tau=\tfrac{1}{ 1+\abs{t_{j,\varepsilon}}^2},\quad d\tau=\tfrac{d\abs{t_{j,\varepsilon}}^2}{\left(1+\abs{t_{j,\varepsilon}}^2\right)^2},
\\
&g_{j,\varepsilon}:S(V)\to S_\varepsilon(V\langle D\rangle),\quad u\mapsto u_{j,\varepsilon}=u'_j+\sqrt{1-\tau}\,e_j^*(u)\,e_j =u'_j+\tfrac{e_j^*(u)}{\left(1+\abs{t_{j,\varepsilon}}^2\right)^{1/2}}\,e_j,
\end{align*}
so that $u_{j,\varepsilon}$ satisfies
$\abs{u_{j,\varepsilon}}^2_\varepsilon=\abs{u}^2$ and
\[
e_j^*(u_{j,\varepsilon})=\sqrt{1-\tau}\,e_j^*(u) =\tfrac{1}{\left(1+\abs{t_{j,\varepsilon}}^2\right)^{1/2}}\,e_j^*(u),
\quad
\abs{t_{j,\varepsilon}}^2\,\left|e_j^*(u_{j,\varepsilon})\right|^2=\tau\,\left|e_j^*(u)\right|^2.
\]
This gives $d\mu_\varepsilon(u_{j,\varepsilon})=d\mu(u)$, and as a consequence~\eqref{6.17} can be rewritten as
\begin{equation}\label{6.29}
I_{p,\varepsilon}(\beta)=\int_{S(V)}\left\langle\theta_{V\langle D\rangle,\gamma,\varepsilon}\cdot u_{j,\varepsilon},u_{j,\varepsilon}\right\rangle^p\wedge\beta(z)\,d\mu(u).
\end{equation}
Finally, a use of polar coordinates with
$\alpha=\Arg(t_{j,\varepsilon})$ shows that
\[
\tfrac{\imath dt_{j,\varepsilon}\wedge d\bar t_{j,\varepsilon}}{\left(1+\abs{t_{j,\varepsilon}}^2\right)^2}=
\tfrac{2\,\abs{t_{j,\varepsilon}}\,d\abs{t_{j,\varepsilon}}\wedge d\alpha}{\left(1+\abs{t_{j,\varepsilon}}^2\right)^2} = d\tau\wedge d\alpha.
\]
A substitution $u\mapsto u_{j,\varepsilon}$ in (\eqref{6.27_1}, \eqref{6.27_2}) yields
\begin{multline}\label{6.30}
\langle\theta_{V\langle D\rangle,\gamma,\varepsilon}\cdot u_{j,\varepsilon}, u_{j,\varepsilon}\rangle
\simeq\langle\theta_{V,\gamma}\cdot u_{j,\varepsilon},u_{j,\varepsilon}
\rangle+
\left(\gamma\,\Theta_{A,h_{A,\delta}}-\rho_j^{-1}\,\Theta_{\Delta_j,h_j}\right)
\,\tfrac{\abs{t_{j,\varepsilon}}^2\,\abs{e_j^*(u)}^2}{ 1+\abs{t_{j,\varepsilon}}^2}\\ +\tfrac{\imath}{ 2\pi}\tfrac{1}{ 1+\abs{t_{j,\varepsilon}}^2}\,\Big\langle
\tfrac{dt_{j,\varepsilon}\,e_j^*(u)}{\left(1+\abs{t_{j,\varepsilon}}^2\right)^{1/2}}+ t_{j,\varepsilon}\,b_j\left(u_{j,\varepsilon}\right)\,,\,
\tfrac{dt_{j,\varepsilon}\,e_j^*(u)}{\left(1+\abs{t_{j,\varepsilon}}^2\right)^{1/2}}+ t_{j,\varepsilon}\,b_j\left(u_{j,\varepsilon}\right)\Big\rangle_{h_j}.
\end{multline}
The last term is a $(1,1)$-form that is a square of a $(1,0)$-form (when $u$ is fixed), hence the expansion of the $p^{\rm th}$ power can involve at most one such factor. Therefore we get
\begin{multline}
\tag{$\theequation^{p}$}
\left\langle\theta_{V\langle D\rangle,\gamma,\varepsilon}\cdot u_{j,\varepsilon}, u_{j,\varepsilon}\right\rangle^p
\simeq\\
\left(
\left\langle\theta_{V,\gamma}\cdot u_{j,\varepsilon},u_{j,\varepsilon}\right\rangle +\left(\gamma\,\Theta_{A,h_{A,\delta}}-\rho_j^{-1}\,\Theta_{\Delta_j,h_j}\right)
\,\tfrac{\left|t_{j,\varepsilon}\right|^2\,\left|e_j^*(u)\right|^2}{1+\left|t_{j,\varepsilon}\right|^2}\right)^p\\ +p\,\tfrac{\imath}{ 2\pi}\tfrac{1}{1+\left|t_{j,\varepsilon}\right|^2}\,\left\langle
\tfrac{dt_{j,\varepsilon}\,e_j^*(u)}{\left(1+\abs{t_{j,\varepsilon}}^2\right)^{1/2}}+ t_{j,\varepsilon}\,b_j\left(u_{j,\varepsilon}\right)\,,\,
\tfrac{dt_{j,\varepsilon}\,e_j^*(u)}{\left(1+\abs{t_{j,\varepsilon}}^2\right)^{1/2}}+ t_{j,\varepsilon}\,b_j\left(u_{j,\varepsilon}\right)\right\rangle\\
\wedge{}\left(
\left\langle\theta_{V,\gamma}\cdot u_{j,\varepsilon},u_{j,\varepsilon}\right\rangle +\left(\gamma\,\Theta_{A,h_{A,\delta}}-\rho_j^{-1}\,\Theta_{\Delta_j,h_j}\right)
\,\tfrac{\left|t_{j,\varepsilon}\right|^2\,\left|e_j^*(u)\right|^2}{ 1+\left|t_{j,\varepsilon}\right|^2}
\right)^{p-1}.
\end{multline}
The integrals involving $b_j(u_{j,\varepsilon})$ are of the form
\[
\int_{S(V)}\tfrac{\bar t_{j,\varepsilon}\,dt_{j,\varepsilon}\wedge
\left\langle e_j^*(u),b_j(u_{j,\varepsilon})\right\rangle }{\left(1+\abs{t_{j,\varepsilon}}^2\right)^{3/2}}\wedge A_{j,\varepsilon}(u),\quad
\int_{S(V)}\tfrac{\left|t_{j,\varepsilon}\right|^2
\left\langle b_j\left(u_{j,\varepsilon}\right),b_j\left(u_{j,\varepsilon}\right)\right\rangle}{ 1+\left|t_{j,\varepsilon}\right|^2}\wedge A'_{j,\varepsilon}(u)
\]
where $A_{j,\varepsilon}(u)$, $A'_{j,\varepsilon}(u)$ are forms with uniformly bounded coefficients in orbifold coordinates. Since $\tfrac{\abs{t_{j,\varepsilon}}^2}{ 1+\abs{t_{j,\varepsilon}}^2}$ is bounded by $1$ and converges to $0$ on $X\setminus\Delta_j$, Lebesgue's dominated convergence theorem shows that the second integral converges to $0$. The second integral can be estimated by the Cauchy--Schwarz inequality. We obtain an upper bound
\[
\left(\int_{S(V)}\tfrac{\left|t_{j,\varepsilon}\right|^2
\langle b_j(u_{j,\varepsilon}),b_j(u_{j,\varepsilon})\rangle }{ 1+\abs{t_{j,\varepsilon}}^2}\wedge A_{j,\varepsilon}(u)\right)^{1/2}
\left(\int_{S(V)}\tfrac{\imath dt_{j,\varepsilon}\wedge d\bar t_{j,\varepsilon}\,
\abs{e_j^*(u)}^2 }{\left(1+\abs{t_{j,\varepsilon}}^2\right)^2}\wedge A_{j,\varepsilon}(u)\right)^{1/2}
\]
where the first factor converges to $0$ and the second one is bounded by Fubini, since $\int_\bC \imath dt\wedge d\bar t/(1+\abs{t}^2)^2<+\infty$. Modulo negligible terms, and changing variables into our new parameters
$(\tau,\alpha)$, we finally obtain
\begin{multline}
\langle\theta_{V\langle D\rangle,\gamma,\varepsilon}\cdot u_{j,\varepsilon}, u_{j,\varepsilon}\rangle^p
\simeq
\langle\theta_{V,\gamma}\cdot u_{j,\varepsilon},u_{j,\varepsilon}\rangle^p\\
+p\,\tfrac{d\tau\wedge d\alpha}{ 2\pi}\,\abs{e_j^*(u)}^2\wedge\Big(
\langle\theta_{V,\gamma}\cdot u_{j,\varepsilon},u_{j,\varepsilon}\rangle +\big(\gamma\,\Theta_{A,h_{A,\delta}}-\rho_j^{-1}\,\Theta_{\Delta_j,h_j}\big)
\,\tau\,\abs{e_j^*(u)}^2\Big)^{p-1}.
\end{multline}
Therefore
\begin{multline}\label{6.32}
\int_{S(V)}\langle\theta_{V\langle D\rangle,\gamma,\varepsilon}\cdot u_{j,\varepsilon}, u_{j,\varepsilon}\rangle^p\wedge\beta\,d\mu(u)
\simeq \int_{S(V)}
\left\langle\theta_{V,\gamma}\cdot u_{j,\varepsilon},u_{j,\varepsilon}\right\rangle^p
\wedge\beta\,d\mu(u)
\\
+\int_{S(V)}p\,\tfrac{d\tau\wedge d\alpha}{ 2\pi}\,\left|e_j^*(u)\right|^2\wedge\left(
\left\langle\theta_{V,\gamma}\cdot u_{j,\varepsilon},u_{j,\varepsilon}\right\rangle +\left(\gamma\,\Theta_{A,h_{A,\delta}}-\rho_j^{-1}\,\Theta_{\Delta_j,h_j}\right)
\,\tau\,\left|e_j^*(u)\right|^2\right)^{p-1}
\\
\wedge\beta\,d\mu(u).
\end{multline}
Since $u_{j,\varepsilon}\to u$ almost everywhere and boundedly, we have
\[
\lim_{\varepsilon\,\to\,0}
\int_{S(V)}\left\langle\theta_{V,\gamma}\cdot u_{j,\varepsilon},u_{j,\varepsilon}\right\rangle^p
\wedge\beta\,d\mu(u)=
\int_{S(V)}
\left\langle\theta_{V,\gamma}\cdot u,u\right\rangle^p\wedge\beta\,d\mu(u).
\]
Here, we have to remember that $\tau=\tau_{j,\varepsilon}$ converges uniformly to $0$ (even in the $C^\infty$ topology), on all compact subsets of $X\setminus\Delta_j$, hence the second integral in~\eqref{6.32} asymptotically concentrates on $\Delta_j$ as $\varepsilon\to 0$. Also, the angle $\alpha=\Arg(t_{j,\varepsilon})$ runs over the interval $]-(1-1/\rho_j)\pi,(1-1/\rho_j)\pi[$. In the easy case
$p=1$, we get
\begin{multline*}
\lim_{\varepsilon\,\to\,0}
\int_{S(V)}\left\langle\theta_{V,\gamma,\varepsilon}\cdot u_{j,\varepsilon}, u_{j,\varepsilon}\right\rangle\wedge\beta\,d\mu(u)\\
\begin{aligned}
&=\int_{S(V)}\left\langle\theta_{V,\gamma}\cdot u,u\right\rangle\wedge\beta\,d\mu(u) +(1-1/\rho_j)\int_{S(V)_{\abs{\Delta_j}}}e_j^*(u)|^2\beta\,d\mu(u)\\
&=\int_{X}\tfrac{1}{r}\,\Tr\theta_{V,\gamma}\wedge\beta +(1-1/\rho_j)\int_{\Delta_j}\tfrac{1}{ r}\beta.
\end{aligned}
\end{multline*}
If we assume $\beta$ closed, this is equal to the intersection product
\[
\tfrac{1}{ r}\big(\rho\gamma\,c_1(A)-c_1(V)+(1-1/\rho_j) c_1(\Delta_j)\big)\cdot\beta
\]
and the final assertion of the proof of Lemma~6.19 is thus confirmed, adding the components $\Delta_j$ one by one (see below). Now, in the general case $p\geq 1$, we will obtain a lower bound of the second integral involving $d\tau\wedge d\alpha$ in~\eqref{6.32} by using a change of variable $h_{j,\varepsilon}\colon S(V)\to S(V)$,
\[
u\mapsto h_{j,\varepsilon}(u) =
\big((1-\tau)\abs{u'}^2+\left|e_j^*(u)\right|^2\big)^{-1/2}\,
\big(\sqrt{1-\tau}\,u'+e_j^*(u)\,e_j\big)
\]
where $\tau=\tau_{j,\varepsilon}$. Observe that the composition $g_{j,\varepsilon}\circ h_{j,\varepsilon}\colon S(V)\to S(V)\to S_\varepsilon(V\langle D\rangle)$ is given by
\[
g_{j,\varepsilon}\circ h_{j,\varepsilon}(u) =
\tfrac{\sqrt{1-\tau}}{
\big((1-\tau)\abs{u'}^2+\left|e_j^*(u)\right|^2\big)^{1/2}}\,u.
\]
Since $(1-\tau)\abs{u'}^2+\abs{e_j^*(u)}^2\leq\abs{u}^2=1$, it is easy to check that $d\mu(h_{j,\varepsilon}(u))\linebreak\geq(1-\tau)^{r-1}\,d\mu(u)$ on the unit sphere, that $\abs{e_j^*(h_{j,\varepsilon}(u))}\geq \abs{e_j^*(u)}$, and finally, that
\begin{multline*}
\left\langle\theta_{V,\gamma}\cdot g_{j,\varepsilon}\left(h_{j,\varepsilon}(u)\right),g_{j,\varepsilon}\left(h_{j,\varepsilon}(u)\right)
\right\rangle\\
=\tfrac{1-\tau}{ (1-\tau)\abs{u'}^2+\left|e_j^*(u)\right|^2}\,
\langle\theta_{V,\gamma}\cdot u,u\rangle
\geq(1-\tau)\,\langle\theta_{V,\gamma}\cdot u,u\rangle.
\end{multline*}
Hence, by a change a variable $u\mapsto h_{j,\varepsilon}(u)$ we find
\begin{multline}\label{6.33}
\int_{S(V)}p\,\tfrac{d\tau\wedge d\alpha}{ 2\pi}\,\abs{e_j^*(u)}^2\\
\wedge\Big(
\left\langle\theta_{V,\gamma}\cdot u_{j,\varepsilon},u_{j,\varepsilon}\right\rangle +\left(\gamma\,\Theta_{A,h_{A,\delta}}-\rho_j^{-1}\,\Theta_{\Delta_j,h_j}\right)
\,\tau\,\left|e_j^*(u)\right|^2\Big)^{p-1}
\wedge\beta\,d\mu(u)\\
\geq\int_{S(V)}p\,\tfrac{d\tau\wedge d\alpha}{2\pi}\,\left|e_j^*(u)\right|^2\wedge\Big((1-\tau)\left\langle\theta_{V,\gamma}\cdot u,u\right\rangle +\left(\gamma\,\Theta_{A,h_{A,\delta}}-\rho_j^{-1}\,\Theta_{\Delta_j,h_j}\right)
\,\tau\,\abs{e_j^*(u)}^2\Big)^{p-1}\\
\hfill\wedge\beta\,(1-\tau)^{r-1}\,d\mu(u).
\end{multline}
Here, we have to remember that $\tau=\tau_{j,\varepsilon}$ converges uniformly to $0$ (even in the $C^\infty$ topology), on all compact subsets of $X\setminus\Delta_j$. Therefore, the last integral concentrates over the divisor $\Delta_j$. If we apply the binomial formula with an index $q'=q-1$, we see that the limit as $\varepsilon\to 0$ is equal to
\begin{multline}\label{6.34}
p~(1-1/\rho_j)\int_{S(V)_{|\Delta_j}}\sum_{q=1}^{p}\binom{p-1}{ q-1}
\langle\theta_{V,\gamma}\cdot u,u\rangle^{p-q}\\
\wedge\big(\gamma\,\Theta_{A,h_{A,\delta}}-\rho_j^{-1}\,\Theta_{\Delta_j,h_j}\big)^{q-1}\,
\abs{e_j^*(u)}^{2q}
\bigg(\int_0^1(1-\tau)^{p-q+r-1}\,\tau^{q-1}\,d\tau\bigg)\wedge
\beta\,d\mu(u).
\end{multline}
We have
\begin{equation}\label{6.35}
\int_0^1(1-\tau)^{p-q+r-1}\,\tau^{q-1}\,d\tau=\tfrac{ (p-q+r-1)!\,(q-1)!}{ (p+r-1)!}
\end{equation}
and the combination of~\eqref{6.29} and (\eqref{6.32}--\eqref{6.35}) implies
\begin{multline}\label{6.36}
\lim_{\varepsilon\,\to\,0}I_{p,\varepsilon}(\beta)\ge
\int_{S(V)}\left\langle\theta_{V,\gamma}\cdot u,u\right\rangle^p\wedge\beta\,d\mu(u) +p(1-1/\rho_j)\,\sum_{q=1}^{p}\tfrac{(p-1)!\,(p-q+r-1)!}{ (p-q)!\,(p+r-1)!}\\
\times{}\int_{S(V)_{|\Delta_j}}
\left\langle\theta_{V,\gamma}\cdot u,u\right\rangle^{p-1-q}\wedge
\left(\gamma\,\Theta_{A,h_{A,\delta}}-\rho_j^{-1}\,\Theta_{\Delta_j,h_j}\right)^q\,
\left|e_j^*(u)\right|^{2q+2}\wedge\beta\,d\mu(u).
\end{multline}
Inductively, formula~\eqref{6.36} requires the investigation of more general integrals
\[
I_{p,p',Y,\varepsilon} =
\int_{S_\varepsilon(V\langle D\rangle)_{|Y}}
\left\langle\theta_{V,\gamma,\varepsilon}\cdot u,u\right\rangle^{p-p'}\wedge
\prod_{1\,\leq\,j\,\leq\,p'}\abs{\ell_j(u)}^2\,\beta\,d\mu_\varepsilon(u)
\]
where $Y$ is a subvariety of $X$ (which we assume to be transverse to the $\Delta_j$'s, and $\ell_j\in C^\infty(Y,V^*)$ with $\abs{\ell_j}=1$, and $\beta\geq_S 0$ is a smooth form of suitable bidegree on $Y$. Not much is changed in the calculation, except that the change of variable $u\mapsto g_{j,\varepsilon}\circ h_{j,\varepsilon}(u)$ applied to $\prod_{1\,\leq\,j\,\leq\,p'}\abs{\ell_j(u)}^2$ introduces an extra factor $(1-\tau)^{p'}$ in the lower bound, entirely compensated by the corresponding factor $(1-\tau)^{p-p'-q}$ appearing in $\langle\theta_{V,\gamma,\varepsilon}\cdot u,u\rangle^{p-p'}$. The binomial formula yields a coefficient $\binom{p-p'-1}{ q-1}$ instead of
$\binom{p-1}{ q-1}$. We thus obtain
\begin{multline}
\lim_{\varepsilon\,\to\,0}I_{p,p',Y,\varepsilon}(\beta)\ge
\int_{S(V)_{|Y}}\left\langle\theta_{V,\gamma}\cdot u,u\right\rangle^{p-p'}\wedge
\prod_{1\,\leq\,j\,\leq\,p'}\left|\ell_j(u)\right|^2\,\beta\,d\mu(u)\\
+(1-1/\rho_j)\,\sum_{q=1}^{p-p'}\tfrac{(p-p')!\,(p-q+r-1)!}{ (p-p'-q)!\,(p+r-1)!}\times{}\\
\int_{S(V)_{|Y\cap\Delta_j}}\kern-4pt
\left\langle\theta_{V,\gamma}\cdot u,u\right\rangle^{p-p'-q}\wedge
\big(\gamma\,\Theta_{A,h_{A,\delta}}-\rho_j^{-1}\,
\Theta_{\Delta_j,h_j}\big)^{q-1}\!
\wedge
\\
\left|e_j^*(u)\right|^{2q}\prod_{1\,\leq\,j\,\leq\,p'}\left|\ell_j(u)\right|^2\,
\beta\,d\mu(u).
\end{multline}
When $D$ contains several components, we apply induction on $N$ and put
\begin{align}
\abs{u}_{h_{V\langle D\rangle,\varepsilon}}^2&=
\abs{u}_{h_{V\langle D\rangle,\varepsilon'}}^2 +\varepsilon_N\,\abs{\sigma_N}^{-2+2/\rho_N}\,
\left|\nabla_N\sigma_N(u)\right|^2_{h_N}\\
\intertext{where}
\\ \tag{\theequation\ensuremath{'}}
\abs{u}_{h_{V\langle D\rangle,\varepsilon'}}^2&=
\abs{u}^2_{h_V}+\sum_{1\,\leq\,j\,\leq\,N-1}\varepsilon_j\,\abs{\sigma_j}^{-2+2/\rho_j}\,
\left|\nabla_j\sigma_j(u)\right|^2_{h_j}.
\end{align}
In this setting, \eqref{6.26} can be rewritten in the form of a decomposition
\begin{multline*}
\left\langle\theta_{V\langle D\rangle,\gamma,\varepsilon}\cdot u,u\right\rangle
\simeq\left\langle\theta_{V\langle D\rangle,\gamma,\varepsilon'}\cdot u,u\right\rangle\\
\begin{aligned}
&+\varepsilon_N\,\abs{\sigma_N}^{-2+2/\rho_N}\,
\big(\gamma\,\Theta_{A,h_{A,\delta}}-\rho_N^{-1}\,\Theta_{\Delta_N,h_N}\big)
\,\left|\nabla_N\sigma_N(u)\right|^2\\
&+\tfrac{1}{2\pi}
\tfrac{\varepsilon_N\,\abs{\sigma_N}^{-2+2/\rho_N}} { 1+\varepsilon_N\,\abs{\sigma_N}^{-2+2/\rho_N}\, \left|\nabla_N\sigma_N\right|^2}
\,\big|\,\nabla_N^2\sigma_N(\xi,u)-(1-1/\rho_N)\,
\sigma_N^{-1}\,\nabla_N\sigma_N(\xi)\nabla_N\sigma_N(u)\big|^2.
\end{aligned}
\end{multline*}
inductively with all intersections
$\Delta_J=\Delta_{j_1}\cap\ldots\cap\Delta_{j_m}$,
$J=\{j_1,\,\dotsc,\,j_m\}\subset\{1,\,\dotsc,\,N\}\,$; we neglect the self-intersection terms, since they are anyway non-negative. We obtain
\begin{multline}
\lim_{\varepsilon\,\to\,0}I_{p,\varepsilon}(\beta)\\\ge
\sum_{J\subset\{1,\,\dotsc,\,N\}}\abs{J}!\kern-5pt
\sum_{\substack{(q_j)\,\in\,(\bN^*)^J
\sum_{j\,\in\,J}q_j\,\leq\,p}}
\tfrac{p!\,\left(p+r-1-\sum_{j\,\in\,J}q_j\right)!}{(p+r-1)!\,\left(p-\sum_{j\in J}q_j\right)!}
\prod_{j\,\in\,J}(1-1/\rho_j)
\int_{z\,\in\,\Delta_J}
\int_{u\,\in\,S(V)_z}
\\
\left\langle\theta_{V,\gamma}(z)\cdot u,u\right\rangle^{p-\sum_{j\,\in\,J}q_j}\wedge
\bigwedge_{j\,\in\,J}\left|e_j^*(u)\right|^{2q_j}\,
\big(\gamma\,\Theta_{A,h_{A,\delta}}-\rho_j^{-1}\,\Theta_{\Delta_j,h_j}\big)^{q_j-1}
\wedge\beta(z)\,d\mu(u)
\end{multline}
where $J=\emptyset$ corresponds to the integral taken over $X$, with a coefficient equal to $1$ in that case. By the Fubini theorem, we get the following lower bound of $I_{k,p,\varepsilon}(\beta)$.

\begin{prop}\label{6.41}
With the same notation as above, assume that
\[
\gamma_s>\gamma_{V,D_s}\bydef\max\left(\max_j(d_j/\rho_{s,j}),\gamma_V\right),\quad 1\leq s\leq k.
\]
and consider the limit $\lim_{\varepsilon\,\to\,0} I_{k,p,\varepsilon}(\beta)$ computed as an iterated limit
\[
\lim_{\varepsilon_{11}\to 0}\ldots\lim_{\varepsilon_{kN}\,\to\,0}
\]
with respect to the lexicographic order $(i,j)<(i',j')$ if $i<i'$ or $i=i'$ and $j<j'$. Then we have the following lower bound, where the summation is taken over all disjoint subsets $J_1,\,\dotsc,\,J_k
\subset\{1,2,\,\dotsc,\,N\}$:
\begin{multline*}
\lim_{\varepsilon\,\to\,0} I_{k,p,\varepsilon}(\beta)\geq{}\kern-8pt
\sum_{\substack{ J_1\amalg\ldots\amalg J_k\\
\subset\{1,\,\dotsc,\,N\}}}
\sum_{\substack{(q_j)\,\in\,(\bN^*)^{J_1\amalg\ldots\amalg J_k}\\
\sum_{j\,\in\,J_s}q_j\,\leq\,p_s}}
\prod_{1\,\leq\,s\,\leq\,k}\kern-4pt\tfrac{\abs{J_s}!\,p_s!\,\left(p_s-\sum_{j\,\in\,J_s}q_j+r-1\right)!}{ \left(p_s+r-1\right)!\,\left(p_s-\sum_{j\,\in\,J_s}q_j\right)!}\kern-2pt
\prod_{j\,\in\,J_s}\kern-2pt\left(1-\tfrac{1}{\rho_{s,j}}\right)\kern-3pt\\
\int_{z\,\in\,\Delta_{J_1\amalg\ldots\amalg J_k}}\int_{(u_s)\,\in\,S(V)^k_z}~
\bigwedge_{1\,\leq\,s\,\leq\,k}\Bigg(
\left\langle\theta_{V,\gamma_s}\cdot u_s,u_s\right\rangle^{p_s-\sum_{j\,\in\,J_s}q_j}
\wedge{}\\
\left. \bigwedge_{j\,\in\,J_s}\left|e_j^*(u_s)\right|^{2q_j}\,
\left(\gamma_s\,\Theta_{A,h_{A,\delta}}-\rho_{s,j}^{-1}\,\Theta_{\Delta_j,h_j}
\right)^{q_j-1}\,d\mu(u_s)\right)\wedge\beta(z).
\end{multline*}
\end{prop}


Our assumptions imply that we can take
$\theta_ {V,\gamma_s}>_G(\gamma_s-\gamma_V-\delta)\Theta_{A,h_A}\otimes\Id_V$ for every $\delta>0$. By Lemma~\ref{lemm:Chern_form_inequalities}$\MK$\eqref{item:Chern_curv_inequality2}, we obtain the simpler and purely cohomological lower bound
\begin{multline}\label{eq:simpler_bound}
\lim_{\varepsilon\,\to\,0} I_{k,p,\varepsilon}(\beta)\\
\geq{}\kern-8pt
\sum_{\substack{ J_1\amalg\,\ldots\amalg J_k\\\subset\{1,\,\dotsc,\,N\}}}
\sum_{\substack{(q_j)\,\in\,(\bN^*)^{J_1\amalg\ldots\amalg J_k}\\
\sum_{j\,\in\,J_s}q_j\,\leq\,p_s}}
\prod_{1\,\leq\,s\,\leq\,k}\kern-4pt\tfrac{\abs{J_s}!\,p_s!\,\left(p_s-\sum_{j\,\in\,J_s}q_j+r-1\right)!}{ (p_s+r-1)!\,\left(p_s-\sum_{j\,\in\,J_s}q_j\right)!}
\prod_{j\,\in\,J_s}\kern-2pt\Big(1-\tfrac{1}{\rho_{s,j}}\Big)\\
\int_{\Delta_{J_1\amalg\ldots\amalg J_k}}~
\bigwedge_{1\,\leq\,s\,\leq\,k}\bigg(
\big((\gamma_s-\gamma_V)\Theta_{A,h_A}\big)^{p_s-\sum_{j\,\in\,J_s}q_j}
\wedge{}\\
\bigwedge_{j\,\in\,J_s}\tfrac{(r-1)!}{\left(q_j+r-1\right)!}\,
\big(\gamma_s\,\Theta_{A,h_{A,\delta}}-\rho_{s,j}^{-1}\,\Theta_{\Delta_j,h_j}
\big)^{q_j-1}\bigg)\wedge\beta(z).
\end{multline}
What is a bit surprising in all these estimates is that, in spite of the fact that we are integrating non-closed and metric dependent forms, the limits of the integrals as $\varepsilon\to 0$ admit rather natural lower and upper bounds that are purely cohomological, and can be expressed solely in terms of well understood Chern classes. This will also be true for the related Morse integrals in~\S~\ref{se:non_proba_estimates}. It could be desirable to have an algebro-geometric explanation of this phenomenon. The algebraic versions of Morse inequalities developed by B.~Cadorel in~\cite{Cad19} might possibly be used in this context.

\begin{rema}\label{6.43}
As mentioned in the course of the proof of~\eqref{6.41}--\eqref{eq:simpler_bound}, we have neglected certain non-negative terms coming from self-intersections $\Delta_j^p$ of the components ($p\geq 2$), by restricting the summation to the family of disjoint subsets $J_1,J_2,\,\dotsc,\,J_k$. It would be interesting to refine the lower bound and to take these terms into account. This might be possible by observing that the iterated limit process, when integrating on $\Delta_j$, involves inductively a few extra terms in~\eqref{6.30}, when we take the limit as $t_{j,\varepsilon}\to\infty$. Those terms are equal to
\[
\left\langle\theta_{V,\gamma}\cdot u'_j,u'_j\right\rangle+
\big(\gamma\,\Theta_{A,h_{A,\delta}}-\rho_j^{-1}\,\Theta_{\Delta_j,h_j}\big)
\,\left|e_j^*(u)\right|^2+\tfrac{\imath}{ 2\pi}\big\langle b_j(u'_j),b_j(u'_j)\big\rangle_{h_j}.
\]
One would then have to evaluate the contribution of $\langle b_j(u'_j),b_j(u'_j)\rangle_{h_j}$ in the integral~$\int_{\Delta_j}$.
\end{rema}

\section{Non probabilistic estimates of the Morse integrals}\label{se:non_proba_estimates}

The non-probabilistic estimate uses more explicit curvature inequalities and has the advantage of producing results also in the general orbifold case. Let us fix an ample line bundle $A$ on $X$ equipped with a smooth hermitian metric $h_A$ such that $\omega_A\bydef\Theta_{A,h_A}>0$, and let $\gamma_V$ be the infimum of values $\lambda\in\bR_+$ such that
\[
\lambda\,\omega_A\otimes\Id_V -
\Theta_{V,h_V} > _G 0,
\]
in the sense of Griffiths. For any orbifold structure $D=\sum_j(1-1/\rho_j)\Delta_j$, Corollary~\ref{coro:curv_estimate} then shows that the $s^{\rm th}$ directed orbifold bundle $V_s\bydef V\langle D^{(s)}\rangle$ (cf. \S~\ref{sse:log_directed}) possesses hermitian metrics $h_{V\langle D^{(s)}\rangle,\varepsilon_s}$ such that the associated curvature tensor satisfies the inequality
\[
\theta_{s,\gamma,\varepsilon}
\bydef
\gamma_s\,\omega_A\otimes\Id_{V\left\langle D^{(s)}\right\rangle} -\Theta_{V\left\langle D^{(s)}\right\rangle,h_{V\left\langle D^{(s)}\right\rangle,\varepsilon_s}}>_G 0,
\]
provided we assume $d_jA-\Delta_j$ nef and take
\begin{equation}
\gamma_s > \gamma_{V,D^{(s)}}
\bydef
\max\left(\max_j\left(d_j/\rho_j^{(s)}\right),\gamma_V\right)
\quad\text{where}~~
\rho_j^{(s)}=\max(\rho_j/s,1).
\end{equation}
In particular, any value
\begin{equation'}\label{7.3'}
\gamma_s > \max\left(s\,\max_j\left(d_j/\rho_j\right),\gamma_V\right).
\end{equation'}
\noindent is admissible, and we can apply the estimates~\ref{6.41}~(b) and~\eqref{eq:simpler_bound} with these values. Instead of exploiting a Monte Carlo convergence process for the curvature tensor as was done in \S~\ref{sse:proba_estimate}, we are going to use a more precise lower bound of the curvature tensor $\Theta_{L_{\tau,k},\varepsilon}$ of the orbifold rank $1$ sheaf associated with $F=\tau A$, $\tau\ll 1$, namely
\[
L_{\tau,k}
\bydef
\cO_{X_k(V\langle D\rangle)}(1)\otimes
\pi_k^*\cO_X(-\tau A).
\]
Our formulas~\ref{prop:orbi_curvature}~(a,b) become
\begin{align*}
\Theta_{L_{\tau,k},\varepsilon}& =
\omega_{r,k,b}(\xi)+g_{k,0,\varepsilon}(z,x,u) -\tau\,\omega_A(z),\quad
\text{where}
\\
g_{k,\gamma,\varepsilon}(z,x,u)& =
\sum_{s=1}^k\tfrac{x_s}{ s}\,
\theta_{s,\gamma,\varepsilon}(u_s),
\\
\theta_{s,\gamma,\varepsilon}(u_s)& =
\tfrac{\imath}{ 2\pi}
\sum_{i,j,\lambda,\mu}c^{(s,\gamma,\varepsilon)}_{ij\lambda\mu}(z)\, u_{s,\lambda}\bar u_{s,\mu}\,dz_i\wedge d\bar z_j.
\end{align*}
Under the assumption~\eqref{7.3'}, we have $g_{k,\gamma,\varepsilon}(z,x,u)\geq 0$, but in general this is not true for $g_{k,0,\varepsilon}(z,x,u)$, so we express $g_{k,0,\varepsilon}(z,x,u)$ as a difference of $g_{k,\gamma,\varepsilon}(z,x,u)$ and of a multiple of $\omega_A$. By definition $\theta_{s,\gamma,\varepsilon}=\gamma_s\,\omega_A\otimes\Id+\theta_{s,0,\varepsilon}$, and we get
\begin{align}\label{7.6}
\Theta_{L_{\tau,k},\varepsilon}&=\omega_{r,k,b}+\alpha_\varepsilon-\beta,
\qquad\text{where}
\\
\tag{\theequation\ensuremath{'}} \alpha_\varepsilon &=g_{k,\gamma,\varepsilon}\geq 0,\quad \beta=
\left(\tau+\sum_{1\,\leq\,s\,\leq\,k}\tfrac{\gamma_sx_s}{ s}\right)\omega_A =\sum_{1\,\leq\,q\,\leq\,k}\tfrac{(\gamma_q+q\tau)\,x_q}{ q}\,\omega_A\geq 0.
\end{align}
Then~\eqref{7.6} and the inequalities used for~\eqref{eq:volume}, especially Lemma~\ref{lemm:forms_ineq} and Proposition~\ref{prop:volume_form}$\MK$\eqref{prop4.3.b}, lead to
\begin{multline}\label{7.7}
\int_{X_k(V\langle D\rangle)\left(L_{\tau,k},\,\leq\,1\right)}
\Theta_{L_{\tau,k},\varepsilon}^{n+kr-1}\\
\begin{aligned}
&=\tfrac{(n+kr-1)!}{n!\,k!^r(kr-1)!}
\int_{z\,\in\,X}
\int_{(x,u)\,\in\,\bDelta^{k-1}\times\left(\bS^{2r-1}\right)^k}
\bOne_{\alpha_\varepsilon-\beta,\,\leq\,1}
(\alpha_\varepsilon-\beta)^n d\nu_{k,r}(x) d\mu(u)\\
&\geq \tfrac{(n+kr-1)!}{n!\,k!^r(kr-1)!}
\int_{z\,\in\,X}
\int_{(x,u)\,\in\,\bDelta^{k-1}\times\left(\bS^{2r-1}\right)^k}
\left(\alpha_\varepsilon^n-n\alpha_\varepsilon^{n-1}\wedge\beta\right) d\nu_{k,r}(x) d\mu(u).
\end{aligned}
\end{multline}
The main point is thus to find a lower bound of the difference $\alpha_\varepsilon^n-n\alpha_\varepsilon^{n-1}\wedge\beta$, hence a lower bound of $\alpha_\varepsilon^n$ and an upper bound of $\alpha_\varepsilon^{n-1}\wedge\beta$. An expansion of $\alpha_\varepsilon^n$ by Newton's multinomial formula yields
\[
\alpha_\varepsilon^n=\sum_{p\,\in\,\bN^k,\,\abs{p}=n}\tfrac{n!}{ p_1!\ldots\,p_k!}\,
\prod_{s=1}^k\Big(\tfrac{x_s}{ s}\,\theta_{s,\gamma,\varepsilon}(u_s)\Big)^{p_s}.
\]
If we assume $k\geq n$ and retain only the monomials for which $p_s=0,1$, we get
\[
\alpha_\varepsilon^n\geq\sum_{1\,\leq\,s_1<\,\ldots\,<s_n\,\leq\,k}\,\tfrac{n!}{s_1\ldots s_n}\,
\prod_{\ell=1}^n x_{s_\ell}\theta_{s_\ell,\gamma,\varepsilon}(u_{s_\ell}).
\]
By Proposition~\ref{prop:volume_form}$\MK$\eqref{prop4.3.a} and an elementary calculation (cf.~\cite[Prop.~1.13]{Dem11}), one gets for every $(p_1,\,\dotsc,\,p_k)\in\bN^k$
\begin{equation}\label{7.9}
\int_{\bDelta^{k-1}}x_1^{p_1}\ldots x_k^{p_k}\, d\nu_{k,r}(x)=\tfrac{(kr-1)!}{ (r-1)!^k}\;
\tfrac{\prod_{1\,\leq\,s\,\leq\,k}(p_s+r-1)!}{\left(\sum_{1\,\leq\,s\,\leq\,k}p_s+kr-1\right)!},
\end{equation}
and in particular, for $k\geq n$, $p_1=\ldots=p_n=1$, $p_{n+1}=\ldots=p_k=0$, we have
\[
\int_{\bDelta^{k-1}}x_{s_1}\ldots x_{s_n}\,d\nu_{k,r}(x) =\int_{\bDelta^{k-1}}x_1\ldots x_n\,d\nu_{k,r}(x) =\tfrac{(kr-1)!\,r^n}{(n+kr-1)!}.
\]
As a consequence, the equality case in (\eqref{6.20}--\eqref{6.21}) implies
%\stepcounter{equation}
\begin{multline*}
M_{n,k,\varepsilon} \bydef
\int_{z\,\in\,X}\int_{(x,u)\,\in
\,\bDelta^{k-1}\,\times\left(\bS^{2r-1}\right)^k}
\alpha_\varepsilon(z)^n\,d\nu_{k,r}(x)\,d\mu(u_1)\ldots\,d\mu(u_k)\\
\begin{aligned}
&\geq\sum_{1\,\leq\,s_1\,<\,\ldots\,<\,s_n\,\leq\,k}
\int_{\bDelta^{k-1}}
\tfrac{n!\,x_{s_1}\ldots x_{s_n}}{s_1\ldots s_n}\,d\nu_{k,r}(x)\\
&\quad\times \int_X\int_{\raise-1pt\hbox{${\textstyle\Pi}$}\kern1pt S(V\langle D^{(s_\ell)}\rangle)}~\bigwedge_{\ell=1}^n
\langle\theta_{s_\ell,\gamma,\varepsilon}(u_{s_\ell}),u_{s_\ell}\rangle\, d\mu(u_{s_\ell})\\
&\geq\sum_{1\,\leq\,s_1\,<\,\ldots\,<\,s_n\,\leq\,k}\tfrac{(kr-1)!}{(n+kr-1)!} \,
\tfrac{n!}{ s_1\ldots s_n}\\
&\quad\times \int_X~\prod_{\ell=1}^n\left(r\gamma_{s_\ell}\,c_1(A)-c_1(V) +\sum_j\left(1-1/\rho^{(s_\ell)}_j\right)c_1(\Delta_j)\right)
\end{aligned}
\end{multline*}
 \begin{equation_}{1}\label{7.10_1}
%\tag{$\theequation_1$}
M_{n,k,\varepsilon} \geq
\tfrac{(kr-1)!}{ (n+kr-1)!} \,
\int_X~\prod_{s=1}^n\bigg(r\gamma_s\,c_1(A)-c_1(V)+\sum_j \left(1-1/\rho^{(s)}_j\right)c_1(\Delta_j)\bigg).
\end{equation_}
If we assume $c_1(V^*)=\lambda_V\,c_1(A)$ and $c_1(\Delta_j)=d_j\,c_1(A)$, the lower bound takes the simpler form
\begin{equation_}{2}\label{7.10_2}
M_{n,k,\varepsilon}\geq\tfrac{(kr-1)!}{(n+kr-1)!}\,
\prod_{s=1}^n\left(r\gamma_s+\lambda_V+\sum_j d_j\left(1-1/\rho^{(s)}_j\right)\right)\,A^n.
\end{equation_} In fact, our lower bounds are obtained by taking into account the single term $s_\ell=\ell$, $1\leq\ell\leq k$ (which is the unique term in the sum when $k=n$). A more refined method is to integrate all monomials
$x_1^{p_1}\ldots\,x_k^{p_k}$ and to use the lower bound~\eqref{eq:simpler_bound} instead of (\eqref{6.20}--\eqref{6.21}). This has the advantage of eventually producing a non-zero contribution, even when $k<n$. We find
\begin{align*}
M_{n,k} &\bydef\lim_{\varepsilon\,\to\,0}\int_{z\,\in\,X}\int_{(x,u)\in
\bDelta^{k-1}
\times\left(\bS^{2r-1}\right)^k}
\alpha_\varepsilon(z)^n\,d\nu_{k,r}(x)\,d\mu(u_1)\ldots\,d\mu(u_k)\\
&\geq\lim_{\varepsilon\,\to\,0}
\sum_{\substack{p\,\in\,\bN^k\\ \abs{p}=n}}
\int_{\bDelta^{k-1}}
\tfrac{n!\,x_1^{p_1}\ldots x_k^{p_k}}{\prod_{s=1}^kp_s!\,s^{p_s}}\,d\nu_{k,r}(x)
\int_X\int_{\raise-1pt\hbox{$\Pi$}\; S\left(V\left\langle D^{(s)}\right\rangle\right)}\bigwedge_{s=1}^k
\left\langle\theta_{s,\gamma,\varepsilon}(u_s),u_s\right\rangle^{p_s}\, d\mu(u_s)
\\
&\ge
\sum_{\substack{p\,\in\,\bN^k\\ \abs{p}=n}}
\tfrac{n!}{\prod_{s=1}^kp_s!\,s^{p_s}}\,
\tfrac{(kr-1)!}{(r-1)!^k}\;
\tfrac{\prod_{1\,\leq\,s\,\leq\,k}\left(p_s+r-1\right)!}{\left(\sum_{1\,\leq\,s\,\leq\,k}p_s+kr-1\right)!}~
\sum_{\substack{ J_1\amalg\ldots\amalg J_k\\
\subset\{1,\,\dotsc,\,N\}}}
\\
&\sum_{\substack{ (q_j)\,\in\,(\bN^*)^{J_1\amalg\ldots\amalg J_k} \\
\sum_{j\,\in\,J_s}q_j\,\leq\,p_s}}
\prod_{1\,\leq\,s\,\leq\,k}\tfrac{\abs{J_s}!\,p_s!\,\left(p_s-\sum_{j\,\in\,J_s}q_j+r-1\right)!}{(p_s+r-1)!\, \left(p_s-\sum_{j\,\in\,J_s}q_j\right)!}\,\prod_{j\,\in\,J_s}\left(1-\tfrac{1}{ \rho^{(s)}_j}\right)
\int_{z\,\in\,\Delta_{J_1\amalg\ldots\amalg J_k}}\\
&
\bigwedge_{1\,\leq\,s\,\leq\,k}\big((\gamma_s-\gamma_V)\Theta_{A,h_A}\big)^{p_s -\sum_{j\in J_s}q_j}
\\
&\wedge\bigwedge_{j\,\in\,J_s}\tfrac{(r-1)!}{ (q_j+r-1)!}\,
\big(\gamma_s\,\Theta_{A,h_A}- \left(\rho^{(s)}_j\right)^{-1}\,\Theta_{\Delta_j,h_j}\big)^{q_j-1},
\end{align*}
thus
\begin{multline}
M_{n,k}\geq\tfrac{n!\,(kr-1)!}{(n+kr-1)!}
\sum_{\substack{p\,\in\,\bN^k\\
\abs{p}=n}}~
\prod_{1\,\leq\,s\,\leq\,k}\tfrac{1}{ s^{p_s}}\,
\sum_{\substack{ J_1\amalg\ldots\amalg J_k
\\
\subset\{1,\,\dotsc,\,N\}}}
\\
\sum_{\substack{(q_j)\,\in\,(\bN^*)^{J_1\amalg\ldots\amalg J_k}
\\ \sum_{j\,\in\,J_s}q_j\,\leq\,p_s}}
\prod_{1\,\leq\,s\,\leq\,k}\tfrac{\abs{J_s}!\,\left(p_s-\sum_{j\,\in\,J_s}q_j+r-1\right)!}{\left(p_s-\sum_{j\,\in\,J_s}q_j\right)!}\,\prod_{j\,\in\,J_s}
\left(1-\tfrac{1}{\rho^{(s)}_j}\right)
\int_{z\,\in\,\Delta_{J_1\amalg\ldots\amalg J_k}}\bigwedge_{1\,\leq\,s\,\leq\,k}
\\
\big((\gamma_s-\gamma_V)\Theta_{A,h_A}\big)^{p_s-\sum_{j\,\in\,J_s}q_j}
\wedge\bigwedge_{j\,\in\,J_s}\tfrac{(r-1)!}{\left(q_j+r-1\right)!}\,
\big(\gamma_s\,\Theta_{A,h_A}- \left(\rho^{(s)}_j\right)^{-1}\,\Theta_{\Delta_j,h_j}\big)^{q_j-1}.
\end{multline}
In particular, if $c_1(\Delta_j)=d_j\,c_1(A)$, we infer
\begin{multline}\label{7.12}
M_{n,k}\geq\tfrac{n!\,(kr-1)!}{(n+kr-1)!}
\sum_{\substack{p\,\in\,\bN^k\\ \abs{p}=n}}
\prod_{1\,\leq\,s\,\leq\,k}\tfrac{1}{ s^{p_s}}\kern-5pt
\sum_{\substack{J_1\amalg\ldots\amalg J_k\\
\subset\{1,\dotsc,N\}}}
\sum_{\substack{(q_j),\,q_j\geq\,1\\ \sum_{j\,\in\,J_s}q_j\,\leq\,p_s}}
\\
\prod_{1\,\leq\,s\,\leq\,k}\!
\Bigg(
\tfrac{\abs{J_s}!\,\left(p_s-\sum_{j\,\in\,J_s}q_j+r-1\right)!}{\left(p_s-\sum_{j\,\in\,J_s}q_j\right)!} (\gamma_s-\gamma_V)^{p_s-\sum_{j\in J_s}q_j}\,
\\
\prod_{j\,\in\,J_s}d_j\left(1-\tfrac{1}{ \rho^{(s)}_j}\right)\,
\tfrac{(r-1)!}{\left(q_j+r-1\right)!}
\left(\gamma_s-\tfrac{d_j}{\rho^{(s)}_j}\right)^{q_j-1}
\Bigg)
\,A^n.
\end{multline}
In the special case $k=1$ and $N\geq n$, by taking $\abs{J}=\abs{J_1}=n$ and $q_j=1$ for all $j\in J$, we find
\begin{equation_}{1}\label{7.12_1}
M_{n,1}\geq\tfrac{n!\,(r-1)!}{(n+r-1)!}
\sum_{J\subset\{1,\,\dotsc,\,N\},\,\abs{J}=n}
\tfrac{n!\,(r-1)!}{ r^n}\,
\prod_{j\,\in\,J}d_j\left(1-\tfrac{1}{ \rho_j}\right)\,A^n.
\end{equation_}
Next, we turn ourselves to the evaluation of the integral of $\alpha_\varepsilon^{n-1}\wedge\beta$. We have
\[
\alpha_\varepsilon^{n-1}\wedge\beta=\sum_{p\,\in\,\bN^k,\,\abs{p}=n-1}
\tfrac{(n-1)!}{p_1!\ldots\,p_k!}\,
\prod_{s=1}^k\left(\tfrac{x_s}{ s}\,\theta_{s,\gamma,\varepsilon}(u_s)\right)^{p_s}
\wedge\beta,
\]
and the upper bound given by (\eqref{6.20}--\eqref{6.21}) provides

\break
\begin{multline*}
M'_{n,k} \bydef
\lim_{\varepsilon\,\to\,0}\int_{z\,\in\,X}
\int_{(x,u)\,\in\,\bDelta^{k-1}\times(\bS^{2r-1})^k} n\,\alpha_\varepsilon(z)^{n-1}\wedge\beta\, d\nu_{k,r}(x)\,d\mu(u_1)\ldots\,d\mu(u_k)\\
\begin{aligned}
&\leq\lim_{\varepsilon\,\to\,0}\sum_{p\,\in\,\bN^k,\,\abs{p}=n-1}\int_{\bDelta^{k-1}} n\,\tfrac{(n-1)!\,x_1^{p_1}\ldots x_k^{p_k}}{
\prod_{s=1}^kp_s!\,s^{p_s}}\,d\nu_{k,r}(x)\\
&\quad\times{}\int_X\int_{\raise-1pt\hbox{${\textstyle\Pi}$}\kern1pt S\left(V\left\langle D^{(s)}\right\rangle\right)}\bigwedge_{s=1}^k
\left\langle\theta_{s,\gamma,\varepsilon}(u_s),u_s\right\rangle^{p_s}\wedge\beta\,
\prod_{s=1}^kd\mu(u_s)\\
&\leq\sum_{p\,\in\,\bN^k,\,\abs{p}=n-1}\int_{\bDelta^{k-1}} n\,\tfrac{(n-1)!\,x_1^{p_1}\ldots x_k^{p_k}}{\prod_{s=1}^kp_s!\,s^{p_s}}\,
\left(\sum_{q=1}^k\tfrac{(\gamma_q+q\tau)\,x_q}{q}\right)\,d\nu_{k,r}(x)\\
&\quad \times \int_X~\bigwedge_{1\,\leq\,s\,\leq\,k}\tfrac{p_s!\,(r-1)!}{\left(p_s+r-1\right)!}\,
\Bigg(r\gamma_s\,\Theta_{A,h_A}-\Tr\Theta_{V,h_V}\\
&\mkern150mu \left.+\sum_j \left(1-1/\rho^{(s)}_j\right)\Theta_{\Delta_j,h_j}\right)^{p_s}\wedge\Theta_{A,h_A}.
\end{aligned}
\end{multline*}
By~\eqref{7.9}, for $\abs{p}=\sum p_s=n-1$, we get
\begin{multline*}
\int_{\bDelta^{k-1}}x_1^{p_1}\ldots x_k^{p_k}\left(\sum_{q=1}^k\tfrac{\gamma_q}{ q} +\tau x_q\right)d\nu_{k,r}(x)
\\
\begin{aligned}
&=
\tfrac{(kr-1)!}{ (r-1)!^k}\;
\tfrac{\prod_{1\,\leq\,s\,\leq\,k}\left(p_s+r-1\right)!}{(n-1+kr-1)!}
\left(\sum_{q=1}^k\tfrac{\gamma_q}{ q}+\tau\sum_{q=1}^k\tfrac{p_q+r}{ n+kr-1}\right)
\\
&=
\tfrac{(kr-1)!}{(r-1)!^k}\;
\tfrac{\prod_{1\,\leq\,s\,\leq\,k}\left(p_s+r-1\right)!}{(n+kr-2)!}
\left(\sum_{q=1}^k\tfrac{\gamma_q}{ q}+\tau\right).
\end{aligned}
\end{multline*}
Therefore, assuming $c_1(\Delta_j)=d_j\,c_1(A)$ and
$c_1(V^*)=\lambda_Vc_1(A)$, we find
\stepcounter{equation}
\begin{equation}
\tag{$\theequation_1$}
\begin{split}
M'_{n,k}
&\leq \tfrac{n!\,(kr-1)!}{ (r-1)!^k\,(n+kr-2)!}\,
\left(\sum_{q=1}^k\tfrac{\gamma_q}{q}+\tau\right)
\sum_{p\,\in\,\bN^k,\,\abs{p}=n-1}\tfrac{\prod_{1\,\leq\,s\leq k}(p_s+r-1)!}{
\prod_{s=1}^kp_s!\,s^{p_s}}\\
&\quad\times
\prod_{1\leq s\leq k}\tfrac{p_s!\,(r-1)!}{(p_s+r-1)!}\,
\Big(r\gamma_s+\lambda_V+\sum_jd_j\left(1-1/\rho^{(s)}_j\right)\Big)^{p_s}\,A^n,
\\
&\leq \tfrac{n!\,(kr-1)!}{(n+kr-2)!}\,
\left(\sum_{q=1}^k\tfrac{\gamma_q}{ q}+\tau\right)\\
&\quad\times{}
\sum_{p\,\in\,\bN^k,\,\abs{p}=n-1}~
\prod_{1\,\leq\,s\,\leq\,k}\tfrac{1}{ s^{p_s}}\,\left(r\gamma_s +\lambda_V+\sum_jd_j\left(1-1/\rho^{(s)}_j\right)\right)^{p_s}\,A^n.
\end{split}
\end{equation}
A simpler (but larger) upper bound is
\begin{equation_}{2}\label{7.14_2}
M'_{n,k}\leq\tfrac{n!\,(kr-1)!}{(n+kr-2)!}\,
\left(\sum_{s=1}^k\tfrac{\gamma_s}{ s}+\tau\right)
\left(\sum_{1\,\leq\,s\,\leq\,k}\tfrac{1}{ s}\,\left(r\gamma_s+
\lambda_V+\sum_jd_j\left(1-1/\rho^{(s)}_j\right)\right)\right)^{n-1}\,A^n.
\end{equation_}
Finally, inequality~\eqref{7.7} translates into
\begin{equation}
\tfrac{1}{ (n+kr-1)!}\int_{X_k(V\langle D\rangle)\left(L_{\tau,k},\leq\,1\right)}
\Theta_{L_{\tau,k},\varepsilon}^{n+kr-1}
\geq\tfrac{1}{ n!\,k!^r(kr-1)!}\left(M_{n,k}-M'_{n,k}\right).
\end{equation}
If we put everything together, we get the following (complicated!) existence criterion for orbifold jet differentials.

\begin{enonce*}
{Existence criterion}
Let $(X,V,D)$ with
$D=\sum_{1\,\leq\,j\,\leq\,N}(1-1/\rho_j)\Delta_j$ be a directed orbifold, and let $A$ be an ample line bundle on $X$. Assume that $D$ is a simple normal crossing divisor transverse to $V$, that $c_1(\Delta_j)=d_j\,c_1(A)$,
$c_1(V^*)=\lambda_V\,c_1(A)$ and let $\gamma_V$ be the infimum of values $\gamma>0$ such that $\Theta_A\otimes\Id_V-\Theta_V\geq_G 0$. Take
\[
\gamma_s=\max\left(\max\left(d_j/\rho_j^{(s)}\right),\gamma_V\right),\quad
\rho_j^{(s)}=\max\left(\rho_j/s,1\right).
\]
Then, a sufficient condition for the existence of (many) non-zero holomorphic sections of multiples of
\[
L_{\tau,k}=\cO_{X_k(V\langle D\rangle)}(1)\otimes\pi_k^*\cO(-\tau A)
\]
on $X_k(V\langle D\rangle)$ is that $M_{n,k}-M'_{n,k}>0$, where $M_{n,k}$ admits the lower bounds~\eqref{7.10_2} or~\eqref{7.12}, and $M'_{n,k}$ admits the upper bound~\eqref{7.14_2}.
\end{enonce*}

\subsection{Compact case (no boundary divisor)}

We address here the case of a compact (projective) directed manifold $(X,V)$, with a boundary divisor $D=0$. By~\eqref{7.10_2} and~\eqref{7.14_2}, we find
\begin{align*}
M_{n,k}&\geq \tfrac{(kr-1)!}{(n+kr-1)!}(r\gamma_V+\lambda_V)^n\,A^n\quad
\text{if $k\geq n$},\\
M'_{n,k}&\leq \tfrac{n!\,(kr-1)!}{(n+kr-1)!}
\Bigg(\tau+\gamma_V\sum_{s=1}^k\tfrac{1}{ s}\Bigg)\Bigg(\sum_{s=1}^k
\tfrac{1}{ s}\big(r\gamma_V+\lambda_V\big)\Bigg)^{n-1}.
\end{align*}
Therefore, for $\tau>0$ sufficiently small, $M_{n,k}-M'_{n,k}$ is positive as soon as $k\geq n$ and $(r\gamma_V+\lambda_V)^n>n!\,\gamma_V(
\sum_{1\,\leq\,s\,\leq\,k}\tfrac{1}{ s})^n(r\gamma_V+\lambda_V)^{n-1}$, that is
\begin{equation}\label{7.17}
k\geq n\quad 
\text{and}\quad
\lambda_V > n!
\left(\sum_{1\,\leq\,s\,\leq\,k}\tfrac{1}{ s}\right)^n\gamma_V - r\gamma_V.
\end{equation}

\begin{exam}
In the case where $X$ is a smooth hypersurface of $\bP^{n+1}$ of degree $d$ and $V=T_X$, we have $r=n$ and $\det(V^*)=\cO(d-n-2)$. We take $A=\cO(1)$. If $Q$ is the tautological quotient bundle on
$\bP^{n+1}$, it is well known that $T_{\bP^{n+1}}\simeq Q\otimes\cO(1)$ and
$\det Q=\cO(1)$, hence
$T^*_{\bP^{n+1}}\otimes\cO(2)=Q^*\otimes \cO(1)=\Lambda^nQ\geq_G0$, and the surjective morphism
\[
T^*_{\bP^{n+1}|X}\to T^*_X=V^*
\]
implies that we also have $V^*\otimes\cO(2)\geq_G 0$. Therefore, we find $\gamma_V=2$ and $\lambda_V=d-n-2$. The above condition~\eqref{7.17} becomes $k\geq n$ and
\[
k\geq n\quad\text{and}\quad d>2\,n!\,\left(\sum_{1\,\leq\,s\,\leq\,k}\tfrac{1}{s}\right)^n-n+2.
\]
This lower bound improves the one stated in~\cite{Dem12}, but is unfortunately far from being optimal. Better bounds---still probably non-optimal---have been obtained in~\cite{Dar16,MTa19}.
\end{exam}

\subsection{Logarithmic case}

The logarithmic situation makes essentially no difference in treatment with the compact case, except for the fact that we have to replace $V$ by the logarithmic directed structure $V\langle D\rangle$, and the numbers
$\gamma_V$, $\lambda_V$ by
\stepcounter{equation}
\begin{align*}
&\gamma_{V\langle D\rangle}=\inf\gamma~~\text{such that}~~
\gamma\Theta_A-\Theta_{V\langle D\rangle}\geq_G 0,
\\
&\lambda_{V\langle D\rangle}~~\text{such that}~~c_1(V^*\langle D\rangle)=\lambda_{V\langle D\rangle}\,c_1(A)\quad
\text{(if such $\lambda_{V\langle D\rangle}$ exists)}.
\end{align*}
We get the sufficient condition
\[
k\geq n\quad\text{and}\quad
\lambda_{V\langle D\rangle}>n!\,\left(\sum_{1\,\leq\,s\,\leq\,k}\tfrac{1}{ s}
\right)^n\gamma_{V\langle D\rangle}-r\gamma_{V\langle D\rangle}.
\]
For $X=\bP^n$, $V=T_{\bP^n}$, and for a divisor $D=\sum\Delta_j$ of total degree $d$ on $\bP^n$, we can still take $\gamma_{V\langle D\rangle}=2$ by {Lemma}~\ref{lemm:Chern_form_inequalities}, and we have $\det(V^*\langle D\rangle)=\cO(d-n-1)$. We get the degree condition
\[
k
\geq n
\text{\quad and\quad} d > 2\,n!\,\left(\sum_{1\,\leq\,s\,\leq\,k}\tfrac{1}{s}\right)^n-n+1.
\]
Again,~\cite{Dar16,MTa19} gave better bounds for this particular logarithmic situation.

\subsection{Case of orbifold structures on projective space}

Let us come to our main target, namely ``genuine'' orbifolds, for which our results are completely new. The situation we have in mind is the case of triples $(X,V,D)$ where $X=\bP^n$, $V=T_X$, $D=\sum (1-1/\rho_j)\Delta_j$ is a normal crossing divisor, with components $\Delta_j$ of degree $d_j$. Set again $A=\cO(1)$. Since $c_1(V^*)=-(n+1)\,c_1(A)$ and $D^{(s)}=\sum_j(1-s/\rho_j)_+\Delta_j$, we have
\[
\lambda_V=-n-1,\quad \det V^*\left\langle D^{(s)}\right\rangle =
\cO_{\bP^n}\left(-n-1+\sum_j d_j\left(1-s/\rho_j\right)_+\right).
\]
Moreover, by {Lemma}~\ref{lemm:Chern_form_inequalities}, we get
\[
\Theta_{V^*\langle D^{(s)}\rangle}+\gamma_s\,\omega_\FS\otimes\Id>_G0
\]
as soon as $\gamma_s>2$ and
$\gamma_s>\max_j(d_j/\max(\rho_j/s,1))$ for all components $\Delta_j$ in $D^{(s)}$. We can take for instance
$\gamma_s>st$ where $t=\max(\max_j(d_j/\rho_j),2)$. By considering the infimum and applying~\eqref{7.10_2} when $r=n$ and $k\geq n$, we find
\[
M_{n,k,\varepsilon}\geq\tfrac{(kn-1)!}{(n+kn-1)!}\,
\prod_{s=1}^n\left(ns\,t-n-1+\sum_j d_j\left(1-s/\rho_j\right)_+\right)\,A^n,
\]
while~\eqref{7.14_2} implies
\begin{equation}\label{7.25}
M'_{n,k}\leq\tfrac{n!\,(kn-1)!}{(n+kn-2)!}\, (kt+\tau)
\left(\sum_{1\,\leq\,s\,\leq\,k}\tfrac{1}{ s}\left(ns\,t-n-1+\sum_jd_j\left(1-s/\rho_j\right)_+
\right)\right)^{n-1}\,A^n.
\end{equation}
If we take $\rho_j\geq\rho>n$, then $(1-s/\rho_j)_+\geq 1-s/\rho$ for $s\leq n$, and as $ns\,t-n-1\geq 0$ and $\sum_{1\,\leq\,s\,\leq\,k}\tfrac{1}{ s}(nst-n-1)\leq nkt$, we get for $\tau>0$ small a sufficient condition
\[
\prod_{s=1}^n\left(\left(1-\tfrac{s}{ \rho}\right)\sum_jd_j\right)> kt\,(n+kn-1)\,n!\,
\left(nk\,t+\left(1+\tfrac{1}{ 2}+\cdots+\tfrac{1}{ k}\right)\sum_jd_j\right)^{n-1}.
\]
For $k=n$, the latter condition is satisfied if
$\sum_j d_j>c_nt\,\prod_{s=1}^n(1-\tfrac{s}{\rho})^{-1}$ with
\[
c_n=n\left(n^2+n-1\right)\,n!\,
\left(1+\tfrac{1}{ 2}+\cdots+\tfrac{1}{ n}+\tfrac{1}{ n^3}\right)^{n-1}.
\]
In fact, $c_1=1$, $c_2=32.5$ and $c_n\geq n^5$ for all $n\in\bN^*$, hence the above requirement implies in any case the inequality $n^2t\leq \tfrac{1}{ n^3}\sum d_j$. The Stirling and Euler--Maclaurin formulas give
\[
c_n\sim (2\pi)^{1/2} n^{n+7/2}\,e^{-n}\,(\gamma+\log n)^{n-1}
\]
as $n\to+\infty$, where $\gamma=0.577215\ldots$ is the Euler constant, the ratio being actually bounded above for $n\geq 3$ by
$\exp((1/2)(1-1/n)/(\gamma+\log n)+13/12n-1/n^2)\to 1$. Let us observe that
\[
\tfrac{1}{ t}=\min\left(\!\min_j\left(\tfrac{\rho_j}{ d_j}\right),\tfrac{1}{ 2}\right).
\]
In this way, we get the sufficient condition
\begin{equation}\label{7.27}
\rho_j\geq \rho>n,\quad
\sum_j d_j\cdot
\min\left(\!\min_j\left(\tfrac{\rho_j}{ d_j}\right),\tfrac{1}{ 2}\right)
\,\prod_{s=1}^n\left(1-\tfrac{s}{\rho}\right)>c_n.
\end{equation}
For instance, if we take all components $\Delta_j$ possessing the same degrees $d_j=d$ and ramification number $\rho_j\geq \rho$, these numbers and the number $N$ of components have to satisfy the sufficient condition
\begin{equation_}{N}\label{7.27_N}
\rho>n,\quad N\min(\rho,d/2)\,\prod_{s=1}^n\left(1-\tfrac{s}{\rho}\right)>c_n.
\end{equation_}
This possibly allows a single component (taking $d,\rho$ large), or $d,\rho$ small (taking $N$ large). Since we have neglected many terms in the above calculations, the ``technological constant'' $c_n$ appearing in these estimates is probably much larger than needed. Notice that the above estimates require jets of order $k\geq n$ and ramification numbers $\rho>n$. Parts ($\rma$) and ($\rma'$) of Theorem~\ref{theo:existence} follow from~\eqref{7.27} and~\eqref{7.27_N}.

\break

\stepcounter{equation}
\textbf{\theequation. Case of jet differentials of order k = 1 (symmetric differentials).}
When $k<n$ or $\rho_j\in{}]1,+\infty]$, estimate~\eqref{7.12} still allows us to obtain an existence criterion. For instance, when $k=1$ and $N\geq n$, \eqref{7.12_1} and~\eqref{7.25} give

\begin{align*}
M_{n,1}&\geq\tfrac{n!\,(n-1)!}{(2n-1)!}
\sum_{J\,\subset\,\{1,\,\dotsc,\,N\},\,\abs{J}=n}
\tfrac{n!\,(n-1)!}{ n^n}\,
\prod_{j\,\in\,J}d_j\left(1-\tfrac{1}{ \rho_j}\right)\,A^n,\\
M'_{n,1}&\leq \tfrac{n!\,(n-1)!}{(2n-2)!}\, (t+\tau)
\left(n\,t-n-1+\sum_jd_j(1-1/\rho_j)\right)^{n-1}\,A^n,
\end{align*}
and we get the non-void existence criterion
\begin{equation}\label{7.29}
\sum_{J\,\subset\,\{1,\,\dotsc,\,N\},\,\abs{J}=n}~
\prod_{j\,\in\,J}d_j\left(1-\tfrac{1}{ \rho_j}\right)> (2n-1)\,t\,
\left(n\,t-n-1+\sum_jd_j\left(1-1/\rho_j\right)\right)^{n-1}
\end{equation}
where $t=\max(\max_j(d_j/\rho_j),2)$. For instance, if all divisors have the same degrees $d_j=d$ and ramification numbers $\rho_j\geq\rho$, condition~\eqref{7.29} is implied by
\[
\binom{N}{ n}\,d^n\left(1-\tfrac{1}{ \rho}\right)^n> (2n-1) \max(d/\rho,2)\,\left((N+n)d\right)^{n-1},
\]
or equivalently, by
\[
\min(\rho,d/2)\,\binom{N}{ n}
\left(1-\tfrac{1}{ \rho}\right)^n>(2n-1)\,(N+n)^{n-1}.
\]
As $j\mapsto(N-j)/(n-j)$ is non-decreasing for $0\leq j<n\leq N$, we have the inequality $\binom{N }{ n}=\prod_{0\leq j<n}\tfrac{N-j}{ n-j}\geq (N/n)^n$, hence
\[
\tfrac{\displaystyle\binom{N}{ n}}{ (2n-1)(N+n)^{n-1}}
\geq \tfrac{N^n}{ n^n(2n-1)(2N)^{n-1}}=\tfrac{N}{ 2^{n-1}\,(2n-1)\,n^n}.
\]
We finally get the sufficient condition
\begin{equation_}{N}\label{7.29_N}
N\geq n,\quad N\min(\rho,d/2)\,\,\left(1-\tfrac{1}{ \rho}\right)^n>2^{n-1}\,(2n-1)\,n^n.
\end{equation_}
Parts ($\rmb$) and ($\rmb'$) of Theorem~\ref{theo:existence} follow from~\eqref{7.29} and~\eqref{7.29_N}. Again, the constant $2^{n-1}\,(2n-1)\,n^n$ is certainly far from being optimal. Answering the problem raised in Remark~\ref{6.43} might help to improve the bounds.

\appendix
\section{A proof of the orbifold vanishing theorem}

The orbifold vanishing theorem is proved in~\cite{CDR20} in the case of boundary divisors $D=\sum(1-1/\rho_j)\Delta_j$ with rational multiplicities
$\rho_j\in{}]1,\infty]$. However, the definition of orbifold curves shows that we can replace $\rho_j$ by $\lceil\rho_j\rceil\in\bN\cup\{\infty\}$ without modifying the space of curves we have to deal with. On the other hand, this replacement makes the corresponding sheaves $E_{k,m}V^*\langle D\rangle$ larger. Therefore, the case of arbitrary real multiplicities $\rho_j\in{}]1,\infty]$ stated in Proposition~\ref{prop:vanishing} follows from the case of integer multiplicities. We sketch here an alternative and possibly more direct proof of Proposition~\ref{prop:vanishing}, by checking that we can still apply the Ahlfors--Schwarz lemma argument of~\cite{Dem97} in the orbifold context. For this, we associate to $D$ the ``logarithmic divisor''
\[
D'=\lceil D\rceil = \sum\Delta_j \geq D,
\]
and, assuming $(X,V,D')$ non-singular, we make use of the tower of logarithmic Semple bundles
\[
X^\rS_k(V\langle D'\rangle)\to X^\rS_{k-1}(V\langle D'\rangle)\to\cdots\to X^\rS_1(V\langle D'\rangle)\to X^\rS_0(V\langle D'\rangle)
\bydef X
\]
(in reference to the work of the British mathematician John Greenlees Semple, see~\cite{Sem54}), where each stage is a smooth directed manifold
$(X^\rS_k(V\langle D'\rangle),V_k\langle D'\rangle)$ defined inductively by
\[
X^\rS_k(V\langle D'\rangle)
\bydef P(V_{k-1}\langle D'\rangle)=
\text{projective bundle of lines of $V_{k-1}\langle D'\rangle$},
\]
and $V_k\langle D'\rangle$ is a subbundle of the logarithmic tangent bundle of $X^\rS_k(V\langle D'\rangle)$ associated with the pull-back of $D'$. Each of these projective bundles is equipped with a tautological line bundle
$\cO_{X^\rS_k(V\langle D'\rangle)}(-1)$ (see~\cite{Dem97} for details), and $V_k\langle D'\rangle$ consists of the elements of the logarithmic tangent bundle that project onto the tautological line, so that we have an exact sequence
\[
0\to T_{X^\rS_k(V\langle D'\rangle)/X^\rS_{k-1}(V\langle D'\rangle)}
\to V_k\langle D'\rangle\to \cO_{X^\rS_k(V\langle D'\rangle)}(-1)\to 0.
\]
We let
$\pi_{k,\ell}\colon X^\rS_k(V\langle D'\rangle)\to X^\rS_\ell(V\langle D'\rangle)$ be the natural projection. Then the top-down projection $\pi_{k,0}\colon X^\rS_k(V\langle D'\rangle)\to X$ yields a direct image sheaf
\[
(\pi_{k,0})_*\cO_{X^\rS_k(V\langle D'\rangle)}(m)
\bydef E^\rS_{k,m}V^*\langle D'\rangle
\subset E_{k,m}V^*\langle D'\rangle.
\]
Its stalk at point $x\in X$ consists of the algebraic differential operators $P(f_{[k]})$ acting on germs of $k$-jets $f\colon(\bC,0)\to(X,x)$ tangent to $V$, satisfying the invariance property
\[
P((f\circ\varphi)_{[k]}) = (\varphi')^m P(f_{[k]})\circ \varphi,
\]
whenever $\varphi\in \bG_k$ is in the group of $k$-jets of biholomorphisms $\varphi\colon(\bC,0)\to(\bC,0)$. By construction, the sheaf of orbifold jet differentials $E_{k,m}V^*\langle D\rangle$ is contained in $E_{k,m}V^*\langle D'\rangle$, and we have a corresponding inclusion
\[
E^\rS_{k,m}V^*\langle D\rangle
\subseteq E^\rS_{k,m}V^*\langle D'\rangle
\]
of the Semple orbifold jet differentials into the Semple logarithmic differentials. A consideration of the algebra $\bigoplus E^\rS_{k,m}V^*\langle D\rangle$ makes clear that there exists a submultiplicative sequence of ideal sheaves $(\cJ_{D,k,m})_{m\,\in\,\bN}$ on $X^\rS_k(V\langle D'\rangle)$, such that the image of $\pi_{k,0}^*\cO_X(E^\rS_{k,m}V^*\langle D\rangle)$ in $\cO_{X^\rS_k(V\langle D'\rangle)}(m)$ is a sheaf
\[
\cO_{X^\rS_k(V\langle D'\rangle)}(m)\otimes\cJ_{D,k,m}.
\]
It is clear that the zero variety of $V(\cJ_{D,k,m})$ projects into the support $\abs{D'}=\abs{D}$ of $D$. We consider a smooth log resolution
\[
\mu_k
\colon
\tilde X_k
\to X^\rS_k(V\langle D'\rangle)
\]
of the ideal $\cJ_{D,k,m}$ in $X^\rS_k(V\langle D'\rangle)$, so that $\mu_k^*(\cJ_{D,k,m})=\cO_{\tilde X_k}(-G_{D,k,m})$ for a suitable effective simple normal crossing divisor $G_{D,k,m}$ on $\tilde X_k$ that projects into $\abs{D}$ in $X$. Denoting $\cO_{\tilde X_k}(1)=\mu_k^*\cO_{X^\rS_k(V\langle D'\rangle)}(1)$, we get

\[
\mu_k^*\left(\cO_{X^\rS_k(V\langle D'\rangle)}(m)\otimes\cJ_{D,k,m}\right)=
\cO_{\tilde X_k}(m)\otimes\cO_{\tilde X_k}\left(-G_{D,k,m}\right).
\]
We denote by $\tilde\pi_{k,\ell}$ the composition
\[
\tilde\pi_{k,\ell}=\pi_{k,\ell}\circ\mu_k:\tilde X_k\to X^\rS_k(V\langle D'\rangle)
\to X^\rS_\ell(V\langle D'\rangle),
\]
and consider especially the projection $\tilde\pi_{k,0}\colon\tilde X_k\to X$. For every entire or local orbifold entire curve
$f\colon\bC\supset\Omega\to(X,V,D)$, the image $f(\Omega)$ is not entirely contained in $\abs{D'}$, and we thus get holomorphic $k$-jet liftings
\[
f_{[k]}
\colon
\Omega
\to X^\rS_k(V\langle D'\rangle)
\text{\quad and\quad}
\tilde f_{[k]}
\colon
\Omega
\to
\tilde X_k.
\]
Moreover, the derivative $f'_{[k-1]}$ of the $(k-1)$-jet lifting
$f_{[k-1]}$ can be seen as a meromorphic section of the logarithmic tautological line bundle $(f_{[k]})^*\cO_{X^\rS_k(V\langle D'\rangle)}(-1)$, since the multipli\-cities of zeroes of $f'_{[k-1]}$ are possibly less than the ones prescribed by the logarithmic condition. The poles are of course contained in $f^{-1}(\abs{D'})$. As a consequence,
$f'_{[k-1]}$ also lifts as a meromorphic section of $(\tilde f_{[k]})^*\cO_{\tilde X_k}(-1)$, which we denote by $\tilde f'_{[k-1]}$. If $\tau_{D'}\in H^0(X,\cO_X(D'))$ is the canonical section of divisor equal to $D'$, we get at worst that
\begin{equation}\label{8.9}
\begin{aligned}
\tau_{D'}(f)\,f'_{[k-1]}&
\in H^0\big(
\Omega, (f_{[k]})^*\left(\cO_{X^\rS_k(V\langle D'\rangle)}(-1)\otimes \pi_{k,0}^*\cO_X(D')\right)
\big)\\
\text{and}\quad 
\tau_{D'}(f)\,\tilde f'_{[k-1]}&
\in H^0\big(\Omega, \left(\tilde f_{[k]}\right)^*\left(\cO_{\tilde X_k}(-1)\otimes\tilde\pi_{k,0}^*\cO_X(D')\right)
\big)
\end{aligned}
\end{equation}
are holomorphic. On the other hand, every local section $P\in H^0(U,E^\rS_{k,m}V^*\langle D\rangle)$ on an open subset $U\subset X$ gives rise in a one-to-one manner to a section
\[
\sigma_P\in H^0\left(U_k,\cO_{X^\rS_k(V\langle D'\rangle)}(m)
\otimes \cJ_{D,k,m}\right),\quad U_k=\pi_{k,0}^{-1}(U)\subset X^\rS_k(V\langle D'\rangle),
\]
by the correspondence
\[
P(f_{[k]}) =
\sigma_P(f_{[k]})\cdot \left(f'_{[k-1]}\right)^m
\]
for every local orbifold curve $f$ contained in $U$. By pulling back to $\tilde X_k$, we get a section
\[
\tilde\sigma_P\in H^0\left(\tilde U_k,\cO_{\tilde X_k}(m)
\otimes \cO_{\tilde X_k}\left(-G_{D,k,m}\right)\right),
\quad \tilde U_k=\mu_k^{-1}(U_k)=\tilde\pi_{k,0}^{-1}(U),
\]
such that
\[
P(f_{[k]})=\tilde\sigma_P\left(\tilde f_{[k]}\right)\cdot \left(\tilde f'_{[k-1]}\right)^m.
\]
However, $P(f_{[k]})$ is a holomorphic function, and we must have a cancellation of the poles of $(\tilde f'_{[k-1]})^m$ for all sections
$\tilde\sigma_P$, which generate the sheaf
$\cO_{\tilde X_k}(m)\otimes \cO_{\tilde X_k}(-G_{D,k,m})$. This means that
\begin{equation}\label{8.11}
\tilde f'_{[k-1]}
\text{ is a holomorphic section of } (\tilde f_{[k]})^*\cO_{\tilde X_k}(-1)\otimes
\cO_\bC\big(\lfloor\tfrac{1}{ m}(\tilde f_{[k]})^*G_{D,k,m}\rfloor\big)
\end{equation}
For any given ample divisor $A$ over $X$, we can find $s=s_{k,m}\in\bN^*$ such that the tensor product
$\cO_X(E^\rS_{k,m}V^*\langle D\rangle)\otimes\cO_X(sA)$ is generated by its global sections over $X$. By taking the pull-back to $\tilde X_k$ and looking at the image in $\cO_{\tilde X_k}(m)$, we conclude that
\begin{equation}\label{8.12}
\cO_{\tilde X_k}(m)\otimes\cO_{\tilde X_k}\left(-G_{D,k,m}\right)\otimes
\tilde\pi_{k,0}^*\cO_X(sA)\quad\text{is generated by sections on $\tilde X_k$}.
\end{equation}
As in~\cite{Dem97}, let us consider for every weight
$\underline a=(a_1,\,\dotsc,\,a_k)\in\bZ^k$ the line bundles
\begin{equation}\label{8.13}
\cO_{X^\rS_k(V\langle D'\rangle)}(\underline a)=
\bigotimes_{1\,\leq\,\ell\,\leq\,k}\pi_{k,\ell}^*\cO_{X^\rS_\ell(V\langle D'\rangle)}(a_\ell),
\qquad
\cO_{\tilde X_k}(\underline a)=\mu_k^*
\cO_{X^\rS_k(V\langle D'\rangle)}(\underline a).
\end{equation}
Since each factor $\cO_{X^\rS_\ell(V\langle D'\rangle)}(1)$ is relatively ample with respect to $\pi_{\ell,\ell-1}$, it is easy to see by induction on $k$ that thee exists a weight $\underline a\in(\bN^*)^k$ and $b\in\bN^*$ such that the line bundle
$\cO_{X^\rS_k(V\langle D'\rangle}(\underline a)\otimes\pi_{k,0}^*\cO_X(bA)$ is~ample. After possibly replacing $(\underline a,b)$ by a multiple, we can find a $\mu_k$-exceptional divisor $H_{D,k}$ on $\tilde X_k$ such that
\begin{equation}\label{8.14}
\cO_{\tilde X_k}(\underline a)\otimes
\cO_{\tilde X_k}(-H_{D,k})\otimes\tilde\pi_{k,0}^*\cO_X(bA)
\end{equation}
is very ample on $\tilde X_k$. Finally, we select $c\in\bN^*$ such that
\begin{equation}\label{8.15}
\cO_X(cA-D')\text{ is very ample on $X$}.
\end{equation}
By taking the tensor product of (\eqref{8.12}--\eqref{8.15}), \eqref{8.15} being raised to a power $t\in\bN^*$, we find that
\begin{multline}\label{8.16}
L_{k,m}
\bydef\\
\cO_{\tilde X_k}(m)\otimes\cO_{\tilde X_k}(\underline a)\otimes
\cO_{\tilde X_k}\left(-G_{D,k,m}-H_{D,k}\right)\otimes
\tilde\pi_{k,0}^*\cO_X((s+b+tc)A-tD')
\end{multline}
is very ample on $\tilde X_k$. We will later need to take
$t=\abs{\underline a}=\sum_\ell a_\ell$, which is of course an admissible choice.

\begin{lemm}
Let $(X,V,D)$ be a projective non-singular directed orbifold, and $A$~an ample divisor on $X$. Then, for every orbifold entire curve $f\colon\bC\to(X,V,D)$ and every section
\[
P\in H^0\big(X,E^\rS_{k,m}V^*\langle D\rangle\otimes\cO_X(-A)\big),
\]
we have $P(f_{[k]})=P(f,f',\dotsc,f^{(k)})=0$.
\end{lemm}
\begin{proof}
As we have already seen for local sections, every global jet differential
\[
P\in H^0\big(X,E^\rS_{k,m}V^*\langle D\rangle\otimes\cO_X(-A)\big)
\]
gives rise to sections
\begin{align*}
\sigma_P&
\in H^0\big(X^\rS_k(V\langle D'\rangle),
\cO_{X^\rS_k(V\langle D'\rangle)}(m) \otimes \cJ_{D,k,m}\otimes\pi_{k,0}^*\cO_X(-A)
\big),
\\
\tilde\sigma_P&
\in H^0\big(
\tilde X_k,
\cO_{\tilde X_k}(m) \otimes \cO_{\tilde X_k}(-G_{D,k,m})\otimes\tilde\pi_{k,0}^*\cO_X(-A)
\big)
\end{align*}
such that
\[
P\left(f_{[k]}\right)=\tilde\sigma_P\left(\tilde f_{[k]}\right)\cdot \left(\tilde f'_{[k-1]}\right)^m
\in H^0\left(\bC,f^*\cO_X(-A)\right).
\]
Assume that $P(f_{[k]})\ne 0$ (so that, in particular $\tilde\sigma_P\ne 0$). We consider a basis $(g_j)$ of sections of $L_{k,m}$ in~\eqref{8.16}, the canonical section $\eta_{D,k}\in H^0(\tilde X_k,\cO_{\tilde X_k}(H_{D,k}))$ and take the products

\begin{multline}\label{8.18}
h_j=\\
g_j\,(\tilde\sigma_P)^{q-1}\,(\tau_{D'})^t~\eta_{D,k}
\in H^0\left(\tilde X_k,
\cO_{\tilde X_k}(mq)\otimes\cO_{\tilde X_k}(\underline a)\otimes
\cO_{\tilde X_k}\left(-qG_{D,k,m}\right)\right)
\end{multline}
where $q=s+b+tc+1$. We now observe, thanks to our choice $t=\abs{\underline a}=\sum a_\ell$, that
\begin{multline}\label{8.19}
h_j\left(\tilde f_{[k]}\right)\cdot \left(\tilde f'_{k-1}\right)^{mq}\cdot \prod_{1\,\leq\,\ell\,\leq\,k} \left(d\tilde\pi_{k,\ell}\left(\tilde f'_{k-1}\right)\right)^{a_\ell}
=\Big(\tilde\sigma_P\left(\tilde f_{[k]}\right)\cdot \left(\tilde f'_{k-1}\right)^m\Big)^{q-1}
\\
\times\Big(g_j\left(\tilde f_{[k]}\right)\cdot \left(\tilde f'_{k-1}\right)^m\cdot
\prod_{1\,\leq\,\ell\,\leq\,k}d\tilde\pi_{k,\ell}\left(\tau_{D'}(f)\,\tilde f'_{k-1}\right)^{a_\ell}\Big)
\times \eta_{D,k}\left(\tilde f_{[k]}\right)
\end{multline}
is a product of holomorphic sections on $\bC$, by~\eqref{8.11} and~\eqref{8.9} combined with~\eqref{8.16} and~\eqref{8.18}, and the fact that $P(f_{[k]})=\tilde\sigma_P(\tilde f_{[k]})\cdot (\tilde f'_{k-1})^m$ is holomorphic with values in $f^*\cO_X(-A)$. The product also takes value in the trivial bundle over $\bC$, and can thus be seen as a holomorphic function. As $j$ varies, these functions are not all equal to zero, and we define a hermitian metric $\gamma(t)=\gamma_0(t)\,\abs{dt}^2$ on the complex line $\bC$ by putting
\[
\gamma_0 =
\left(
\sum_je^{\psi\left(\tilde f_{[k]}\right)}
\abs*{h_j\left(\tilde f_{[k]}\right)\cdot\left(\tilde f'_{k-1}\right)^{mq}
\cdot
\prod_{1\,\leq\,\ell\,\leq\,k} d\tilde\pi_{k,\ell}\left(\tilde f'_{k-1}\right)^{a_\ell} }^2
\right)^{\tfrac{1}{mq+\abs{\underline a}}},
\]
where $\psi$ is a quasi-plurisubharmonic potential on $\tilde X_k$ which will be chosen later. Notice that $\gamma_0(t)$ is locally bounded from above and almost everywhere non-zero. Since~\eqref{8.19} only involves holomorphic factors in the right hand side, we get
\begin{equation}\label{8.21}
\imath\ddbar\log\gamma_0\ge
\tfrac{1}{ mq+\abs{\underline a}}\,\left(\tilde f_{[k]}\right)^*\left(\tilde\omega_k+\imath\ddbar\psi\right)
\end{equation}
where $\tilde\omega_k=\imath\ddbar\log\abs{g_j}^2$ is a K\"ahler metric on $\tilde X_k$, equal to the curvature of the very ample line bundle $L_{k,m}$ for the projective embedding provided by $(g_j)$. (In fact, \eqref{8.21} could be turned into an equality by adding a suitable sum of Dirac masses in the right hand side). Of course, $\psi$ will be taken to be an $\omega$-plurisubharmonic potential on $\tilde X_k$. We wish to get a contradiction by means of the Ahlfors--Schwarz lemma (see e.g.~\cite[Lemma~3.2]{Dem97}), by showing that $\imath\ddbar\log\gamma_0\geq A\gamma$ for some $A>0$, an impossibility for a hermitian metric on the entire complex line. Since $\psi$ is locally bounded from above, by~\eqref{8.19} and the inequality between geometric and arithmetic means, we have
\begin{equation}\label{8.22}
\gamma_0(t)
\leq C\,
\left(\sum\left|h_j\left(\tilde f_{[k]}(t)\right)\right|^2\right)^{\tfrac{1}{mq+\abs{\underline a}}}
\left|f'_{[k-1]}(t)\right|_{\log}^2
\end{equation}
where $C>0$ and the norms $\abs{h_j}^2$ and $\abs{f'_{[k-1]}(t)}_{\log}^2$ are computed with respect to smooth metrics on
$\cO_{\tilde X_k}(mq)\otimes\cO_{\tilde X_k}(\underline a)\otimes
\cO_{\tilde X_k}(-qG_{D,k,m})$ and on the logarithmic tautological line bundle $\cO_{X^\rS_k(V\langle D'\rangle)}(-1)$, respectively. The term $\abs{h_j}^2$ is bounded, but one has to pay attention to the fact that $\abs{f'_{[k-1]}(t)}_{\log}^2$ has poles on $f^{-1}(\abs{D'})$. If we use local coordinates $(z_1,\,\dotsc,\,z_n)$ on $X$ such that $\Delta_j=\{z_j=0\}$, we have

\[
\left|f'_{[k-1]}\right|_{\log}^2\sim
\left|f'_{[k-1]}\right|_{\omega_k}^2+\sum_j\abs{f_j}^{-2}\,\left|f'_j\right|^2
\]
in terms of a smooth K\"ahler metric $\omega_{k-1}$ on $X^\rS_k(V\langle D'\rangle)$. What saves us is that $h_j$ contains a factor $\tau_{D'}(f)^t$ that vanishes along all components $\Delta_j$. Therefore~\eqref{8.22} implies the existence of a number $\delta>0$ such that
\[
\gamma_0(t)\leq C'\,\Big(\,\left|f'_{[k-1]}(t)\right|_{\omega_{k-1}}^2 +\sum_j\abs{f_j}^{-2+2\delta}\,\left|f'_j\right|^2\Big).
\]
Since the morphism $\tilde\pi_{k,k-1}$ has a bounded differential and $f'_{[k-1]}(t)=d\tilde\pi_{k,k-1}(\tilde f'_{[k]}(t))$, we infer
\begin{equation'}\label{8.22''}
\gamma_0(t)\leq C''\,\left(\,\left|\tilde f'_{[k]}(t)\right|_{ \tilde\omega_k}^2 +\sum_j\abs{f_j}^{-2+2\delta}\,\left|f'_j\right|^2\right).
\end{equation'}
By~\eqref{8.21} and~\eqref{8.22''}, in order to get a lower bound $\imath\ddbar\log\gamma_0\geq A\gamma$, we only need to choose the potential $\psi$ so that
\begin{equation}\label{8.23}
\sum_j\abs{f_j}^{-2+2\delta}\,\abs{f'_j}^2\leq C'''
\,\left(\tilde f_{[k]}\right)^*\left(\tilde\omega_k+\imath\ddbar\psi\right).
\end{equation}
If $\tau_j\in H^0(X,\cO_X(\Delta_j))$ is the canonical section of divisor $\Delta_j$, \eqref{8.23} is achieved by taking
$\psi=\varepsilon\sum_j \abs{\tau_j\circ\tilde\pi_{k,0}}^{2\delta}$, for any choice of a smooth hermitian metric on $\cO_X(\Delta_j)$ and $\varepsilon>0$ small enough. In some sense, we have to take a suitable orbifold K\"ahler metric $\tilde\omega_k+\imath\ddbar\psi$ on $\tilde X_k$ to be able to apply the Ahlfors--Schwarz lemma. It might be interesting to find the optimal choice of $\delta>0$, but this is not needed in our proof.
\end{proof}


\begin{proof}[End of the proof of Proposition~\ref{prop:vanishing}]
We still have to extend the vanishing result to the case of non-necessarily $\bG_k$-invariant orbifold jet differentials
\[
P\in H^0\big(X,E_{k,m}V^*\langle D\rangle\otimes\cO_X(-A)\big).
\]
One can then argue by using the $\bG_k$-action on jet differentials
\[
(\varphi,P)\mapsto\varphi^*P,\qquad \left(\varphi^*P\right)(f_{[k]})
\bydef P((f\circ\varphi)_{[k]})\circ\varphi^{-1},\quad
\varphi\in\bG_k.
\]
This action yields a decomposition
\[
\left(\varphi^*P\right)\left(f_{[k]}\right) =
\sum_{\substack{
\alpha\,\in\,\bN^k\\
\abs{\alpha}_w=m \phantom{\textstyle|} }} \left(\varphi^{(\alpha)}
\circ\varphi^{-1}\right)\,P_\alpha(f_{[k]}),\quad P_\alpha\in H^0\left(X,E_{k,m_\alpha}V^*\langle D\rangle\otimes\cO_X(-A)\right)
\]
where $\alpha=(\alpha_1,\,\dotsc,\,\alpha_k)\in\bN^k$,
$\varphi^{(\alpha)}=(\varphi')^{\alpha_1}(\varphi'')^{\alpha_2}\ldots (\varphi^{(k)})^{\alpha_k}$,
$\abs{\alpha}_w=\alpha_1+2\alpha_2+\dotsb+k\alpha_k$ is the weighted degree, and $P_\alpha$ is a homogeneous polynomial of degree
\[
m_\alpha
\bydef
\deg P_\alpha=m-(\alpha_2+2\alpha_3+\dotsb+(k-1)\alpha_k)=
\alpha_1+\alpha_2+\dotsb+\alpha_k.
\]
In particular $\deg P_\alpha<m$ unless $\alpha=(m,0,\dotsc,0)$, in which case
$P_\alpha=P$. If the result is known for degrees${}<m$, then all
$P_\alpha(f_{[k]})$ vanish for $P_\alpha\ne P$ and one can reduce the proof to the invariant case by induction, as the term $P_\alpha$ of minimal degree is invariant. The proof makes use of induced directed structures, and is purely formal and group theoretic. Essentially, the argument is that $P$ becomes an invariant jet differential when restricted to the subvariety of the Semple $k$-jet bundle consisting of germs $g_{[k]}$ of $k$-jets such that $P_\alpha(g_{[k]})=0$ for
$P_\alpha\ne P$. Singularities may appear in this subvariety, but this does not affect the proof since the induced directed structure is embedded in the non-singular logarithmic Semple tower. We refer the reader to~\cite[\S~7.E]{Dem20} and~\cite[Theorem~8.15]{Dem20} for details.
\end{proof}


\bibliography{Campana}
\end{document}
